\documentclass[../main.tex]{subfiles}
\begin{document}
%==========================================TIÊU ĐỀ TẬP========================================
\begin{table}[h]
  \begin{adjustwidth}{-1cm}{}
    \begin{tabular}{>{\centering\arraybackslash}p{6cm}|>{\raggedright\arraybackslash}p{12.5cm}}
      \multirow{5}{6cm}
      % Cột bên trái: ÉPISODE và số tập
      {\\[-40pt]
        \flaregothic\fontsize{20pt}{20pt}\selectfont \centering
        {\contour{errorone}{\textcolor{black}{ÉPISODE}}}\color{black}\\
        \burbank\fontsize{72pt}{72pt}\selectfont
        \includegraphics[height=3cm]{../../../../Image/Book?/number/num0.png}
        \\[18pt]
        % \flaregothic\fontsize{14pt}{14pt}\selectfont
        % \color{epattr}BIENVENUE, \uline{\color{Banpire} Mel} !
        % \flaregothic\fontsize{18pt}{18pt}\selectfont
        % \color{epattr}ÉPISODE PERDU
        \color{black}
      }
       &
      % Dòng thứ nhất: tên tập tiếng Nhật
      {\fotskip\fontsize{18pt}{18pt}\selectfont hololiveERROR
      }                                                                                                  \\[6pt]
      % Dòng thứ hai: tên tập tiếng Anh
       & {\cafeteria\fontsize{18pt}{18pt}\selectfont The Samsara's Call}                                                                                      \\[4pt]
      % Dòng thứ ba: tên tập tiếng Việt
       & {\vnmsans\fontsize{18pt}{18pt}\selectfont Tiếng gọi từ vòng lặp}                                                                                     \\[8pt]
      % Dòng thứ tư: Nhân vật xuất hiện
       & {\caviar{\fontsize{14pt}{14pt}\selectfont Track used: \href{https://www.youtube.com/watch?v=qq0BDRhBCF8}{The Nighmare - theo5970} (from Wind Line)}}
      \\[6pt]
      % Dòng cuối cùng: Thời gian diễn ra của tập (thường sẽ là ngày phát hành)
       & {\continuum\fontsize{16pt}{16pt}\selectfont Ngày phát hành: 07/01/2022}                                                                              \\[8pt]
    \end{tabular}
  \end{adjustwidth}
\end{table}
%===========================================NỘI DUNG==========================================
\color{black}

\begin{center}
  \uline{(Câu chuyện này lấy từ phần game demo của hololive ERROR. Phần bản dịch được tham khảo từ bản dịch tiếng Việt của \href{https://www.youtube.com/@namakenekoji2023}{Namake Nekoji} và Google Dịch.\\Các bức ảnh được lấy trong phần game demo, được chơi bởi \href{https://www.youtube.com/@weiweilim}{laowei老威}.)}
\end{center}

\pov{???}

\dots

\dots

Một giọng nói nữ vọng lại từ đâu đó trong không gian mờ ảo.

``Các cậu thấy trường học có vui không?''

Giọng nói ấy nghe như đang vang vọng từ rất xa, nhưng đồng thời cũng như thể đang thì thầm ngay bên tai.

\dots

``Tớ dám chắc là các cậu thấy vui.''

Âm sắc trong giọng nói ấy thật dịu dàng, ấm áp\dots\ và kỳ lạ thay, nó khiến tôi cảm thấy nhói đau.

*BÍNH BONG\dots\ BÍNH BONG\dots*

Tiếng chuông báo hết giờ vang vọng, như một tiếng vọng xa xăm từ quá khứ.

\dots

``\dots Đã hết tiết rồi đó. Thôi nào, ta hãy cùng về thôi.''

\dots

Giọng nói ấy lại cất lên. Một giọng nói quen thuộc đến kỳ lạ.

Tôi biết giọng này. Tôi chắc chắn là mình biết.

Nhưng\dots

\dots tại sao tôi lại không thể nhớ ra chủ nhân của giọng nói ấy?

Tôi có thể cảm nhận được sự thân thuộc trong từng âm điệu, từng câu chữ.

Nhưng cái tên ấy\dots\ nó cứ tuột khỏi tâm trí tôi như cát chảy qua kẽ tay.

\dots Tại sao?

Tại sao tôi không thể nhớ?

Tại sao!?

ĐIỀU NÀY KHÔNG CÓ LÝ GÌ CẢ!

\dots

Đầu tôi nhức nhối.

Giọng nói kia vẫn văng vẳng đâu đây\dots

Như một bóng ma ám ảnh, nó vừa gần gũi đến độ tôi có thể chạm tới\dots

\dots nhưng lại xa xôi như một giấc mơ tan biến khi tỉnh giấc.

\dots

Một ý nghĩ chợt lóe lên.

\dots Khoan đã.

\dots Tôi là ai?

Tâm trí tôi trống rỗng.

Hoảng loạn dâng lên như thủy triều.

\dots Không thể nào.

\dots Làm ơn, hãy nhớ ra đi!

Tôi cố gắng lục lọi trong ký ức\dots

Nhưng chỉ thấy một màn sương mù dày đặc.

\dots Chuyện quái gì đang xảy ra vậy?

Tôi không thể nhớ nổi\dots\ bản thân mình là ai\dots

Làm sao có thể? Làm sao một người lại có thể quên đi chính danh tính của mình--

\dots

\dots Khoan.

Còn một điều kỳ lạ hơn nữa.

Tôi đang ở đây một mình, phải không?

Vậy mà ngoài giọng nói bí ẩn kia\dots

\dots tôi không nghe thấy bất kỳ âm thanh nào khác.

Không một tiếng động.

Không một hơi thở.

Chỉ có sự tĩnh lặng đến rợn người.

\dots

\dots

\ct

``\dots -chan! Tỉnh lại đi, Nii-chan!''

Lại một giọng nói khác vang lên. Cũng là giọng con gái, và\dots\ quen thuộc nữa.

\dots

\dots Nhưng khác hẳn với giọng nói đầu tiên. Giọng này có gì đó gấp gáp hơn, run rẩy, như thể người nói đang sợ hãi điều gì đó.

``\textbf{Nii-chan! Mau dậy đi! Đừng bỏ em lại ở đây một mình mà!}''

\dots

\dots Một tiếng khóc?

Và ``Nii-chan''\dots? Có nghĩa là tôi là anh trai của người đó à?

\dots

\dots Vậy\dots\ tôi không chỉ nghe thấy, mà còn cảm nhận được rằng còn có một người khác đang ở đây. Không, tôi chắc chắn là đang không ở một mình.

Ai đó thực sự đang cố lay tôi dậy.

\dots

Tay tôi cảm nhận được một lực kéo nhẹ, lắc lư nơi vai. Ấm áp. Gần gũi.

\dots

Tôi có thể mở mắt ra.

\dots

Tôi từ từ mở mắt. Ngay lập tức, một tia sáng trắng nhòe chiếu thẳng vào khiến tôi phải nheo mắt lại theo phản xạ. ``S-Sáng quá\dots\ Ặc\dots''

Tôi đưa tay lên che bớt ánh sáng. Trong quầng sáng mờ mờ ấy, tôi thấy dáng người của một cô gái đang cúi xuống, hai tay vẫn đang bấu lấy vai tôi.

Tôi rên rỉ khẽ một tiếng, mắt cố điều chỉnh để thích nghi dần với ánh sáng. Mọi thứ trước mắt vẫn còn mờ mịt, nhưng rồi\dots

\dots tôi thấy gương mặt của cô ấy. Cô ấy đang khóc.

``May quá trời! Anh tỉnh rồi!''

Cô ấy bật thốt lên, nước mắt vẫn còn lăn dài trên má, nhưng ánh mắt đã ánh lên niềm vui mừng đến mức tôi thấy như nghẹn lại.

Cô ấy mặc một bộ đồng phục học sinh: áo sơ-mi trắng ngắn tay, trước ngực là một chiếc nơ màu xanh dương, váy ca-rô cùng tông kéo dài tới gần đầu gối. Tóc cô ấy có vẻ hơi rối, buộc đuôi ngựa bằng một chiếc nơ đỏ.

Tôi cũng đang mặc đồng phục khá giống với cô ấy, chỉ khác là thay nơ bằng cà vạt, và dĩ nhiên là\dots\ quần dài thay vì váy. Cả hai món đều là màu xanh nước biển giống hệt.

Tôi cũng chẳng rõ từ khi nào chúng tôi lại mặc như thế này. Mà đúng hơn\dots\ từ bao giờ tôi đã ở đây?

Tôi gượng ngồi dậy, vẫn hơi choáng váng. Cô gái ấy liền sà tới, vòng tay ôm tôi thật chặt.

Tôi có thể cảm nhận được hơi ấm từ cơ thể cô ấy. Thật rõ ràng.

Cơ mà\dots

``\dots Khuôn mặt này\dots\ quen quen,'' tôi buột miệng.

Nghe vậy, cô gái đang ôm tôi phụng phịu ngay tức thì: ``Mồ! Từ nãy tới giờ anh vẫn chưa tỉnh sao? Em tưởng anh chết luôn rồi đó!''

\dots Có lẽ tôi thực sự vẫn chưa tỉnh táo hoàn toàn.

Đầu tôi nặng trĩu. Mọi thứ như vừa rơi xuống sau một cơn mơ dài.

Cô ấy vội lau nước mắt trên mặt mình, rồi nhìn tôi cười. Một nụ cười nhẹ nhõm, và nói: ``Là em đây, Haragen Saikyou (\ruby{原}{はら}\ruby{阮}{げん}\ruby{崔}{さい}\ruby{強}{きょう}, \ruby{Ha-ra}{Nguyên}-\ruby{ghen}{Nguyễn}\ \ruby{Xai}{Thôi}-\ruby{kiêu}{Cường})! Em gái song sinh của anh!''

\dots

\dots Hả?

Tôi có em gái à? Mà lại còn là\dots\ song sinh?

Chuyện gì đang diễn ra vậy?

``Trời ạ, Nii-chan,'' cô ấy lắc đầu, ``nhìn mặt anh là biết liền anh chẳng nhớ gì rồi!''

``\dots Thật sự anh không nhớ nổi chuyện gì đã xảy ra cả,'' tôi đáp, thành thật. Cứ cho cô gái tên Saikyou này là em gái tôi đi, nhưng\dots\ mọi thứ vẫn quá mơ hồ.

``Anh là Haragen Kyuen (原阮久遠), anh có biết không?''  cô em gái (?) nói, mặt tỏ vẻ đắc ý, như thể đoán được tôi đang nghĩ gì. ``Tên của anh viết là Cửu Viễn đó!''

\dots

Được rồi. Tóm lại\dots

Tôi tên là Haragen Kyuen.

Cô gái này tên là Haragen Saikyou, em gái sinh đôi của tôi.

Tôi là anh lớn. Cô ấy là em gái.

Chúng tôi\dots là anh em sinh đôi, cùng họ Haragen.

Nếu tôi không nhầm, chữ ``gen'' (\ruby{阮}{げん}) này khá hiếm gặp trong tiếng Nhật\dots\ và hình như là ``Nguyễn'' theo âm Hán Việt thì phải?

Và\dots

\dots

\dots Tôi không nhớ gì khác. Không một thứ gì rõ ràng hơn.

``\dots Ta đang ở đâu đây?'' tôi hỏi, cố gắng ngồi thẳng lên một lần nữa.

Cô ấy nhìn quanh, rồi đáp, có phần lúng túng: ``Em cũng vừa mới tỉnh lại thôi. Mải gọi anh dậy quá nên cũng chưa kịp để ý\dots''

Bây giờ tôi mới để ý tới mọi thứ xung quanh. Bên cạnh, ngoài có cô em gái tôi đang nhìn với vẻ lo âu\dots

\img{e0a}

\dots chúng tôi còn nhìn thấy bóng dáng của một cô gái đang đi về phía trước, ở cách chúng tôi khá xa. Cách cô ấy đi đung đưa như một con lật đật, dường như tỏ vẻ khá vô tư.

\dots

\dots

\dots Tại sao tôi thấy cô ấy trông có vẻ quen quen nhỉ?
Hình như tôi đã từng thấy cô ấy ở đâu rồi\dots

Tôi lớn giọng, định gọi cô nàng đang khuất đi: ``\textbf{N-Này, cô gì đó ơi!}''

\dots

\dots

Không có tiếng trả lời. Cô gái ấy vẫn thong dong đi về phía trước, cho đến khi khuất hẳn khỏi tầm mắt chúng tôi.

\dots

\dots Cơ mà, thực sự đấy, chúng tôi đang ở đâu đây?

Nhìn xung quanh, tôi thấy đây là một hành lang của một trường học cấp ba. Một bên có những ô cửa sổ khá lớn chiếu ánh sáng mặt trời vào, bên dưới có lấp ló những tủ đồ xếp ngăn nắp.

Phần sàn được lát bê tông khá chắc, trên đó in những hình bóng của các ô cửa sổ do ánh sáng chiếu vào. Không có gì quá đặc biệt. Trần nhà được lát bằng các gạch ô vuông, trên đó có treo những chiếc đèn huỳnh quang và một hộp loa thông báo.

Đối diện với các ô cửa sổ là những lớp học, mỗi lớp có hai cửa trượt ra vào, phía trước mỗi lớp là một tấm bảng lục để đính thông báo của trường. Hoặc ít nhất là tôi nghĩ vậy. Tôi cá là mỗi hành lang sẽ có đâu đó khoảng bốn, năm lớp học như thế, lớp nào cũng giống lớp nào, và\dots\ không có đánh số lớp. Kỳ lạ.

Nhìn vào bên trong một trong những lớp ấy, tôi thấy những vật dụng mà một lớp học bình thường hay có: một cái bảng phấn xanh, bốn hàng bàn ghế ngồi một người, mỗi hàng có từ năm đến sáu chiếc ghế, vài ba chiếc quạt trần, máy chiếu, hệ thống đèn huỳnh quang trên trần và một chiếc đồng hồ cơ treo trên tường.

Nhưng điều kỳ quái là\dots\ chiếc đồng hồ đó thay vì quay theo chiều kim đồng hồ như lẽ thường, nó lại đang quay ngược lại. Tại sao lại như thế?

\dots

\dots

\dots Tôi cũng chẳng biết nữa. Chỉ riêng việc chúng tôi đang ở đây thôi đã là điều không thể lý giải được rồi, huống chi là giải thích những hiện tượng kỳ lạ như vậy.

``Có vẻ như là\dots'' tôi thở mạnh, nói trong sự ngập ngừng. Saikyou (từ giờ tôi sẽ gọi em ấy là như vậy, dù tôi không quá chắc chắn) nối tiếp lời tôi: ``\dots một hành lang trường học nhỉ?''

``Nhưng nếu thế thì\dots\ mọi người đâu cả rồi?'' tôi nói. Lớp học mà tôi ngó vào ban nãy hoàn toàn không có lấy một bóng người. Không. Có. Một. Người. Nào. Cả trên hành lang cũng vậy. Vắng như chùa Bà Đanh.

Thậm chí, ngoài tiếng chúng tôi nói chuyện ra, không có lấy một âm thanh nào khác. Cũng không còn giọng nói nào nữa, trừ cái giọng của một cô gái mà tôi thoáng nghe ban nãy, giờ cũng đã im lìm hẳn.

``Hmm\dots\ Lạ thật nhỉ, Nii-chan.''

Quả thật như em ấy nói, mọi chuyện đang diễn ra ngay trước mắt rất kỳ lạ, đến mức đáng sợ.

Chúng tôi bị lạc vào một ngôi trường bỏ hoang nào đó mà chính chúng tôi còn không biết ư?

Đây thực sự là đâu? Tại sao chúng tôi lại ở đây?

Và cô gái mà chúng tôi vừa nhìn thoáng qua kia là ai!?

Đó thực sự là những câu hỏi hóc búa mà chúng tôi cần những lời giải đáp thỏa đáng. Từ một cái gì đó. Một điều gì đó.

Hoặc một ai đó.

\dots

\dots

\dots

Chúng tôi im lặng một hồi lâu, chờ đợi xem có một thứ gì đó xảy ra không.

\dots

\dots

Tuyệt nhiên không có lấy một cái gì xảy ra cả. Em gái tôi cất lời: ``Nii-chan, em nghĩ chúng ta nên đi thôi nhỉ?''
``Nhưng vấn đề là đi đâu?'' tôi hỏi, thở dài trong sự tuyệt vọng.

``Trước tiên thì\dots\ ta nên tìm cách rời khỏi đây,'' em tôi đáp lại với chút lưỡng lự.

Đó có lẽ là giải pháp tốt nhất.

Chúng tôi quay đầu lại. Sau lưng là một cánh cửa trượt, thứ duy nhất khác biệt giữa bốn bức tường trống trơn.

Ánh sáng trắng lóa chói vào mắt đến nỗi tôi không nhìn được gì ở phía bên kia cả, nhưng tôi có cảm tưởng cứ như thể bên kia là một thế giới hoàn toàn khác vậy.

Saikyou bước lại gần, thử kéo cánh cửa.

*LẠCH CẠCH\dots\ LẠCH CẠCH\dots*

Nó không nhúc nhích dù chỉ một chút.

``Nii-chan, em không mở được cái này!'' em gái tôi nói, có vẻ hơi bực bội.

Tôi bước lên. ``Để anh.''

Hai chúng tôi cùng thử kéo cửa.

*LẠCH CẠCH\dots*

Không hề xi nhê.

Lạ thật, nó có vẻ không hề bị khóa, mà đúng hơn là đã bị kẹt chết từ bên ngoài, hoặc có thứ gì đó giữ nó lại.

Xung quanh cũng chẳng có gì có thể dùng để phá cửa, không gạch đá, không xà beng, thậm chí cả một cái ghế nhựa cũng không có.

``Chúng ta bị nhốt ở đây à\dots'' tôi lẩm bẩm.

Một hành lang dài hiện ra phía trước. Ánh sáng mờ từ các cửa sổ rọi vào làn khói lững lờ bay nhẹ, khiến nơi này mang vẻ ma mị khó tả.

``Chắc phải tìm một lối khác thôi,'' Saikyou nói, ngập ngừng, ``hy vọng có cầu thang hay gì đó\dots''

Tôi gật. ``Ừ, đi thôi.''

Chúng tôi bắt đầu bước đi. Tiếng bước chân lộp cộp, lộp cộp vang vọng khắp hành lang vắng lặng, mỗi âm vang nghe như vọng từ một nơi khác, một nơi không hẳn là ở ``thời hiện tại.''

Không khí ở đây khiến tôi sởn da gà. Thỉnh thoảng, tôi lại nghe như có tiếng máy móc chạy ngầm, hoặc là tiếng gió rít, không thể phân biệt rõ.

Saikyou đi sát bên tôi, hai tay nắm chặt lấy tay áo tôi như thể chỉ cần buông ra là em ấy sẽ lạc mất.

Bỗng em dừng lại, cúi đầu nhìn xuống chân. ``Nii-chan\dots\ cái này là gì vậy?''

Tôi cúi xuống theo. Một mảnh giấy.

Tôi nhặt nó lên. Nét chữ nguệch ngoạc, chữ to chữ nhỏ như được viết trong lúc hoảng loạn. Dù vậy, vẫn đủ đọc.

``{\vnmsans Vụ việc Trường Trung học Aogami bị sập: Ký sự số 1.}''

Tôi khựng lại. Trường Aogami?

Tôi nhìn lại bộ đồng phục mình đang mặc. Quả đúng như Saikyou-chan nói lúc nãy, có một phù hiệu được gắn ngay trên túi áo: là hình thêu chìm, hơi sờn, nhưng vẫn còn đọc được: ``Trường Trung học Aogami.''

Vậy\dots\ nơi này là ngôi trường được nhắc tới trong mảnh giấy đó?

Khoan đã\dots

Nếu đây là Aogami, thì vụ trường bị sập là chuyện gì?

Tôi nhìn quanh, tất cả mọi thứ trông đều nguyên vẹn, dù hơi cũ kỹ. Không có dấu hiệu gì của đổ sập cả.

Tôi nhìn vào mảnh giấy ra đọc tiếp.

``{\vnmsans Tháng Bảy, năm Bình Thành thứ--}'' Phần sau đã bị gạch đen kín. Không đọc được số.

Chỉ biết đây là thời kỳ Heisei, từ tháng 1/1989 đến tháng 4/2019. Nếu đây là ký sự cũ thì chắc chắn vụ việc đã xảy ra trong quá khứ. Khả năng cao là sẽ ở đâu đó khoảng thập niên 90.

Tôi tiếp tục đọc:
\txtbx{
  {\arial "Tại thị trấn Aogami, đã xảy ra một tai nạn thảm khốc khiến cho ngôi trường của thị trấn bị đổ sập. Cảnh sát bước đầu xác định nguyên nhân là do nổ khí ga.\\Tuy nhiên, vẫn còn tồn tại nhiều điều uẩn khúc mà chưa có lời lý giải rõ ràng."}}
Tôi nhíu mày. Nổ khí ga? Có vẻ như đây không đơn giản chỉ là một tai nạn.

Và tại sao tôi và Saikyou lại tỉnh dậy ở chính ngôi trường đã từng bị đổ sập ấy?

``Em thấy chuyện này\dots\ có gì đó không đúng,'' Saikyou-chan lên tiếng, giọng nhỏ nhưng chắc nịch.

Tôi gật đầu. ``Anh cũng vậy. Mà nếu đây là `ký sự số 1'\dots\ thì có lẽ còn những bản khác nữa.''

Tôi gấp mảnh giấy lại, bỏ vào túi áo. ``Biết đâu chúng sẽ giải đáp cho chúng ta phần nào.''

*BỘP*

Một tiếng động khẽ vang lên khi một vật nhỏ rơi xuống từ mảnh giấy kia. Theo phản xạ, tôi cúi xuống nhặt nó lên.

Một đoạn dây, dài chừng hai phần ba gang tay, chỉ rộng khoảng hai phân, với chín lỗ tròn đục dọc theo chiều dài của nó. Mỗi lỗ rộng chừng nửa phân. Chất liệu là cao su, một đầu còn được gắn một trục kim loại nhỏ, trông có vẻ tinh xảo một cách kỳ lạ.

``Nii-chan, anh nghĩ đây là cái gì đây?'' Saikyou đưa tay ra, vẻ mặt vừa tò mò vừa nghi ngờ.

Tôi cầm lấy, xoay xoay nó trong lòng bàn tay.

``\dots Trông giống như dây đeo của một cái đồng hồ.''

Một thứ như vậy\dots\ tại sao lại xuất hiện ở đây? Trong một nơi như thế này?

Nó rơi ra từ đằng sau một mảnh giấy đã ố màu thời gian. Và nó trông không giống món đồ bị ai đó làm rơi vô tình. Mà giống như\dots\ được đặt vào đúng chỗ đó, để chờ ai đó tìm thấy.

\dots Có thể chỉ là trùng hợp. Nhưng trực giác mách bảo tôi rằng tôi nên giữ nó lại.

Không rõ vì sao, nhưng từ sâu trong lòng, tôi có cảm giác rằng thứ dây đeo này quan trọng hơn vẻ ngoài của nó rất nhiều. Nó gợi tôi nhớ tới một cái gì đó. Một cái gì đó đang nằm ngoài tầm với của ký ức tôi hiện giờ.

Tôi lặng lẽ bỏ nó vào túi quần.

``Hy vọng mình không vừa ăn cắp đồ của ai đó\dots'' tôi lẩm bẩm, nửa đùa nửa thật. Saikyou mỉm cười nhẹ, rồi nhẹ nhàng đặt mảnh giấy trở lại chỗ cũ như thể chúng tôi chưa từng động vào.

``Ta nên đi tiếp thôi nhỉ?'' tôi nói. Em tôi gật đầu. Và thế là chúng tôi lại bước tiếp, qua hành lang phủ bụi và ánh sáng mờ ảo, những dấu chân nhẹ in lên nền gạch cũ kỹ.

Nhưng dù vậy, những thông tin từ mảnh giấy đó vẫn luôn thôi thúc tôi về vụ việc đã xảy ra ở đây. Một vụ nổ khí ga khiến cả trường sập đổ\dots\ Nhưng tại sao chúng tôi lại ở đây? Trong một ngôi trường đã bị phá hủy?

``Vụ nổ\dots'' tôi khẽ thì thầm.

``Bị sập\dots'' Saikyou-chan lặp lại, giọng đều đều nhưng hơi rùng mình.

Những từ đó cứ vang vọng trong đầu tôi như một âm vang xa xăm từ một giấc mơ cũ mà tôi không tài nào nhớ nổi.

\dots Rốt cuộc nơi này đã xảy ra chuyện gì?

Một hiện tượng siêu nhiên? Hay chỉ là một tai nạn?

Chúng tôi vẫn chưa có câu trả lời nào. Và có thể, trước khi tìm ra sự thật, chúng tôi còn phải đối mặt với nhiều bí ẩn hơn thế.

``Chúng ta cần tập trung vào việc thoát khỏi nơi này,'' tôi tự nhủ. ``Không cần quan tâm đến mấy chuyện linh tinh này.''

Nhưng rõ ràng, những câu chuyện ``linh tinh'' ấy lại cứ tìm đến chúng tôi\dots

``Nii-chan, em tìm thấy một mảnh giấy nữa này,'' Saikyou-chan bất ngờ lên tiếng, cắt ngang dòng suy nghĩ của tôi.

Tôi quay đầu lại. Em ấy đang đứng trước một chiếc bảng thông báo treo lệch trước cửa một phòng học, tay cầm một mảnh giấy được ghim cẩn thận bằng một chiếc đinh rỉ sét.

``{\vnmsans Ghi chú của nhà báo - tờ số 1,}'' em ấy đọc to, mắt vẫn không rời khỏi tờ giấy.

Một bản ghi chép của\dots\ một nhà báo từng đến đây?

``Anh có nghĩ là thứ này có liên quan tới đoạn bài báo mà ta đọc khi nãy không?'' em hỏi.

Tôi không chắc. Có thể. Mà cũng có thể chẳng liên quan gì. Nhưng dù gì, nếu nó chứa thông tin, thì cũng đáng để xem qua.

``Để xem họ viết gì\dots'' tôi lẩm bẩm.

\txtbx{
  {\arial "Về việc xảy ra những hiện tượng kỳ lạ ở khu vực vùng núi - Thị trấn Aogami.\\
    Tôi đã được (cấp trên) phân công đến vùng này để thu thập thông tin.\\
    Tất cả chỉ vì một lời nói vu vơ của tổng biên tập mà tôi bị lệnh đưa tới đây. Phiền toái thật.\\
    \dots Kiểu, họ bắt tôi phải tìm hiểu một câu chuyện hão huyền nhảm nhí, rồi còn đếch thèm đặt bản thân của họ vào vị trí của tôi. Như thế này chẳng phải quá quắt lắm sao?"}}

Tôi nhíu mày.

Người này rõ ràng đang bực bội với nhiệm vụ mà họ nhận được. Trông không giống một bản ghi chép chuyên nghiệp cho lắm, mà giống như một bản trút giận cá nhân hơn. Nhưng ít nhất, nó xác nhận một điều: trước khi chúng tôi đến đây, đã có người từng điều tra nơi này. Và người đó cũng từng hoài nghi, giống như chúng tôi bây giờ.

``Tội thật,'' tôi lẩm bẩm. ``Đã thế còn bị cấp trên sai đi vào một chỗ quái gở như thế này nữa.''

``\dots Nếu mà là em chắc em sẽ cãi lại họ rồi đấy,'' Saikyou lắc đầu, mồ hôi rịn trên trán.

``Chắc chắn ai cũng sẽ muốn làm vậy thôi,'' tôi cười nhẹ. ``Nhưng mà, phận làm cấp dưới là vậy. Trừ khi họ muốn bị đá khỏi chỗ làm.''

``C-Cũng phải\dots'' em ấy đáp, cười khì một tiếng.

Chúng tôi rẽ sang phải, theo hướng hành lang dẫn tới một cánh cửa trượt. Em gái tôi bước tới trước, khẽ kéo nhẹ cánh cửa.

*XOẠCH*

``Ồ, cái cửa này mở được này!'' em reo lên, giọng có phần ngạc nhiên.

Lạ thật. Nếu cửa này mở được\dots\ chẳng lẽ cánh cửa bị khóa ban nãy lại chính là cửa ra?

\dots Hoặc cũng có thể, cả hai đều không phải.

\dots Hoặc cả hai đều có liên quan.

Chúng tôi nhìn nhau một lát, rồi em tôi bước hẳn vào trong. Còn tôi, vẫn đứng lại một chút, liếc nhìn mảnh dây đồng hồ trong túi quần qua một khe hở nhỏ.

Một cảm giác kỳ lạ chạy dọc sống lưng tôi, rằng chiếc dây này, rồi đây, sẽ mở ra cánh cửa dẫn tới một điều gì đó quan trọng hơn tất cả những gì tôi từng biết.

\ct

Trước mắt chúng tôi là một hành lang khác.

Rất giống với hành lang mà chúng tôi vừa đi qua.

Vẫn là những ô cửa sổ dài khung gỗ, vẫn là những lớp học không đánh số, vẫn là ánh sáng đèn huỳnh quang mờ nhạt phủ một lớp trắng bệnh lên mặt sàn. Vẫn là những chiếc loa rè treo trên cao, những kệ tủ sát tường đã khép chặt, không một chi tiết hay âm thanh nào khác.

Mọi thứ không khác một ly.

Và chính vì vậy, cảm giác bất an ngày càng lớn dần trong tôi. Có cái gì đó\dots\ sai sai.

Một ngôi trường bình thường có thể có vài hành lang giống nhau, điều đó đúng. Nhưng ở đây, cái sự \emph{giống nhau} ấy lại quá mức hoàn hảo. Như thể chúng tôi không hề di chuyển chút nào, mà chỉ đang giậm chân tại chỗ trong một vòng lặp vô tận.

Tôi rùng mình.

``Em thấy mình tự dưng bị nổi da gà\dots'', Saikyou khẽ nói, càng áp sát vào người tôi hơn. Tôi có thể cảm nhận rõ hơi thở của em ấy đang dồn dập hơn bình thường. Không rõ là do trời lạnh hay là do sự hiện diện của cái gì đó bất thường đang bám riết lấy chúng tôi.

\dots Lộp cộp, lộp cộp\dots

Tiếng bước chân vang lên khe khẽ. Một hành lang dài mà vắng vẻ, âm thanh bước chân như phản lại chính mình, gợn gợn như vọng ra từ một không gian rỗng không.

Những lớp học tiếp tục trống rỗng, không một bóng người. Tất cả những gì còn sót lại trong đó là những bộ bàn ghế xếp ngay ngắn, những bức tường trắng không treo bất kỳ tấm bảng nào, và cửa sổ thì bị che phủ bởi những tấm rèm tối màu.

Tôi bắt đầu có cảm giác đây không phải là một nơi bị bỏ hoang, mà đúng hơn là một nơi được cố tình giữ nguyên hiện trạng. Không ai sống, nhưng cũng không ai dọn dẹp, không ai chạm vào, không có một vết bụi nào. Như thể nó được bảo quản trong trạng thái tĩnh.

Tôi thở dài. Thật sự không có ai ở đây sao?

``\dots Nii-chan,'' em khẽ gọi. Tôi ngoảnh lại.

``Liệu chúng ta có thật sự\dots\ ở một chiều không gian nào đó mà không có lấy một bóng người không?''

Tôi đã suýt cười, suýt nữa thì buông một câu đùa rằng em ấy đang xem quá nhiều phim. Nhưng tôi không thể. Với bầu không khí này, mọi câu nói đều hóa thành tiếng vọng rỗng.

\dots Biết đâu em ấy lại đúng?

Tôi chưa kịp nghĩ thêm, thì em gái tôi lại lên tiếng: ``Với cả, em tìm được thêm một thứ tài liệu bên cửa sổ này.''

Em ấy giơ lên một mảnh giấy khác, gấp tư, như vừa được lượm lên từ kẽ giữa hai tấm kính. Khuôn mặt em tỏ rõ sự lo lắng. Không giống kiểu tò mò như trước nữa. Có vẻ, càng đi sâu vào hành lang này, những thứ em nhặt được càng khiến em rối bời.

``{\vnmsans Tin nóng: Hiện tượng kỳ lạ ở thị trấn Aogami - tin thứ nhất}'' em đọc to tiêu đề.

Tôi khựng lại. Lại là Aogami?

Thị trấn nơi ngôi trường này tọa lạc. Bài báo này giống như một mảnh ghép khác được ai đó cố tình rải rác dọc lối đi. Như thể chúng tôi đang lần theo dấu vết của chính một kẻ nào đó từng ở đây trước chúng tôi.

``Để em đọc nó lên nhé?'' cô em gái của tôi nói. Tôi gật đầu.

Nhưng chỉ vài giây sau, em ấy nhăn mặt, xoay tờ giấy lại, cố gắng nhìn thật kỹ các dòng chữ.

``Nii-chan, em không đọc được mấy chữ trên đây.''

``Lại gì nữa?''

``Phần ghi tháng và năm bị ai đó gạch bung bét đây này.''

Giống hệt với mẩu tin đầu tiên.

Cứ như thể\dots\ có người cố tình che giấu thời gian xảy ra sự kiện.

``Chắc là em sẽ đọc tháng năm thành các chữ X ha,'' Saikyou hắng giọng, rồi bắt đầu đọc:

\txtbx{{\arial
    "Tháng X, năm 2018.\\
    Điều tra bí ẩn, phát hiện một vật thể bay không xác định tại thị trấn UFO.\\
    Tại thị trấn Aogami, một vật thể bay không xác định (UFO) đã được người dân địa phương nhìn thấy trên bầu trời vào lúc đêm khuya.\\
    Khi nhắc đến UFO, người ta sẽ thường nghe kể tới những câu chuyện liên quan tới việc sẽ đến bắt và đưa con người đi. Mặc dù đây chỉ là những câu chuyện được truyền miệng, tuy nhiên, cũng không thể phủ nhận rằng nó có thể liên quan tới số lượng người mất tích bí ẩn đang tăng lên tại thị trấn này."}}

Một đĩa bay à?

Ngay giữa một thị trấn nhỏ như Aogami?

Nghe như cốt truyện của một bộ phim giả tưởng hạng ba chiếu khuya. Bạn biết đấy, UFO, đĩa bay, người ngoài hành tinh xâm chiếm Trái Đất, bắt cóc con người làm thí nghiệm, vân vân mây mây?

Nhưng việc nó xuất hiện ở thế giới thực ư? Thật sự phải rất hiếm lắm mới có thể xảy ra đấy.

\dots Hoặc, đây có lẽ là một trò đùa nào đó của một đứa ất ơ vào giữa đêm để dọa nạt cư dân tại đây? Một dạng ``thí nghiệm xã hội''? Một ai đó đã dựng nên kịch bản hoang đường này để khiến người ta hoảng loạn?

Không. Không giống lắm.

Cái sự im lặng chết chóc trong hành lang, những tờ giấy cũ rích được giấu kỹ, ánh sáng nhợt nhạt, hành lang lặp lại như bản sao\dots\ tất cả đều gợi nên một điều gì đó cố ý. Có lẽ, nó còn vượt xa cả trò đùa.

Tôi nhìn sang Saikyou. Em ấy cũng trầm ngâm, tay vô thức nắm lấy tờ giấy.

``Em nghĩ chắc chúng ta không nên bận tâm quá về chuyện này đâu nhỉ?''

Tôi gật đầu. Dù lòng vẫn còn hoài nghi, nhưng mục tiêu của chúng tôi lúc này là tìm lối ra. Những mẩu tin kia, nếu thực sự có giá trị, có lẽ sẽ tự động dẫn đường cho chúng tôi.

Em gái tôi nhẹ nhàng đặt lại mảnh giấy vào đúng nơi em nhặt được, rồi cả hai chúng tôi tiếp tục bước vào lòng hành lang trước mắt.

\ct

Chúng tôi cứ thế đi tiếp, dọc hành lang dài ngập trong ánh sáng mờ đục của đèn trần và ánh sáng bên ngoài cửa sổ. Không ai nói lời nào. Dường như ngay cả tiếng bước chân của chúng tôi cũng bị hút vào một lớp không khí dày đặc vô hình.

Bất chợt, một thứ ánh sáng lóe lên từ phía trước khiến chúng tôi khựng lại theo phản xạ. Đó là một ánh sáng yếu, nhấp nháy từ chiếc bảng thông báo treo trên vách tường của một lớp học, cũng trống trơn như tất cả những căn phòng khác. Ánh sáng ấy đến từ một màn hình nhỏ, có vẻ như được lắp gắn tạm bợ và cũ kỹ, phát ra một ánh sáng trắng xanh lạnh lẽo.

Trên bảng là một mẩu giấy, giống như trước đó.

``{\vnmsans Ghi chú của nhà báo - tờ số 2}'' tôi đọc to.

Lại là một tờ ghi chú nữa. Tay nhà báo bí ẩn ấy dường như đã để lại các mẩu giấy khắp nơi trong ngôi trường này như thể muốn dẫn dắt người đọc lần theo dấu vết. Liệu lần này, hắn sẽ tiết lộ điều gì nữa đây?

\txtbx{\arial "Trường trung học Aogami. Một ngôi trường phổ thông dự bị đại học."}

``Dự bị đại học?'' tôi lặp lại. ``Ý hắn là nơi này giống như một kiểu trường luyện thi, với mục tiêu chuẩn bị cho học sinh vào đại học?''

``T-Thế thì khác gì mấy cái lò luyện thi ở ngoài đâu,'' Saikyou lẩm bẩm.

Tôi tiếp tục đọc phần còn lại của ghi chú:

\txtbx{\arial "Đây là một khu vực mà cảnh sát thường tuần tra xung quanh, nhưng dạo gần đây lại nhiều người đến nỗi làm tôi phát sợ.\\
  Theo như tôi tìm hiểu, gần đây có rất nhiều vụ mất tích, mà chủ yếu là các nữ sinh. Nghe thật ảo ma thế nào ấy.\\
  Tôi biết rằng việc phỏng vấn các học sinh nơi đây cũng chẳng hề thoải mái gì, nhưng tôi sẽ cố gắng hết sức để tránh sự nghi ngờ về phía chúng ta."}

Tôi nhìn chăm chăm vào những dòng chữ, cố gắng nắm bắt sự thay đổi trong giọng điệu của người viết. Có một cái gì đó không ổn\dots

``Tờ ghi chú này\dots\ không biết có cùng một người viết ra với tờ trước không,'' tôi lẩm bẩm. ``Nhưng nếu có, thì dường như hắn đã thực sự bước vào điều tra\dots\ và có vẻ như không hoàn toàn tự nguyện.''

``Tránh sự nghi ngờ về phía chúng ta'\dots?'' Saikyou-chan nhắc lại một cách khó hiểu. ``Nii-chan, em thấy hắn không chỉ điều tra\dots\ mà còn có dính líu gì đó thì phải\dots''

Rồi em ấy im bặt. Một thoáng sau, cô bé bỗng rùng mình, mặt tái mét như chợt nhận ra điều gì đó đáng sợ.

``Nhưng mà\dots\ nếu những vụ mất tích ở đây chủ yếu là các nữ sinh\dots'' Em ấy càng lúc càng run rẩy, hai tay bấu chặt lấy vạt áo tôi. ``Em không muốn bị chia cắt với anh đâu\dots!''

Tôi đưa tay xoa nhẹ đầu em. Lòng bàn tay tôi cảm nhận rõ sự run rẩy từ làn tóc mềm của em ấy.

``Không sao đâu, Saikyou-chan,'' tôi thì thầm. ``Anh ở đây mà.''

Chúng tôi quay lại với hành lang. Trước mặt là một cánh cửa trượt khác. Tôi nắm tay nắm, cố kéo\dots

*LẠCH CẠCH*

Cánh cửa không hề nhúc nhích.

``Chậc\dots\ không phải cánh nào cũng mở được rồi,'' Tôi rít nhẹ một tiếng.

Thật quá quắt. Không chỉ bị nhốt ở nơi này, không rõ tại sao, không nhớ được chuyện gì, mà còn không thể di chuyển tự do. Như thể có ai đó đang giật dây sau tấm màn sân khấu mà chúng tôi là những con rối không tự chủ.

Tôi quay sang định nói gì đó với Saikyou-chan thì\dots

``N-N-Nii-chan\dots'' em ấy khẽ thốt lên, bàn tay siết lấy tay áo tôi. ``C-Có\dots\ ai đó\dots\ ở trong lớp học lúc nãy\dots''

Tôi quay phắt lại. Qua ô cửa kính mờ của lớp học mà chúng tôi vừa đọc ghi chú, một bóng người đang đứng bên trong.

Cơn lạnh chạy dọc sống lưng tôi.

Đó là một cô gái.

Tôi chắc chắn là lớp học ấy trống rỗng, tôi đã nhìn vào nó trước đó rồi. Nhưng giờ đây\dots\ có ai đó đang ở trong đó, đứng bất động như một cái bóng.

\img{e0b}

Cô gái ấy đứng quay mặt ra phía cửa sổ, hơi nghiêng người. Dáng người thon gọn trong bộ đồng phục học sinh, với một dải nơ cài lệch sang bên trái. Mái tóc dài xõa xuống gần đến thắt lưng, nhưng màu sắc và chi tiết hoàn toàn bị nuốt chửng trong một màu đen đặc sệt, như thể cô ấy chỉ là một hình nhân bị cắt bóng từ một mảnh giấy đen.

``\hl{Giống quá\dots}'' tôi thì thầm, tim đập thình thịch. ``\hl{Giống như cô gái mà mình đã nhìn thấy lúc mới tỉnh lại\dots}''

Tôi đang định cất tiếng gọi thì bỗng nhiên\dots

*XOẸT!*

Cơ thể cô gái bị nhiễu, giống như hình ảnh trong một trò chơi điện tử bị lỗi hình ảnh, vỡ vụn thành những dòng nhiễu điện tử hóa.

Và rồi\dots\ cô ấy biến mất. Không một dấu vết.

Chỉ còn lại khoảng trống lạnh lẽo giữa lớp học, và cơn ớn lạnh đang lan dần nơi gáy tôi.

*XOẠCH*

Một tiếng động khác vang lên từ phía sau. Cánh cửa trượt ở hành lang mà chúng tôi thử mở lúc nãy tự động bật mở.

Không ai chạm vào nó.

``S-Sợ quá, Nii-chan\dots\ Nơi này giống như bị ám vậy\dots!'' Saikyou-chan run lên, gần như muốn bật khóc.

Tôi nuốt khan. Thực sự, thứ tôi cảm thấy không phải là nỗi sợ thông thường.

Mà là một loại trực giác nguyên thủy: có thứ gì đó đang quan sát chúng tôi từ trong bóng tối.

Liệu nơi này chỉ đơn giản là một ngôi trường bỏ hoang?

Hay là một thứ gì hoàn toàn khác?

\pov{Haragen Saikyou}

Chúng tôi hối hả bước qua cánh cửa trượt vừa mở ra, để rồi lại thấy mình đang ở trong một hành lang khác không khác. Cũng cái ánh đèn lạnh lẽo ấy. Cũng cái mùi ẩm mốc, cái cảm giác nghẹt thở mà chẳng hiểu tại sao cứ luẩn quẩn nơi đây. Thành thực mà nói, tôi chẳng thấy nổi một sự khác biệt nào giữa hành lang này với cái hành lang trước đó cả. Mọi thứ dường như đang lặp lại, như thể có một bàn tay vô hình nào đó đang dàn dựng lại từng bước chân chúng tôi.

Nhưng có một điều chắc chắn không giống như trước: lần này, không ai còn bước chậm nữa. Sau những gì vừa chứng kiến, những bước chân chúng tôi đã chuyển sang gấp gáp, hấp tấp, như thể chỉ cần chậm lại một giây thôi, thứ gì đó sẽ vồ lấy mình từ phía sau.

Tôi vẫn quàng chặt tay vào anh trai, tay bám cứng lấy ống tay áo của anh như một cái phao cứu sinh. Tôi không muốn buông ra. Không dám. Sau cái bóng người trong lớp học vừa rồi, ai mà chẳng sợ?

Bình thường thì tôi không phải kiểu người nhát gan như thế này. Tôi mạnh mẽ lắm chứ. Thậm chí không ngán mấy đứa con trai côn đồ dám đụng tới anh em chúng tôi. Nhưng ở cái nơi khỉ ho cò gáy, tối om và toàn sương mù này, tôi chẳng thể ngăn được hai chân mình run cầm cập.

Vì những hiện tượng kỳ quái à? Hay vì cái cảm giác vô hình, đè nén, như thể có ai đó đang lặng lẽ dõi theo từng cử động của mình?

Hay là bởi vì, tất cả nạn nhân trong các vụ mất tích kỳ bí tại ngôi trường này đều là nữ sinh\dots\ như tôi?

Tôi lắc đầu, tự gạt đi cái suy nghĩ đó. Không! Không được nghĩ linh tinh nữa! Phải rời khỏi nơi này. Phải thoát ra.

Chúng tôi bước càng lúc càng nhanh, bỏ mặc phía sau cái bóng tối đang dần nuốt trọn mọi thứ. Không ai nói gì. Chỉ có tiếng bước chân vọng lại trên nền gạch lạnh và\dots\ và một điều gì đó không ổn.

Tôi suýt bật hét lên khi một cái bóng vụt qua ngoài cửa sổ. Nhanh đến mức tôi không kịp nhận ra đó là thứ gì, chỉ biết nó trông như một hình người\dots\ không, như một con ma!

``\textbf{Gyaaah!!}'' tôi hét toáng lên, cả người bất giác nhào vào người anh trai.

Nii-chan dừng lại, quay đầu sang tôi với vẻ hoảng hốt: ``Gì thế!?''

Tôi run rẩy chỉ tay về phía ô cửa kính phủ đầy sương mờ, thốt lên: ``M-Ma kìa! Em vừa thấy nó!''

Anh ấy lập tức nhìn theo hướng tôi chỉ. Một thoáng im lặng. Rồi anh cau mày lại: ``Có thấy gì đâu?''

Tôi biết, bình thường anh hay trêu tôi những lúc thế này. Nhưng lần này, tôi thấy anh cũng hơi lưỡng lự. Sau những gì chúng tôi thấy lúc trước, chính anh cũng chẳng còn đủ bình tĩnh để đùa cợt như mọi khi nữa.

Tôi bám chặt tay áo anh hơn.

\dots Mà\dots\ có phải là do tôi không, mà sao tôi còn nghe thấy\dots

*LẠCH CẠCH*

\dots

*LẠCH CẠCH*

Âm thanh ấy. Rất nhỏ. Như tiếng móng tay gõ vào kính vậy. Tôi quay phắt lại. Không thấy gì. Nhưng âm thanh vẫn cứ vẫn cứ ở đó.

Tôi đâu có tưởng tượng chứ?

Không, không\dots\ nơi này thật sự có gì đó rất sai.

Cảm giác như\dots\ mình đang bị theo dõi. Bị săn đuổi.

Tôi nuốt khan.

Tôi rất muốn kéo tay Nii-chan chạy, chạy thật xa khỏi cái hành lang chết tiệt này. Nhưng\dots\ chạy đi đâu bây giờ? Tất cả mọi chỗ đều giống nhau. Và nếu tôi rời khỏi anh ấy, nhỡ đâu\dots\ tôi sẽ trở thành người tiếp theo thì sao?

Không được. Không thể được. Dù có chuyện gì xảy ra, tôi cũng sẽ không bao giờ rời xa anh ấy. Không bao giờ!

\dots

\dots Phải bình tĩnh lại. Bình tĩnh, Saikyou.

Không được để nỗi sợ xâm chiếm.

Chúng tôi sẽ ổn thôi. Chỉ cần bám sát Nii-chan, từng bước, từng bước\dots

Rồi chúng tôi sẽ rời khỏi nơi này.

Và sau đó\dots\ sau đó\dots

\dots Chúng tôi sẽ\dots

\dots

\dots

Ờm\dots\ chúng tôi sẽ\dots\ làm gì nhỉ?

Thật lòng mà nói, tôi cũng chẳng biết nữa. Có lẽ\dots\ Có lẽ cách tốt nhất bây giờ là cứ ở gần anh ấy mọi lúc.

Vì chỉ cần tôi còn nắm lấy tay anh\dots\ thì ít nhất tôi sẽ không biến mất.

Tôi thấy Nii-chan khựng lại ở một góc phòng, rồi cúi xuống nhặt lên một mẩu giấy đã ngả màu. Anh ấy lật qua một lượt rồi lẩm bẩm: ``Hừm\dots\ Lại là một bài báo nữa.''

Tôi buông anh ấy ra, tò mò: ``Đâu? Cho em xem với!''

Anh đưa tờ giấy đó cho tôi, vẻ mặt đăm chiêu. Tôi ngước mắt đọc lướt qua hàng tiêu đề: ``{\vnmsans Vụ việc Trường Trung học Aogami bị sập: Ký sự số 2.}''

Tôi nhận ra đây là phần tiếp theo của bài báo mà chúng tôi từng nhặt được ở hành lang đầu tiên. Không rõ ai đã để lại những tờ này ở khắp nơi như thế, và cũng chẳng có ngày tháng gì được ghi kèm để xác nhận tính liên tục, nhưng bằng trực giác, tôi cảm thấy chúng liên quan đến nhau.

Tôi liếc nhìn Nii-chan một cái, rồi khẽ nói: ``Em đọc cho anh nghe nhé.''

Tôi hắng giọng, cố gắng phát âm rõ ràng từng chữ giữa ánh sáng mờ ảo:

\txtbx{\arial Vụ việc ngôi trường trung học tại thị trấn Aogami gần đây đã dấy lên nhiều lời đồn đại từ phía cư dân địa phương. Tuy nhiên, vẫn còn vô số khúc mắc chưa được lý giải.
  Theo một số nhân chứng, nguyên nhân không đến từ sự cố nổ khí gas như báo cáo chính thức. Có ý kiến cho rằng\dots}

Tôi dừng lại. Phần tiếp theo gần như không thể đọc được nữa, những dòng chữ ngoằn ngoèo như nét bút bị nghẽn mực, vết bẩn loang lổ chồng lên nhau khiến tôi không tài nào nhận ra nổi. Thật khó chịu.

Tôi khẽ nhăn mặt, gấp tờ giấy lại, lẩm bẩm: ``Đúng là báo lá cải\dots\ Chẳng có lấy một thông tin cụ thể. Viết kiểu này thì làm sao mà giải mã được chuyện gì đang xảy ra ở đây cơ chứ.''

Tôi đưa trả mảnh giấy cho anh trai mình, rồi vừa đặt nó lại lên kệ tủ vừa lầm bầm thêm vài câu chửi thề trong đầu. Mấy tờ giấy ngu ngốc này chỉ tổ làm tôi bối rối hơn thôi.

Thế rồi, tôi ngẩng lên, nói khẽ: ``Nii-chan, đi tiếp thôi. Đừng ở đây nữa.''

Không đợi anh trả lời, tôi đã vội sải bước về phía trước, như thể muốn giễu cợt lại nỗi sợ của chính mình. Một phần tôi muốn chọc anh ấy một chút cho bớt căng thẳng, nhưng có lẽ sâu bên trong, tôi chỉ đang cố trốn chạy khỏi cảm giác bất an đang dần len lỏi khắp cơ thể mình.

Thật lạ. Mới nãy tôi còn bảo phải bám sát lấy anh ấy không rời. Mà giờ lại bỏ đi trước? Tôi có điên rồi không? Điên thật mất.

Và tôi đã sai. Sai hoàn toàn.

Chỉ mới bước ra được vài bước\dots\ khoảng ba bước gì đó thôi\dots

``Hi hi\~{}''

Tôi đông cứng người lại. Có gì đó không đúng. Tôi nghe rõ tiếng cười. Rất khẽ, rất đều, nhưng vọng ra từ một lớp học.

Tôi vội vàng quay ngược trở lại, ôm chầm lấy cánh tay Nii-chan theo phản xạ, thì thào: ``C-Có tiếng cười\dots\ ai đó đang cười ở trong phòng kia!''

Cả hai chúng tôi rón rén bước đến gần cánh cửa phòng học đang khép hờ. Tôi cố len người qua khe cửa nhìn vào bên trong.

Tôi thấy họ.

Năm cô gái. Đứng rải rác ở những vị trí khác nhau trong phòng, không ai nhìn ai, không ai di chuyển, chỉ khẽ\dots\ cười. Mỗi người trong số họ mang một dáng vẻ kỳ dị đến rợn người, như thể được khắc họa ra từ một cơn ác mộng không có hồi kết.

\img{e0c}

Một người trong số đó đứng sát bên khung cửa sổ, mái tóc dài che khuất gần hết khuôn mặt, nhưng không giống hệt với ``nó'' mà chúng tôi đã thấy ở lớp trước. Kiểu tóc có phần khác biệt.

Một người khác có hai bím tóc nhỏ trông như đang lơ lửng, đang nghiêng đầu nhìn ra bên ngoài như thể đang thì thầm với gió.

Một cô gái ngồi gần đó có đôi tai\dots\ cáo? Hay tai mèo? Tôi không biết. Tôi cũng chẳng chắc mình có thực sự nhìn thấy chúng không nữa. Hoặc cũng có thể chỉ là bờm để trang trí tóc.

Gần cô ấy là một cô gái có mái tóc cột lệch một bên, đứng xoay lưng lại phía tôi, bất động.

Và cuối cùng, một người có mái tóc xù, xõa tung với hai búi tóc nhỏ nhô lên như sừng. Cô ấy đang quay mặt về phía bảng đen.

Tất cả đều cười, không một lời, không một tiếng thì thầm, không một cử động thừa.

Tôi nín thở. Nii-chan cũng vậy.

*XOẸT*

Chớp mắt một cái, tất cả biến mất.

Không tiếng động. Không khói mờ. Không gió lùa.

Chỉ còn lại một căn phòng trống, bụi phủ đầy bàn ghế, như thể chưa từng có ai ở đó.

Tôi nhìn lại tay mình vẫn đang bám chặt lấy cánh tay anh trai. Dù tôi không run, nhưng tôi biết, trái tim tôi đã đánh rơi một nhịp. Hay là mấy nhịp liền.

``Thôi nào, Nii-chan, ta đi tiếp thôi-''

Tôi quay lại gọi anh trai, định kéo tay anh ấy để tiếp tục bước qua khu vực hành lang u tối và lạnh ngắt kia, nhưng\dots

Tôi sững người.

Nii-chan đang khom người, ôm đầu như thể vừa bị ai đó đập một cú cực mạnh. Anh thở dốc từng cơn, mồ hôi túa ra trên trán, rồi bất ngờ ngã quỵ xuống sàn.

``Đau\dots\ quá\dots'' anh rên lên, giọng lạc đi.

Tôi hoảng loạn quỳ sụp xuống bên cạnh, luống cuống đỡ lấy vai anh ấy: ``N-Nii-chan!? Anh bị sao vậy!?''

Tôi cố giữ lấy đầu anh, lo lắng nhìn vào khuôn mặt đang nhăn nhó vì đau đớn. Không giống với những lần chóng mặt nhẹ thông thường. Nii-chan trông như thể đang vật lộn với một cơn đau trong tâm trí, không chỉ ở thể xác.

Anh ấy siết lấy cánh tay tôi, rồi khẽ nói, giọng lạc đi: ``Không biết tại sao\dots\ nhưng\dots''

Tôi vội đáp: ``Nhưng sao?''

``\dots Anh thấy\dots\ bốn người bọn họ\dots\ quen lắm. Anh đã từng gặp họ ở đâu đó rồi\dots''

Tôi khựng lại.

Gì cơ? Nii-chan đã gặp họ từ trước?

Không thể nào. Không thể nào như thế được.

Tôi và anh là hai anh em ruột, suốt ngày đi đâu cũng có nhau, kể cả lúc mua kem hay đi xem phim ma. Còn mấy cô gái đó, rõ ràng tôi chưa từng gặp họ trong đời! Không lẽ\dots?

Tôi ngước nhìn anh trai, hơi siết chặt tay lại. Không lẽ anh ấy đã gặp họ lén lút ở đâu đó mà không nói với tôi?

Hay tệ hơn\dots\ họ là bạn anh ấy, một nhóm riêng biệt nào đó mà tôi không được phép biết đến?

Tôi bặm môi, trong lòng dấy lên một cảm giác khó chịu mơ hồ.

Nii-chan đúng là xấu tính mà. Tôi cũng muốn có bạn chơi cùng nữa cơ\dots\ Hức\dots

Nhưng mà không phải lúc để nghĩ tới chuyện đó.

Hít sâu một hơi, tôi gạt đi những cảm xúc hỗn độn đang cuộn trào trong ngực. Những thứ vừa rồi không bình thường một chút nào cả. Cái cách mà anh ấy phản ứng, cái cách anh ấy ôm đầu\dots

Không lẽ đây cũng là một loại hiện tượng tâm linh?

Anh ấy\dots\ bị ảnh hưởng từ những hồn ma kia sao?

``\dots Anh dám chắc là\dots\ họ cũng nằm trong số những người đã mất tích\dots'' Nii-chan nói khẽ, giọng run rẩy, khi anh cố đứng dậy với một tay vẫn còn giữ lấy thái dương.

Tôi nhanh chóng đưa tay đỡ lấy anh để anh không ngã xuống lần nữa. Có lẽ tôi sẽ phải tra hỏi anh kỹ hơn về chuyện ``quen biết lén lút'' này sau.

``Anh nói vậy là vì sao?'' tôi hỏi, thấp giọng.

Anh chỉ lắc đầu. Không trả lời. Im lặng.

Có lẽ chính anh cũng không hiểu rõ cảm giác ấy đến từ đâu. Không chắc chắn, chỉ là một sự quen thuộc mơ hồ, như một đoạn ký ức lướt ngang qua giấc mơ, không thể níu lại, không thể gọi tên.

Là ai? Là ở đâu? Khi nào?

Trong lúc tôi đang bối rối vì điều đó, một vật gì đó trên bảng thông báo phía cửa lớp lọt vào mắt tôi.

Tôi buông tay anh ra một chút, chỉ vào đó: ``Em lại tìm thấy một mẩu giấy khác rồi, dán trên bảng kìa.''

Dù trông vẫn còn hơi choáng, Nii-chan cũng lảo đảo bước tới tấm bảng cùng tôi. Tôi thấy ánh mắt anh lướt qua dòng chữ trên tờ giấy, lẩm bẩm đọc lên: ``{\vnmsans Ghi chú của nhà báo - Tờ số 3.}''

Tôi thở dài. Lại là tay nhà báo kia nữa.

Nhưng lần này, khi tôi bắt đầu đọc đoạn nội dung ghi trên đó, một điều kỳ lạ xảy ra, gã không còn viết với cái giọng châm biếm hay cay cú như những lần trước nữa.

\txtbx{\arial "Có vẻ như, tôi cảm nhận được rằng các học sinh trong ngôi trường sẽ biến mất sau một ngày cụ thể theo một cách nào đó.\\
  Khi đó, tôi có thể cảm nhận được nỗi sợ hãi của bọn chúng vào thời điểm điều đó xảy ra.\\
  Hơn nữa, tôi được biết rằng tại ngôi trường này, có một lời đồn được truyền tai nhau rằng:\\
  'Khi một học sinh từ nơi khác được chuyển tới, những điều xui rủi sẽ xảy ra.'\\
  Dường như cứ vài năm, lại có một học sinh chuyển trường đến đây, và điều đó khiến cả trường xôn xao. Những lời đồn đại cứ thế mà lan rộng."}

Tôi im lặng trong vài giây sau khi đọc xong.

Lần này, không có những câu than phiền vô nghĩa. Không có cái giọng điệu cà khịa hay mỉa mai nữa.

Tay nhà báo kia dường như đã nghiêm túc rồi.

Và đó là lúc đầu tôi bắt đầu xoay mòng mòng với những câu hỏi dồn dập:
\begin{center}
  ``Tại sao những học sinh đó lại biến mất?''\\
  ``Lời đồn đó xuất phát từ đâu?''\\
  ``Người chuyển trường là ai?''\\
  ``Tại sao `người chuyển tới' lại bị xem là dấu hiệu của xui xẻo?''\\
  ``Và tất cả những điều này có liên quan gì đến việc ngôi trường này bị sập, nếu không phải vì vụ nổ khí gas như báo cáo chính thức?''\\
\end{center}

Tôi bắt đầu cảm thấy rằng những gì chúng tôi đang trải qua vượt xa những trò dọa ma thông thường. Đây không còn đơn thuần là hiện tượng tâm linh nữa.

Mà là một chuỗi mắt xích đen tối, sâu sắc và cố tình.

Một điều gì đó đã xảy ra, và nó vẫn chưa kết thúc.

\dots Mà, có lẽ chúng tôi cũng không thể biết được sự thật. Chúng tôi còn chẳng rõ vì sao bản thân mình lại đang ở đây, trong ngôi trường này, bị bủa vây bởi những ký ức không thuộc về mình.

Nii-chan vẫn đang ôm lấy đầu mình, nhưng giờ anh ấy đã có thể thở một cách bình thường. Anh ấy lắc đầu một cách dứt khoát, rồi\dots\ trời ạ, còn làm mấy động tác khoa trương như mấy nhân vật trong phim hoạt hình dành cho trẻ con, kiểu lắc đầu như chong chóng để ``làm bay hết đau nhức'' ấy.

Nii-chan đúng là hết thuốc chữa.

``Đầu anh thế nào rồi?'' tôi hỏi, nhẹ đặt tay mình lên trán anh ấy.

Anh ấy thở dài một tiếng, rồi đáp lại: ``Ừ\dots\ Anh nghĩ là ổn rồi.''

``Thật không đấy?''

``\dots Chắc thế. Đầu vẫn còn hơi nhói một chút.''

``Ồ\~{} Mà này, em vẫn chưa xử anh vụ mấy cô bạn mà anh bảo là quen quen đâu đấy\~{}'' Tôi vờ nghiêm nghị, khoanh tay nhìn anh ấy với ánh mắt nghi ngờ. Dù sao thì, tôi là em gái mà, tôi cũng có quyền tra hỏi chứ nhỉ?

Anh tôi thở ra một tiếng như thể muốn xua đi phiền phức, cố giải thích: ``Đã bảo là anh cảm giác như đã từng thấy họ ở đâu đó rồi\dots\ Anh còn không chắc là anh đã gặp họ khi nào nữa là.''

Tôi khịt mũi, nửa đùa nửa thật. ``Ừ, em biết mà. Nii-chan sẽ không bao giờ phản bội lại cô em gái nhỏ đáng yêu này đâu nhỉ?'' tôi châm chọc, rồi huých khuỷu tay vào người anh một cái không nhẹ chút nào.

``Ái da- Đau đấy!'' Nii-chan kêu lên, lùi hẳn lại phía sau một bước. Tôi cười khoái chí. Có ai nói với bạn rằng trêu chọc người khác cũng là một kiểu thể hiện tình thân không? Không ác ý đâu, chỉ là một chút ``trả đũa nhẹ nhàng'' thôi.

``Rồi, rồi, anh xin. Vậy chúng ta đi tiếp chứ?'' anh nói, xoa sườn, lắc đầu ngao ngán.

``Rất sẵn lòng!'' Tôi lại níu lấy cánh tay anh, như một cái móc sắt nhỏ dính chặt vào thân tàu, rồi cả hai chúng tôi tiến về phía cánh cửa trượt tiếp theo.

Nhưng mà\dots\ ủa? Sao cánh cửa này lại bị che phủ bởi một lớp đen sì vậy? Tôi không thể nhìn thấy bất cứ thứ gì ở phía bên kia cả. Có phải ai đó đã dùng giấy đen, hay băng dính gì đó để bịt kín lại không?

\dots Thôi kệ, chứ còn cách nào khác đâu.

Tôi đẩy cánh cửa ra.

*XOẠCH*

Cánh cửa trượt vẫn mở ra bình thường, không có gì kỳ lạ. Nhưng ngay khi bước qua đó\dots

\dots một hành lang mới hiện ra trước mắt chúng tôi.

\img{e0d}

Về mặt cấu trúc thì không có gì quá khác biệt: tường gạch, sàn gỗ, cửa lớp học ở bên trái, và các ô cửa sổ ở bên phải đều có mặt. Vẫn ở vị trí đó, không thay đổi.

Nhưng có gì đó\dots\ rất sai.

Trần nhà võng xuống, như thể đang bị thứ gì đè nặng lên. Những vết rạn chân chim chạy dài khắp nơi, và phần mép trần thì tróc từng mảng, trông như thể chỉ cần một cái thở mạnh cũng đủ khiến nó sụp xuống.

Ngay khi cả hai chúng tôi mới bước được một bước\dots

*RẦM!* *RỘP\dots\ RỘP\dots*

Một loạt các tấm trần rơi phịch xuống, làm tung bụi mù mịt. Cả hai anh em tôi ho sặc sụa.

``*Khụ\dots\ Khụ khụ!* Có vẻ như chỗ này\dots\ đang gặp vấn đề lớn đấy\dots'' Nii-chan nói, vừa che miệng vừa nghiêng đầu quan sát mớ trần vừa rơi xuống.

``*Khụ khụ\dots\ Em cũng nghĩ vậy\dots'' tôi đáp, dùng tay phẩy phẩy không khí trước mặt để đuổi đám bụi bay tứ tung kia đi. Cổ họng tôi ngứa rát, mắt cay xè.

''Chắc là ta phải đi đứng cẩn thận thôi. Không chừng chỗ này có thể đổ sập bất cứ lúc nào,'' anh ấy nói thêm.

''\dots Hy vọng là không phải như vậy, Nii-chan.''

Không thể tin nổi đây là một phần của cùng ngôi trường. So với những hành lang trước, dù có vẻ lạnh lẽo và u ám, thì ít nhất vẫn còn nguyên vẹn, kiên cố và không tạo cảm giác như sắp sập đến nơi.

Còn nơi này\dots\ giống như một vùng đất bị lãng quên. Hoang tàn. Vặn vẹo.

Tôi tự hỏi: liệu những gì chúng tôi đang thấy có liên quan đến vụ sập trường mà báo đã đưa tin? Có thể nào, đây là dấu tích, hay là một biểu hiện khác của ``\emph{thảm họa}'' đó, theo đúng nghĩa đen và cả nghĩa bóng?

Nghĩ đến đây, tôi rùng mình. Không biết là do bụi, do lạnh, hay là do nỗi sợ đang thấm dần vào từng tế bào.

Chúng tôi bắt đầu bước nhanh hơn, không nói một lời nào. Từng bước chân nghe rõ mồn một trong hành lang vắng lặng như tờ, nơi mà cái chết có thể rơi xuống từ trần bất cứ lúc nào.

Tôi ngó nhìn hai bên. Tất cả các lớp học đều bị che kín bằng thứ gì đó giống như màn đen. Không thể thấy được gì phía sau. Và cũng không thể mở ra. Giống như một phần thực tại đã bị khóa chặt lại, không cho phép bất kỳ ai bước vào.

Nơi này không giống như những hành lang bình thường. Nó không chỉ là hư hỏng. Nó đang mục nát. Nó đang chết dần.

Và tôi\dots\ tôi bắt đầu sợ rằng mình sẽ bị chôn sống cùng với nó.

Bất chợt, một tia sáng nhỏ lướt qua tầm mắt, lóe lên nhẹ từ góc kệ tủ bên tay trái, khiến cả hai chúng tôi cùng khựng lại.

Tôi nghiêng người, cố xác định nguồn ánh sáng đó. Là ánh phản chiếu\dots\ từ một mảnh giấy có in mực bóng? Hay là từ băng keo? Dù là gì đi nữa, nó vẫn nằm đó, nổi bật trên nền bụi phủ bạc màu xung quanh.

Chúng tôi tiến lại gần.

\dots Lại là một mẩu giấy nữa. Nhìn thoáng qua, tôi gần như thở dài. Dạo này tôi thấy bản thân bắt đầu quen với việc cứ vài chục bước lại vấp phải một mẩu giấy kỳ lạ nào đó.

Tôi nhặt nó lên, thở ra một tiếng chán nản. Dòng tiêu đề được in đậm bằng mực đỏ: ``{vnmsans Tin nóng: Hiện tượng kỳ lạ ở thị trấn Aogami - tin thứ hai.}''

Chậc. Lại là cái trang tin tức vớ vẩn đó, cái trang mà lần trước đã từng la toáng lên về đĩa bay với chẳng UFO. Đã vậy, phần ngày tháng ở đầu trang còn bị gạch nguệch ngoạc như thể ai đó cố tình che giấu, hoặc xóa bỏ điều gì đó.

Tôi chẹp miệng một tiếng, đưa mẩu giấy lên sát mắt. Tôi dám chắc là nó sẽ chỉ khiến cho đống câu hỏi trong đầu tôi càng thêm rối tung hơn. Nhưng rồi vẫn đọc:

\txtbx{\arial "Một con tàu hỏa bí ẩn đã được nhìn thấy tại ga Aogami vào đêm qua.\\
  Một số người dân nói rằng họ chưa từng nghe tiếng còi tàu nào to đến thế trong đời.\\
  Tuy nhiên, cho đến hiện tại, sự xuất hiện của đoàn tàu này vẫn chưa có lời giải thích chính thức."}

Tôi khựng lại một lúc, đọc lại đoạn cuối.

Tàu hỏa? Sau đĩa bay, giờ là tàu hỏa?

Xuất hiện rồi biến mất. Không manh mối. Không ai chịu trách nhiệm. Không báo cáo chính thức.

Thị trấn này\dots\ đang ngày càng giống một ảo ảnh méo mó của chính nó. Mọi thứ như đang diễn ra song song với thực tại mà chúng tôi biết, nhưng không trùng khớp với bất kỳ logic thông thường nào.

Tôi định quay sang để hỏi Nii-chan xem anh có thấy thứ này vô lý như tôi không, thì chợt phát hiện anh ấy cũng đang cầm một tờ giấy khác. Ánh mắt tôi mở to.

``Anh lấy nó ở đâu vậy?'' tôi hỏi, giọng ngạc nhiên lộ rõ.

``Trên bảng tin kia kìa,'' anh trả lời, tay chỉ về phía bảng thông báo gắn trước một lớp học bên cạnh. ``Anh nghĩ anh cũng vừa tìm thấy một thứ đáng để suy nghĩ.''

Tôi bước lại gần anh. ``Nó viết gì vậy?''

Anh ngước nhìn tôi, rồi đọc lớn:

\txtbx{\vnmsans "Vụ việc trường trung học Aogami bị sập: Ký sự số 3."\\
  \arial "Một số người dân tại thị trấn Aogami đã lập luận rằng có một trận động đất lớn xảy ra đúng vào thời điểm ngôi trường sụp đổ.\\
  Tuy nhiên, theo dữ liệu giám sát từ Cục Khí tượng Quốc gia, hoàn toàn không có ghi chép nào về một cơn địa chấn tại khu vực này vào thời điểm đó.}

Tôi nuốt nước bọt. Một cảm giác lạnh tràn lên dọc sống lưng.

``\dots Cái này thì thực sự là kỳ quái đấy,'' tôi lẩm bẩm. ``Một trận động đất không được ghi nhận? Làm sao chuyện đó có thể xảy ra?''

Tôi bắt đầu nghĩ nhanh. Có ba khả năng mà tôi có thể hình dung ra ngay lúc này.

Thứ nhất, có lẽ thứ mà người ta cho là động đất chỉ là một chấn động thứ cấp xảy ra sau vụ sụp đổ, và vì nó không đủ mạnh nên không được ghi nhận là động đất chính thức. Điều này có thể xảy ra.

Thứ hai, động đất thực sự đã xảy ra, nhưng vì lý do nào đó, dữ liệu về nó đã bị gỡ bỏ hoặc bị che đậy. Một ai đó, hoặc một tổ chức nào đó đã nhúng tay vào. Nếu là như thế thật\dots\ thì vụ việc này thậm chí còn nghiêm trọng và đen tối hơn nhiều so với chúng tôi tưởng.

Và\dots\ thứ ba.

Tôi không muốn tin vào điều này. Tôi không phải kiểu người tin vào những thứ siêu nhiên. Nhưng\dots

\dots còn nhớ cái mảnh ghi chú trước đó, nơi có đề cập đến một lời nguyền? Nếu như điều đó là thật, nếu như nơi này thực sự đang bị điều khiển bởi một thế lực nào đó vượt ra ngoài sự hiểu biết của con người\dots

\dots thì phải chăng, tất cả những gì đang xảy ra đều là một phần của kế hoạch? Một chuỗi những sự kiện đã được sắp đặt từ trước?

Một chiếc đĩa bay. Một đoàn tàu biến mất. Một trận động đất không tồn tại. Những nữ sinh mất tích. Một ngôi trường bị sập không rõ lý do. Và giờ là chính chúng tôi, đang lang thang trong chính nơi đó, không lối thoát.

``\textit{Ai là người đang điều khiển tất cả chuyện này?}'' tôi lặng lẽ nghĩ. ``\textit{Là người? Là `nó'? Hay là cái gì đó còn khủng khiếp hơn?}''

Chúng tôi cẩn thận đặt hai mẩu giấy trở lại chỗ cũ, như thể sợ rằng bất kỳ thay đổi nhỏ nào cũng có thể làm lệch đi một mắt xích của chuỗi hiện tượng này.

Cả hai lặng im trong vài giây.

Chúng tôi không biết mình đến đây bằng cách nào. Không biết những gì đã xảy ra trước đó là thật hay ảo. Những hiện tượng xoay quanh những cô gái mất tích, những bài báo lờ mờ, những mẩu giấy tản mát, tất cả cứ như một bức tranh bị xé vụn, rải khắp nơi trong ngôi trường này.

Và bức tranh ấy, nếu có thể ghép lại, sẽ kể cho chúng tôi nghe một câu chuyện.

Một câu chuyện mà có lẽ, không phải ai cũng muốn nghe.

Nhưng lúc này, điều duy nhất mà chúng tôi cần làm\dots\ là thoát khỏi nơi này.

Tôi nắm chặt tay áo anh trai mình, và cả hai cùng nhau tiếp tục tiến về phía trước, trong một ngôi trường mà từng bước chân vang vọng như lời cảnh báo lặp đi lặp lại trong đầu: ``\emph{Không phải mọi cánh cửa đều nên được mở.}''

Nhắc đến cánh cửa, chúng tôi nhìn thấy có một cánh cửa trượt đã được mở hé.

``Cơ mà, tại sao nó lại hé nhỉ?'' tôi nghiêng đầu hỏi.

``Anh cũng chẳng biết. Có vẻ như ai đó quên không khóa lại,'' Nii-chan đáp, rồi nhún vai.

``\dots Chắc là không vấn đề gì đâu nhỉ? Được cái là chúng ta không cần tốn công mở nó ra\dots''

Trong lòng hơi hí hửng, chúng tôi vội vã chạy lại gần cánh cửa đang mở. Thế nhưng\dots

*XOẠCH*

\dots một tiếng động khô khốc vang lên. Ngay khi chúng tôi vừa tới nơi, cánh cửa trượt đột ngột đóng sầm lại, như thể có ai đó từ bên trong vừa giật nó về.

\gif{e0e}{1}{36}

Tôi và anh trai khựng lại.

``Nó\dots\ nó tự đóng lại rồi kìa\dots'' tôi lắp bắp.

``Chậc, lại là một cái gì đó ma quái nữa đây,'' Nii-chan gắt lên, vẻ khó chịu.

Tôi bước đến trước cánh cửa vừa đóng, đưa tay kéo thử nó. ``Hnnngh\dots''

*LẠCH CẠCH\dots\ LẠCH CẠCH\dots*

Không được. Nó bị khóa chặt rồi. Không, không phải khóa nữa, mà là như thể nó đã hóa thành một phần của bức tường vậy. Không có lấy một khe hở nào để luồn móng tay vào.

``Chết tiệt\dots'' tôi thì thầm. ``Không thể nào. Bây giờ\dots\ chẳng lẽ lại phải quay ngược trở lại?''

Tôi ngó quanh. Cái hành lang phía sau dài hun hút, tối đen và trống rỗng. Chúng tôi đã đi quá xa rồi. Tôi không chắc liệu đường đó có còn an toàn nữa không.

``Vậy thì\dots\ cửa sổ thì sao?'' Nii-chan nói, ngó ra cái khung kính bên cạnh lớp học.

Tôi nhìn theo ánh mắt anh ấy. Cửa sổ trông cũng bình thường thôi, nhưng cái gì ở đây mà bình thường chứ?

``\dots Thôi thì cứ thử vậy,'' tôi thở dài.

Tôi lùi lại một khoảng lấy đà, rồi quay sang nói: ``Nii-chan, tránh ra một bên.''

Anh ấy nghe lời, nhích sang bên phải. Tôi hít một hơi thật sâu.

``\textbf{Hiyaa!!}''

Tôi lao đến cánh cửa sổ, tung chân đá thật mạnh vào khung kính\dots

*RẦM!*

\dots Một tiếng động lớn vang lên.

Tôi bật ngược trở lại, ngã chỏng gọng, mông đáp thẳng xuống sàn.

``Đau quáaaa\dots'' tôi rên lên, tay xoa chỗ vừa va phải. Nii-chan vội cúi xuống đỡ tôi dậy, phủi lớp bụi bám trên váy tôi.

``Có vẻ như cửa sổ cũng không khả thi nhỉ.''

``Khỉ thật chứ\dots\ Cái cửa sổ này làm bằng kính chống đạn à?'' tôi càu nhàu.

Mồ hôi bắt đầu rịn ra trên trán tôi, lạnh toát. Vậy thì còn đường nào để thoát nữa?

Tôi đang vò đầu bứt tai thì chợt nghe tiếng anh trai: ``\dots Anh thấy có một mẩu báo nữa dán trên lớp học kia,'' Nii-chan chỉ về phía bên trái hành lang.

Tôi quay lại nhìn theo, thấy một tờ giấy ngả màu vàng úa được dán chặt trên cửa một lớp học gần đó.

Tôi cắn môi, cáu tiết. ``\textbf{Giờ không phải lúc đọc báo nữa đâu! Mau cùng em nghĩ cách thoát ra khỏi đây đi!}''

Tôi không chịu đựng được nữa rồi.

Tôi ghét cái cảm giác mắc kẹt này.

Tôi không muốn ở lại đây thêm một phút nào nữa.

Tôi muốn rời khỏi nơi quỷ quái này!

\dots

\dots

Tôi không muốn trở thành một nạn nhân khác trong ngôi trường này! Làm ơn\dots\ ai đó\dots\ mau cứu chúng tôi với!

\pov{Haragen Kyuen}

Tôi nhìn cô em gái mình đang co rúm lại trong hoảng loạn, mái tóc rối tung, tay vò đầu, mắt mở to như sắp bật khóc đến nơi. Em ấy đang thật sự sợ hãi. Không cần phải là người thân thiết mới nhận ra điều đó, nhìn một lần là đủ biết.

Saikyou\dots\ đã thật sự bắt đầu phát điên rồi.

Cũng chẳng thể trách được. Từng bước đi trong cái hành lang quái đản này đều như đang dẫn chúng tôi vào một cơn ác mộng có thật. Từng mẩu báo, từng cánh cửa, từng cái bóng mơ hồ phía cuối hành lang\dots\ Tất cả đều như được dựng sẵn, như thể có một bàn tay vô hình nào đó đang điều khiển mọi chuyện.

Nếu là tôi trong hoàn cảnh của em ấy, tôi có khi còn gào thét lên từ lâu rồi.

Tôi bước đến bên em ấy, kiễng người xuống, đặt hai tay lên vai em nhẹ nhàng. ``Được rồi, Saikyou\dots\ Chúng ta sẽ tìm cách rời khỏi đây nhanh thôi,'' tôi nói, giọng nhỏ nhưng dứt khoát.

Em ấy ngẩng lên, đôi mắt vẫn đẫm đầy vẻ sợ hãi. ``T-Thật chứ?''

Tôi ngần ngại. Nếu là trả lời thật lòng, thì tôi không chắc. Tôi cũng bắt đầu cảm thấy không ổn. Không, phải nói là vô cùng bất an. Không chỉ vì Saikyou\dots\ mà vì chính tôi cũng có cảm giác như mình đang bị theo dõi.

Thậm chí, trong thoáng chốc, tôi có cảm giác rằng tôi cũng đang là mục tiêu của cái gọi là ``sự kiện mất tích'' đó.

Tôi chưa từng là người mê tín. Nhưng bầu không khí nơi này, sự kỳ lạ quá đỗi của từng chi tiết nhỏ nhặt - cái cách mà mọi thứ xảy ra theo một thứ trật tự vô hình nào đó - khiến cho sống lưng tôi lạnh toát.

Không. Tôi không thể để Saikyou thấy được điều đó. Là một người anh, ít nhất tôi cũng phải là chỗ dựa cho em ấy, dù bản thân có đang run lẩy bẩy bên trong.

Tôi siết nhẹ vai em ấy và mỉm cười: ``Ừ. Hy vọng là như thế.''

Tôi đỡ em ấy đứng dậy. Như một phản xạ, Saikyou lại bám chặt lấy cánh tay tôi.

``Sợ quá thì cứ bám lấy tay anh nha?'' tôi nói, cố trấn an em bằng một nụ cười ấm áp. Em ấy không đáp, chỉ khẽ gật đầu, nhưng như vậy cũng là đủ rồi.

Tôi quay đầu nhìn xung quanh. Cánh cửa trượt thì đã đóng chặt như biến thành một phần của bức tường. Cửa sổ thì chắc chắn chẳng khác gì một bức tường kính bọc thép. Hành lang thì kéo dài bất tận, tối om như cổ họng một con quái vật đang há rộng. Chúng tôi\dots\ quả thật, đang mắc kẹt.

``Trước mắt, anh nghĩ nên đọc thử mẩu giấy kia đã,'' tôi nói, chỉ vào tờ giấy dán trên bảng xanh hồi nãy. ``Biết đâu có manh mối gì đó giúp ta hiểu chuyện này hơn một chút.''

Em tôi hơi nhíu mày, rõ ràng vẫn chưa muốn nghĩ gì ngoài chuyện thoát thân. Nhưng cuối cùng, em ấy cũng gật đầu. Vì thật ra, chúng tôi đâu còn lựa chọn nào khác.

Tôi bước tới gần, nhẹ tay gỡ mảnh giấy đã ngả màu ra khỏi tấm bảng. Dòng tiêu đề lờ mờ: ``{\vnmsans Vụ việc trường trung học Aogami bị sập: Ký sự số 4.}''

Chỉ có đúng hai dòng, nhưng lạnh sống lưng:

\txtbx{\arial Những điểm bất thường trong vụ việc trường trung học bị sập càng được nêu bật ra. Một trong số đó là việc bản danh sách học sinh đăng ký học tại đây lại không trùng khớp với danh sách các học sinh bị mất tích.}

Danh sách học sinh\dots\ không trùng khớp với danh sách mất tích?

Một là bản danh sách đăng ký ban đầu bị sai sót.

Hai là danh sách học sinh mất tích đã bị cố ý làm sai lệch.

Hoặc, có thể còn một khả năng thứ ba. \emph{Rằng những người ``không nằm trong danh sách học sinh'' vẫn có thể mất tích.}

Tôi bỗng thấy da gà nổi khắp người.

Như tôi chẳng hạn. Tôi không phải là học sinh trường này. Tôi chỉ đi theo em gái mình, bị kéo vào cái hành lang quái gở này, vô tình mặc bộ đồng phục này, và thậm chí còn chẳng nhớ nổi bản thân mình là ai. Vậy thì\dots\ nếu tôi cũng biến mất, liệu có ai biết không?

Liệu chúng tôi có trở thành ``người không tồn tại'' như bao nhiêu cái tên đã từng bị quên lãng ở đây?

Tôi cảm giác hơi thở mình chậm lại. Không thể để Saikyou thấy tôi đang run, nhưng trong lòng tôi là một cơn sóng dữ. Cảm giác bất an lúc đầu giờ đã lớn hơn nhiều.

Tôi siết nhẹ tay em gái mình. Em ấy ngẩng lên nhìn tôi, như để chắc rằng tôi vẫn còn ở đây. Tôi gật đầu, như để trấn an em\dots\ và cả chính mình.

``Đi thôi,'' tôi nói. ``Nếu những mẩu giấy này xuất hiện liên tục, thì có thể ai đó đang cố gửi gắm điều gì đó cho chúng ta.''

\dots Hoặc đang dẫn chúng tôi đến gần hơn với cái bẫy kế tiếp.

*RỘP!*

Một âm thanh chát chúa bất thình lình vang lên trong không gian im phăng phắc khiến cả hai anh em tôi giật bắn người.

Chúng tôi quay phắt lại.

Một cánh cửa trượt, chính là cửa của lớp học nơi tấm bảng có đính bài báo vừa rồi, đang mở toang.

\dots Khoan đã\dots\ chẳng phải tôi nhớ rõ ràng là cánh cửa này, cùng với các lớp trước đó đã bị khóa từ trước rồi sao?

Tôi liếc nhìn Saikyou. Em ấy cũng đang nhìn cánh cửa với ánh mắt run rẩy, như thể nó vừa tự động mở ra là chuyện đương nhiên.

Dù gì thì, cửa đã mở. Đồng nghĩa với việc chúng tôi có thể vào xem xét bên trong lớp học ấy. Có lẽ sẽ tìm được điều gì đó có ích. Một manh mối, một đường lui, một ký ức còn sót lại?

Tôi nghĩ vậy.

\dots Nhưng em gái tôi lại không nghĩ thế.

``C-Cái gì vậy, Nii-chan\dots\ Đáng sợ quá\dots\'' Em run rẩy, nắm chặt tay tôi không rời, thân người như dán chặt vào tôi. Em ấy vừa mới lấy lại chút bình tĩnh sau những cú sốc liên tiếp. Chấn động tâm lý chưa hề nguôi ngoai, giờ lại thêm một hiện tượng kỳ quặc như thế này\dots\

Tôi cúi xuống, cố giữ bình tĩnh trong giọng nói:
``Có lẽ\dots\ lớp học này sẽ cho ta một manh mối gì đó để thoát khỏi đây?''

``Em không nghĩ vậy đâu\dots\''

Tất nhiên rồi. Còn ai tin nổi mấy thứ ``manh mối'' trong cái nơi bị ma ám này nữa?

Tôi thở dài. ``Vậy em nghĩ chúng ta còn cách nào khác sao?''

Saikyou-chan khựng lại.

Cặp mắt em ấy đảo quanh hành lang một cách vô thức, như đang tìm kiếm một con đường khác. Nhưng không có gì cả ngoài những bức tường đóng kín, bóng tối mơ hồ và\dots\ cánh cửa duy nhất vừa tự động mở ra như đang chờ đón chúng tôi bước vào.

Tôi biết em ấy đang tuyệt vọng. Tôi cũng vậy.

Và nếu chúng tôi cứ đứng đây mà không làm gì\dots\ thì cũng chỉ là đang chờ tới lượt mình bị "mất tích".

``\dots Chúng ta vào thôi,'' tôi nói, cố giữ giọng vững vàng hơn vẻ mặt đang dần cứng lại.

``Sẵn sàng chưa?''

Em ấy không trả lời, chỉ khẽ gật đầu một cái như sắp khóc. Tôi nuốt khan.

``Vậy\dots\ ta vào nha.''

Từng bước, từng bước một, chúng tôi rón rén tiến về phía cánh cửa lớp học đang mở hé.

Khoảng cách giữa chúng tôi và cái lỗ hổng kia ngày một ngắn lại. Không khí nơi đây trở nên đặc quánh, như thể đang bước vào một lớp sương vô hình. Tim tôi đập dồn từng nhịp một, chẳng rõ là vì sợ hãi, hay vì một linh cảm bất tường đang lớn dần trong lòng.

Và rồi, ngay khi chúng tôi chỉ còn cách cánh cửa một bước chân\dots

*ROẠT*

Một bóng người bất ngờ ló đầu ra từ bên trong lớp học.

Tôi và Saikyou giật nảy người, theo phản xạ lùi thụt về phía sau.

Đó là một cô gái.

\gif{e0f}{1}{72}

Mái tóc dài màu nâu óng (ít nhất là dưới ánh sáng le lói ở đây), xõa tới tận lưng, trên đỉnh đầu cột hờ bằng một chiếc nơ đỏ sậm. Đôi mắt xanh lam lạnh lẽo nhìn xoáy vào chúng tôi, gương mặt không biểu cảm, như thể hoàn toàn trống rỗng. Mà cũng có thể là\dots\ quá bình thản để mà đáng sợ.

Cô ấy mặc một bộ đồng phục học sinh đen tuyền, khác hẳn với kiểu đồng phục của trường tôi. Thiết kế đơn giản, cổ áo buộc khăn xanh nước biển cổ điển, mờ nhạt, và\dots\ gần như đã thuộc về một thời đại khác.

Tôi cau mày. Dáng người cô ấy nhỏ, gọn, hơi nghiêng về phía trước\dots\ Như một con lật đật biết đi. Làm tôi nhớ đến cô gái mà tôi định gọi tới ngay sau khi tôi đứng dậy.

Saikyou thì thào: ``Cái quái gì thế\dots? Cô là ai\dots?''

Không có tiếng trả lời. Cô ta chỉ đứng đó, nhìn trừng trừng.

Một sự im lặng khủng khiếp bao trùm.

Cô ấy không cử động. Chẳng có biểu cảm nào. Cũng không chớp mắt.

Hay là\dots\ cô ấy đang cười? Một nụ cười nhếch mép, mơ hồ đến mức khiến tôi nổi da gà. Tôi nhíu mắt cố nhìn rõ hơn.

*RỘP!*

Cánh cửa đóng sầm lại trước mắt chúng tôi, nhanh và bất ngờ như lúc nó mở ra.

Ngay trong khoảnh khắc ấy\dots

*Ù Ù Ù\dots*

Một cơn địa chấn đột ngột ập tới.

Cả hành lang rung chuyển dữ dội. Từ trên trần nhà, những mảng thạch cao dột nước đổ xuống ầm ầm, bụi bay mù mịt, cuộn lên như một cơn sương mù trắng xóa, phủ trùm lên mọi thứ.

Không kịp suy nghĩ gì thêm, tôi quay sang ôm chặt lấy cô em, hét lớn: ``\textbf{Chạy! Chạy thôi!!}''

Tôi biết rõ đây là cách \emph{ngu ngốc} nhất để đối phó với một cơn động đất. Nhưng trong cơn hoảng loạn như thế này, lựa chọn duy nhất chính là bản năng sinh tồn.

``C-Chuyện gì vậy, Nii-chan!?''

``Anh không biết! Nhưng\dots\ nhìn kìa!!''

Cánh cửa trượt cuối hành lang, nơi Saikyou từng cố gắng mở trước đó bỗng nhiên bật mở.

Không cần thêm một giây suy nghĩ, cả hai anh em tôi cắm đầu chạy xuyên qua làn bụi, phóng thẳng về phía cánh cửa vừa hé.

Ngay khi vừa lách mình qua được khung cửa, tôi đập mạnh nó lại phía sau, đóng chặt cửa thật mạnh.

*RẦM!*

Cánh cửa khép lại đúng lúc mặt đất ngừng rung.

Cơn động đất\dots\ bỗng chốc dừng hẳn.

Phía sau cánh cửa, chỉ còn lại một sự im lặng chết người.

Tôi và Saikyou đứng thở hổn hển trong bóng tối của hành lang mới, không ai nói gì.

Bầu không khí lại lặng như tờ, như mọi chuyện ban nãy chỉ đơn thuần là một cơn ác mộng ngắn. Nhưng tôi biết rõ, đó mới chỉ là bắt đầu.

Chúng tôi thở dốc, tim đập như muốn nhảy khỏi lồng ngực, mắt dán chặt vào khung cảnh phía sau, nơi hành lang vừa bị tàn phá đến không thể nhận ra.

Các ô trần lả tả rơi như sung chín, bụi mù bay trắng xóa cả không gian, còn phần trần nhà thì sụp đổ từng mảng, để lộ ra những khung thép gỉ sét, vặn vẹo như những ngón tay của một con quỷ khổng lồ đang cố vươn ra khỏi lớp bê tông.

Hai từ thôi: kinh hoàng. Và điều khiến tôi sợ hãi hơn cả, là cái cảm giác\dots\ đây chưa phải là tận cùng.

``\textbf{Vừa rồi là cái gì vậy, Nii-chan!?}'' Saikyou hét lên, gần như sắp bật khóc.

``Không lẽ nào\dots\ đó là động đất? Nhưng mà\dots'' tôi ấp úng, cũng không thể tìm ra một lời giải thích nào hợp lý hơn.

``Tại sao\dots\ lại là chúng ta chứ!?'' giọng em tôi run rẩy. Tôi không biết trả lời thế nào.

Không khí đang dần trở nên ngột ngạt. Bụi từ những mảng tường rơi tạo nên mùi khét nghẹt của bê tông bị ép nát trộn lẫn với mùi gỉ sắt, mùi ẩm mốc, và cả\dots\ một chút gì đó tanh tưởi mà tôi không muốn nghĩ đến.

Tôi nuốt nước bọt. Một suy nghĩ bỗng lóe lên trong đầu, có thể đây chính là khoảnh khắc ngôi trường Aogami năm xưa đổ sập.

Và nếu đúng như vậy\dots

\textbf{*Ù Ù Ù!!!*}

Trước khi tôi kịp nói ra giả thuyết của mình, một đợt dư chấn khác lại ập tới. Lần này còn dữ dội hơn cả cơn trước.

Sàn nhà rung chuyển dữ dội đến mức tôi cảm thấy cả cơ thể mình như bị nhấc bổng lên rồi thả rơi lại xuống sàn. Saikyou mất đà, ngã quỵ xuống sàn, may mà tôi kịp đỡ em ấy lại.

``Saikyou! Em không sao chứ!?''

``\textbf{Cái gì đang xảy ra vậy hả!?!}'' giọng em tôi vỡ ra vì hoảng loạn.

Tôi gào lên giữa cơn chấn động đang ngày càng dữ dội: ``Không biết! Nhưng mà\dots\ \textbf{phải chạy! Phải ra khỏi đây thôi!!}''

Tôi nắm chặt lấy tay em gái và kéo em chạy về phía trước. Không còn thời gian để suy nghĩ nữa. Bản năng sinh tồn đã lấn át lý trí.

Từng mảnh kính từ cửa sổ vỡ tung rơi như mưa, lấp lánh trong ánh sáng lờ mờ như hàng ngàn lưỡi dao sắc lạnh.

Những tấm trần cuối cùng cũng rơi rụng sạch sẽ, để lại những thanh sắt nhọn hoắt đang lơ lửng giữa không trung như gọng vuốt của quái vật.

Tệ hơn nữa, mặt sàn bắt đầu nứt toác. Những vết nứt như miệng cá mập há ngoác ra, từ từ nuốt chửng mặt đất. Dưới đó, chỉ có một màu đen đặc quánh, không đáy, không hình dạng, như địa ngục đang há miệng.

Tôi nắm chặt tay Saikyou-chan hơn, kéo em né qua từng lỗ sập, từng khe nứt đang lan ra với tốc độ chóng mặt.

*KẸT\dots\ RẦM!! XOẢNG!!*

Tiếng thép gãy. Tiếng trần nhà đổ. Tiếng kính vỡ. Tất cả hòa vào nhau như một bản giao hưởng hỗn loạn, giao hưởng của sự tận diệt.

``\textbf{Kia kìa! Cửa ra!}'' tôi hét lên, chỉ tay về phía cánh cửa trượt phía cuối hành lang.

Saikyou-chan la lên thất thanh: ``\textbf{N-Nii-chan! Phía trước có mấy cái hố to lắm!}''

Tôi nhìn xuống, thấy mặt sàn chỉ còn lại một dải chéo hẹp, nối giữa hai hố sâu ngòm đen đặc.

``Không có cách nào khác! Nhảy thôi!!''

Chúng tôi nhảy. Mắt tôi chỉ thấy khoảng trống dưới chân, cảm giác như đang rơi tự do giữa không trung. Chúng tôi vừa kịp tiếp đất bên kia, và ngay sau đó, dải sàn duy nhất ấy cũng sụp đổ ngay sau lưng.

Thở dốc. Tôi quay đầu lại nhìn hành lang vừa thoát ra khỏi, giờ trông chẳng khác gì một bãi chiến trường đổ nát.

Cả hành lang chẳng còn gì nguyên vẹn. Gạch đá, kính vỡ, cốt thép lòi ra như xương, và những hố đen vô tận vẫn còn thở phì phò như thể chưa no máu.

Lớp học bên cạnh cũng chẳng khá hơn: bàn ghế nằm lộn xộn, đồ treo trần thì hoặc sụp hẳn, hoặc treo lủng lẳng như xác chết chưa kịp tháo xuống.

\begin{figure}[H]
  \centering
  \begin{subfigure}[t]{0.45\textwidth}
    \centering
    \includegraphics[width=8cm]{../../../../Image/Book?/e0g.png}
  \end{subfigure}
  \quad
  \begin{subfigure}[t]{0.45\textwidth}
    \centering
    \includegraphics[width=8cm]{../../../../Image/Book?/e0h.png}
  \end{subfigure}
\end{figure}

``\dots H-Hết chưa vậy\dots?'' tôi thì thầm, giọng run lẩy bẩy.

Saikyou cũng gật đầu thẫn thờ: ``C-Có vẻ là rồi đó, Nii-chan\dots\ Em cứ tưởng em sắp chết đến nơi\dots''

Tôi nén một tiếng chửi thề trong cổ họng. ``\dots Định mệnh cái cơn động đất khỉ gió này chứ!''

``Nhưng ít nhất thì\dots\ chúng ta cũng đến được cánh cửa, đúng không?''

Em gái tôi chỉ về cánh cửa trượt trước mặt, nơi duy nhất chưa bị nuốt chửng. Nhưng bên kia\dots\ chỉ là một màu đen sâu hút, như thể nó là ngưỡng cửa giữa thế giới này và một cõi khác.

``C-Chắc\dots\ chúng ta phải đi tiếp thôi.'' Tôi lẩy bẩy bước về phía nó.

``Đ-Để em mở cửa\dots!''

Thế nhưng, ngay khi tay em chạm vào cánh cửa\dots

\dots Một loạt tiếng cười khúc khích vang lên.

Tiếng cười con gái. Lảnh lót. Éo éo. Mảnh và dai như kim khâu xuyên qua màng nhĩ.

Nó không chỉ vọng từ bên trong lớp học. Mà là khắp nơi. Trên trần. Dưới sàn. Cả từ bên trong bức tường.

Toàn thân chúng tôi đông cứng lại. Rồi\dots

*Ù Ù Ù\dots!!!*

Một cơn địa chấn khác lại bắt đầu.

``\textbf{Cẩn thận!!}'' tôi hét lên, lấy thân mình che cho Saikyou-chan. Một tấm trần rơi trúng vai tôi. Nhẹ thôi, nhưng cũng đủ để khiến tôi rùng mình.

*RẮC!*

Tiếng gãy vỡ. Ngay dưới chân chúng tôi, mặt sàn bắt đầu nứt ra lần nữa.

``T-Ta phải làm gì bây giờ, Nii-chan!?''

``\textbf{Mở cửa! Nhanh lên!!}''

Nhưng đã quá muộn.

*RẮC\dots\ RẮC\dots!!*

\textbf{*RẦM!!!*}

Cả mảng sàn nơi chúng tôi đang đứng đổ sập, không ai trong chúng tôi có thể phản ứng kịp.

Tôi và em gái tôi cùng hét lên trong tuyệt vọng:

``\textbf{Nii-chan!!}''

``\textbf{Saikyou!!}''

Chúng tôi rơi vào hố đen sâu thẳm, nơi ánh sáng bị nuốt chửng như chưa từng tồn tại.

Tiếng cười ma quái vẫn vang vọng, quấn lấy hai anh em tôi như dây gai xiết cổ.

Ánh sáng mờ nhạt phía trên ngày càng xa\dots\ càng nhạt\dots\ rồi tắt hẳn.

Và khi bóng tối nuốt trọn lấy mọi thứ\dots

Tôi chỉ còn một suy nghĩ trong đầu: Phải chăng\dots\ đây là cái chết?

Chết\dots\ ở một nơi mà chúng tôi còn không biết nó là gì, không biết ai tạo ra, không biết vì sao mình bị cuốn vào\dots

Một cái chết\dots\ phi lý và trớ trêu đến tuyệt vọng.

Phải chăng\dots\ cuộc hành trình của hai anh em chúng tôi\dots\ đã kết thúc\dots

\dots ngay khi vừa mới bắt đầu?

\dots

\dots

\dots

\begin{center}
  \uline{???}
\end{center}

\dots

Khụ khụ\dots

C-Chuyện quái gì vừa diễn ra thế này\dots

\dots

À phải rồi. Trận động đất đó. Cái lỗ sập ấy.

Tôi tưởng mình đã chết cơ. Nhưng mà\dots

\dots Tôi vẫn còn sống.

Tôi vẫn có thể suy nghĩ, thở được, cảm nhận được. Toàn thân ê ẩm, nhưng ít nhất tôi vẫn có thể cử động tay chân.

Chỉ có điều\dots\ tôi chẳng thấy gì cả. Trước mắt tôi là một không gian mờ mờ, chỉ lấp loáng ánh sáng trắng lạnh lẽo như ánh huỳnh quang hắt xuống từ trần.

Tôi cố gắng lồm cồm bò dậy, tay dụi mắt để nhìn rõ hơn. Không. Đây không phải là mơ. Mọi giác quan đều rõ rệt, từ mùi bụi bặm xộc vào mũi đến tiếng thở hổn hển của chính tôi vang vọng trong không gian kín.

Và rồi, tôi nhớ ra một điều quan trọng hơn. Cô em gái của tôi.

``Saikyou!? Em đâu rồi!?'' tôi thốt lên. Trong lòng trỗi dậy một nỗi hoảng loạn mơ hồ.

``Khụ\dots\ khụ\dots\ Nii-chan\dots\ ở đây tối quá\dots''

Một giọng nói quen thuộc vang lên ngay bên cạnh, và tôi cảm thấy có thứ gì đó túm chặt lấy cánh tay trái của mình.

Tôi giật bắn người, quay phắt sang rồi thở phào nhẹ nhõm. Cô bé đang ôm chặt lấy tay tôi, ánh mắt rơm rớm nước mắt. Đôi vai run run. Nhưng vẫn là Saikyou.

Tạ ơn trời đất. Ít nhất tôi không cô đơn trong chốn kỳ quái này.

``Em ổn chứ? Có bị thương ở đâu không?''

Em gái tôi nhìn xuống bộ đồng phục phủ đầy bụi, rồi ngẩng đầu lên lắc lắc: ``Chắc là ổn\dots\ không thấy đau gì lắm\dots''

``Không gãy xương, không chảy máu chứ?''

``Không hề! Đây này!'' em ấy nói rồi dang rộng hai tay, xoay một vòng để tôi kiểm tra. Đúng là ngoại trừ vài vết trầy nhẹ và tóc hơi rối, trông Saikyou vẫn ổn.

Nhưng như thể vẫn chưa đủ để tôi tin, em đột ngột nói thêm: ``Nếu Nii-chan không tin thì\dots\ em có thể cho anh xem bên trong luôn cũng được\dots''

Tôi lập tức hơ hai tay, mặt nóng bừng lên. ``T-Thôi! Không cần! \textbf{Anh tin rồi!! Tin rồi mà!!!}''

\dots Phải rồi. Tôi quên mất rằng khi chỉ có hai anh em, Saikyou thường trở nên cực kỳ, ờm, thoải mái, hay nói trắng ra là dâm đãng. Không ít lần em ấy vô tư chạy vào phòng tôi trong khi đang thay đồ hay tắm rửa, rồi giả vờ ngây thơ như không có chuyện gì.

Cũng may là chỉ có tôi mới phải chịu đựng cái thói kỳ quặc này. Nếu có người lạ, Saikyou lập tức trở nên kín đáo và cẩn trọng đến lạ.

\dots Mà thôi, chẳng phải lúc để nghĩ mấy thứ đó.

Tôi hít một hơi thật sâu.

Tình hình bây giờ khá là\dots\ quái đản.

Chúng tôi đang ở đâu vậy? Sau cú rơi lúc nãy, tôi nhận ra mình và Saikyou không còn ở hành lang sập đổ kia nữa.

Không. Đây là một nơi hoàn toàn khác.

Một hành lang. Vẫn là hành lang của một ngôi trường\dots\ nhưng có gì đó sai sai.

\img{e0i}

Tường gỗ. Sàn gỗ. Trần gỗ. Những bóng đèn huỳnh quang mờ nhạt. Các lớp học dọc hai bên hành lang đều dùng cửa trượt bằng gỗ kiểu cũ, với ô lưới nhỏ. Cửa gỗ cũ kỹ, bảng hiệu phía trên đã tróc sơn, chỉ còn vài nét chữ mờ nhòe chẳng thể nào đọc nổi. Những ô kính nhỏ trên cửa đã vỡ tan tành từ đời nào, để lại khung lưới gỗ lởm chởm như miệng một con thú đang chực chờ nuốt trọn bất cứ kẻ nào bước vào.

Toàn bộ không gian đậm nét hoài cổ, mang phong cách kiến trúc học đường\dots\ khoảng thế kỷ 20?
Không giống chút nào với cấu trúc hiện đại mà chúng tôi đã thấy ở trường trung học Aogami, hay ít nhất là từ môi trường trước đó.

``Không lẽ chúng ta vừa bị\dots\ dịch chuyển?'' tôi lẩm bẩm.

Saikyou ngẩng lên, nhìn quanh một lượt rồi rùng mình nói nhỏ: ``Không phải là\dots\ giống như mấy trò trong C*rpse P*rty à\dots?''

Tôi cũng đang nghĩ vậy. Cái không khí lạnh lẽo, u tối, hoang vu, và\dots\ mùi cũ kỹ đến rợn người. Tôi dám cá\dots\ đây không còn là ``trường Aogami'' nữa.

Chúng tôi thử xoay người lại. Phía sau là một cánh cửa sắt màu đỏ, gỉ sét nặng, với tay cầm bị han gỉ đến mức rộp lên như da người bị bỏng. Đằng sau nó\dots\ hoàn toàn là bóng tối.

Chúng tôi thử mở nó ra.

``Nặng quá, Nii-chan\dots\ Nó kẹt rồi\dots''
Cả hai cùng dùng sức, nhưng nó không hề nhúc nhích.

Tôi thử ngẩng lên trần, hy vọng nhìn thấy cái ``lỗ'' đã đưa chúng tôi đến đây. Nhưng không có gì. Không lỗ thủng, không dấu vết. Trần gỗ liền mạch như chưa từng có gì xảy ra.

Tôi siết chặt tay lại. Chúng tôi thực sự đã bị nhốt ở đây rồi.

``Nơi quái nào thế này chứ\dots'', tôi lầm bầm. Cơn khó chịu bắt đầu xâm lấn đầu óc.

Mọi thứ xung quanh quá tối. Chúng tôi gần như chẳng nhìn thấy được gì nếu không nhờ ánh đèn lập lòe từ xa, mà ánh sáng ấy yếu đến mức chẳng khác nào châm chọc.

``À, phải rồi,'' Saikyou reo lên. ``Chúng ta có điện thoại mà! Sao mình không dùng đèn pin của nó nhỉ?''

Tôi giật mình: ``Ờ\dots\ đúng nhỉ. Sao anh không nghĩ ra nhỉ?''

Tôi vội thò tay vào túi quần. May quá, điện thoại vẫn còn nguyên. Cả hai chúng tôi bật đèn pin cùng lúc. Hai luồng sáng mỏng chiếu rọi vào hành lang u tối phía trước.

Chúng tôi bắt đầu lần bước dọc theo hành lang, ánh đèn điện thoại yếu ớt hắt ra chỉ vừa đủ xé một đường mỏng vào màn đêm đặc quánh đang nuốt trọn mọi thứ. Không gian nơi đây rộng lớn một cách bất thường, sâu hun hút đến mức tôi không thể đoán nổi nơi tận cùng. Mỗi bước chân lại kéo theo tiếng cót két vang vọng trên sàn gỗ đã mục, âm thanh ấy khô khốc và lạnh lẽo, như thể bên dưới lớp ván này ẩn giấu thứ gì đó đang âm thầm chờ chúng tôi sơ sẩy.

Thỉnh thoảng, một vài cánh cửa lớp học bên trái hé mở. Tôi liếc đèn pin vào trong, nhưng chỉ thấy bóng tối đặc quánh, vài tia sáng yếu ớt bị lớp bụi dày đặc phản chiếu lại, cùng những hàng bàn ghế đã mục, phủ rêu mốc, nằm bất động như một nghĩa trang gỗ.

Mọi thứ đều lặng im\dots\ quá lặng im.

Nhưng mới đi được chưa đầy chục bước, Saikyou khẽ kéo tay tôi, giọng nhỏ như sợ đánh thức thứ gì đó: ``Này, Nii-chan\dots''

Tôi cau mày: ``Sao nữa?''

Em ấy chỉ xuống sàn: ``Hình như ở dưới sàn\dots\ lại có mảnh giấy kìa.''

Tôi lia đèn xuống. Quả thật, ngay trước mũi chân là một tờ giấy ố vàng, chữ in đã phai nhưng vẫn đọc được. Thật lạ là nó dán chặt xuống nền gỗ, giống như đã ở đó từ rất lâu. Chúng tôi từng thấy một mảnh tương tự ở hành lang trước khi tỉnh lại\dots\ và tôi ghét cái cảm giác quen thuộc này.

Tôi cúi xuống, đọc thành tiếng: ``{\vnmsans Những hiện tượng kỳ lạ ở thị trấn Aogami: Ký sự số 1.}''

Tim tôi hơi chùng xuống. Lần này, nội dung không chỉ nhắc đến trường, mà mở rộng ra cả thị trấn.

\txtbx{\arial Thị trấn Aogami nằm trên một vùng núi, thuộc thành phố XX.}

Như dự đoán, phần tên thành phố đã bị cố ý gạch bỏ bằng mực đen. Có vẻ nơi này được ai đó muốn che giấu hoàn toàn, hoặc đã bị xóa khỏi bản đồ từ lâu.

\txtbx{\arial Trái với cái tên thanh khiết của nó, những hiện tượng kỳ lạ xảy ra ở đây vốn diễn ra thường xuyên.\\
  Người ta đồn rằng: nếu một học sinh chuyển đến mà cảm thấy không được chào đón, nỗi bất hạnh xung quanh đây sẽ ập tới.}

Phần cuối mới thực sự khiến da tôi lạnh buốt. Y hệt như nội dung mà gã nhà báo kia đã viết.

Saikyou ngước nhìn tôi, mắt mở to: ``Nghĩa là\dots\ nếu một ai đó tới trường Aogami, rồi bị xa lánh\dots\ thì chuyện khủng khiếp sẽ xảy ra sao, Nii-chan?''

Tôi gật chậm: ``Dường như là vậy. Nhưng anh thấy vẫn còn điều gì đó bất thường ở đây.''

Chúng tôi đứng im tại chỗ, xâu chuỗi những gì đã biết: học sinh mất tích, thảm họa bất ngờ, tin đồn về trường, việc trường đổ sập\dots\  tất cả dường như mắc vào cùng một mạch. Nhưng để tìm ``mẫu số chung'' chính xác thì\dots\ xin lỗi, tôi đâu phải \href{https://www.youtube.com/@cosu}{C* S*}!

Tuy vậy, tôi cũng đã bắt đầu hình thành một giả thuyết, cho đến khi Saikyou bất chợt nói, giọng trầm xuống: ``Có lẽ anh sẽ không tin\dots\ nhưng em cũng có một giả thuyết.''

Tôi im lặng, ra hiệu cho em ấy tiếp tục.

``Em nghĩ\dots\ đây là một lời nguyền.''

Câu đó khiến tôi hơi rùng mình, bởi bản thân tôi cũng đang nghi ngờ điều tương tự. Hình ảnh cô gái đứng lặng trong lớp học trước khi động đất xảy ra lại ùa về.

``Lời nguyền\dots\ từ bao giờ?''

``Em nghĩ nó bắt nguồn từ quá khứ\dots\ mười, hai mươi năm trước\dots\ hoặc xa hơn. Có thể đã từng có một học sinh chuyển tới đây. Lúc đầu được chào đón, nhưng rồi\dots''

Trong đầu tôi chợt hiện ra một cảnh mơ hồ: một nữ sinh đứng giữa sân trường, ban đầu được bạn bè vây quanh, nhưng dần dần từng người một rời đi\dots\ đến cuối cùng chỉ còn cô ấy đơn độc giữa khoảng sân trống.

``\dots Bị xa lánh hoàn toàn.'' Cả tôi và Saikyou vô thức đồng thanh.

Tôi nói tiếp: ``Và rồi\dots\ có thể cô ấy đã tự sát. Nhưng linh hồn không siêu thoát, quay lại ám trường\dots\ cho tới bây giờ.''

Nghe thì giống truyện ma cũ rích, nhưng giả sử đây là sự thật thì sao? Vấn đề, nạn nhân đó là ai?

Cả hai chúng tôi đồng thời nghĩ đến một khả năng: ``\dots Có thể là\dots\ cô gái trong lớp học lúc trước?''

Không loại trừ.

Nhưng tất cả chỉ là suy đoán. Để thoát khỏi nơi này, có lẽ chúng tôi sẽ phải tìm ra thêm những mảnh ghép, và tôi có cảm giác, những gì sắp tới sẽ còn tệ hơn rất nhiều.

Chúng tôi đặt mảnh giấy cũ kỹ lại đúng vị trí ban đầu, như thể sợ rằng chỉ cần động sai một chút là sẽ khiến cả nơi này nổi giận. Giống như những lần trước. Rồi lại tiếp tục đi, chiếc đèn pin trong tay tôi quét từng khoảng tối trước mặt. Ánh sáng vàng nhạt của nó hòa vào những tia sáng xanh mờ, yếu ớt chớp tắt từ các bóng huỳnh quang treo dọc hành lang.

Không hiểu vì lý do gì mà đám đèn huỳnh quang này vẫn còn hoạt động được sau từng ấy năm\dots\ hoặc có lẽ-- không, không, tôi phải tự nhắc mình rằng đây có thể là một không gian khác, một tầng thực tại chắp vá, không cần tuân theo những quy luật bình thường. Suy nghĩ quá nhiều về mấy chi tiết như vậy chỉ tổ khiến mình tự rối trí.

Chúng tôi bước qua vài cánh cửa đóng kín, thì bất ngờ bắt gặp một lớp học với cánh cửa mở toang. Nó thật sự mở, chứ không khóa cứng như những căn phòng khác.

``Có nên vào không, Nii-chan?'' Saikyou rọi ánh đèn vào bên trong, giọng vừa tò mò vừa cảnh giác.

Tôi nheo mắt nhìn vào khoảng tối kia. ``\dots Anh không thấy có gì quá bất thường\dots\ ít nhất là ở mức độ nhìn bằng mắt thường.'' Tôi cũng chẳng tin lắm vào lời mình, nhưng thôi, chân đã lỡ tiến lại gần, rút lui bây giờ thì kỳ.

Bên trong, không khí phảng phất mùi gỗ mục và bụi cũ. Bàn ghế gỗ xếp ngay ngắn nhưng phủ đầy tro bụi thời gian, bảng xanh vẫn treo ở đầu lớp, bên cạnh chiếc bàn giáo viên cũ kỹ. Trên trần, bốn bóng huỳnh quang chụm lại thành một cụm, chỉ duy nhất một chiếc còn lập lòe, ánh sáng đủ để khắc họa rõ những mảng tối đậm hơn ở mọi góc.

\img{e0j}

Trên bức tường đối diện, một chiếc đồng hồ kim đứng im lìm ở mốc 10:10. Điện thoại của tôi báo 10:45 tối. Có lẽ nó đã chết từ lâu\dots\ hoặc nếu nghĩ theo hướng tâm linh, đó có thể là thời điểm mà một chuyện kinh khủng nào đó đã xảy ra.

Cảm giác lạnh sống lưng thoáng qua, nhưng tôi vẫn cùng Saikyou rà soát khắp lớp. Từng dãy bàn, từng ngăn kéo, thậm chí cả những khoảng trống dưới gầm bàn giáo viên đều bị lục tung.

``Có vẻ chẳng có gì\dots'' tôi buông một tiếng thở dài sau vòng kiểm tra thứ ba.

``\dots Không hẳn,'' em gái tôi lên tiếng từ góc phòng, giọng khẽ nhưng có phần hứng khởi. Em cầm lên một vật hình tròn, phẳng, ánh sáng đèn pin hắt lên bề mặt phản chiếu mờ mờ.

Tôi lại gần, nhận ra đó là một mặt kính gắn trên khung thép không gỉ màu đen. Chính giữa có một hình vuông màu đen, kẹp giữa hai hình chữ nhật trắng. Trên viền kim loại, mấy dòng chữ khắc mờ nhưng tôi vẫn đọc được một cái tên: ``Gear Frontier.''

Cái tên đó làm tôi khựng lại. Tôi chắc mình đã từng nghe qua\dots\ phải rồi, hình như đó là tên một dòng đồng hồ đeo tay. Vậy thì\dots\ thứ này có thể khớp với phần dây đeo mà chúng tôi đã nhặt được trước đó.

Nếu đúng là vậy, thì chiếc đồng hồ này bị tách rời thành nhiều mảnh, rải khắp nơi. Nhưng vì mục đích gì? Và tại sao lại xuất hiện ở đây, trong một căn lớp học dường như bị bỏ quên hàng thập kỷ?

Trông nó có vẻ như cũng không đến từ thời đại này. Nhưng thôi kệ.

Tôi cất mặt kính vào túi, để chung với dây đeo, cảm giác như vừa nhặt được một mảnh ghép quan trọng, hoặc là một mảnh ghép nguy hiểm. Ngay khoảnh khắc tôi đứng thẳng dậy, bỗng có cảm giác như trong lớp học trở nên im ắng hơn hẳn, như thể thứ gì đó vừa nhận ra hành động của chúng tôi\dots

Em gái tôi tiếp tục: ``Em còn tìm thấy ba mảnh giấy trong đây nữa.''

Em xòe ra ba tờ giấy, hai tờ trông như được xé vội từ quyển vở cũ, giấy đã ngả vàng và hơi quăn mép, còn tờ còn lại lại là một mẩu báo, giấy mỏng và có mùi ẩm ẩm đặc trưng của thứ đã bị bỏ quên quá lâu trong môi trường ẩm thấp.

``Đọc thử cho anh cái phần mẩu báo lên xem nào.''

Em đưa hai mảnh giấy còn lại cho tôi và bắt đầu đọc phần tiêu đề trên mẩu báo. ``{\vnmsans Những hiện tượng kỳ lạ ở thị trấn Aogami: Ký sự số 2.} Em nghĩ đây là phần tiếp theo của cái bài trước đó, Nii-chan.''

Saikyou đọc to, giọng vang lên trong không gian trống trải:

\txtbx{\arial "Dựa trên những gì mà chúng tôi tiếp cận được, ở trong phòng Mỹ thuật của trường trung học Aogami, có một nữ sinh đã chết do bị mất máu quá nhiều.\\ Có học sinh cho rằng, cô gái này đã sử dụng máu của chính mình để vẽ một bức tranh nào đó mà cô cho là một 'kiệt tác'.\\
  Những người sống ở gần khu vực này đồn rằng vào ban đêm tại căn phòng này, nữ sinh đó vẫn đang lảng vảng ở đâu đó quanh đây."}

Tôi im lặng, cảm nhận một lớp sương mỏng vô hình vừa len vào tâm trí mình.

``Phòng Mỹ thuật à\dots'' tôi lẩm bẩm, mắt vẫn dán vào mẩu báo, như thể có thể nhìn xuyên qua những con chữ mà thấy được khung cảnh của căn phòng ấy.

``Nhưng tại sao cô ấy lại muốn vẽ một bức tranh bằng máu như vậy?'' Saikyou thắc mắc, nhíu mày.

Quả thật, bức tranh mà cô ấy gọi là kiệt tác là gì? Nó có giá trị đến mức phải trả bằng tính mạng sao? Và liệu nó có liên quan gì đến cô gái vừa chuyển trường mà chúng tôi thấy nhắc tới trong những tài liệu trước?

Tôi không tin đó là một người. Linh cảm bảo tôi rằng họ là hai cá thể hoàn toàn khác nhau, một người bị xa lánh nhưng sống, và một người đã chết vì mất máu. Nhưng nếu vậy\dots\ thì tại sao cả hai câu chuyện đều bị gắn với trường này, và đều vương mùi bất an như thế?

Dù gì thì, vẫn chưa có câu trả lời. Tôi vẫn còn hai mẩu giấy viết tay mà em tôi đưa cho. Cả hai đều như lời đồn được truyền miệng giữa học sinh, văn phong tuỳ tiện, nét chữ vội vã.

Tôi đưa cho em ấy một mẩu và nói: ``Em đọc cái của em, anh sẽ đọc cái của anh. Đọc xong rồi so sánh.''

Saikyou gật đầu.

Tôi mở mẩu giấy đầu tiên.
``Hừm\dots\ {\vnmsans 'Lời đồn trong trường học - Mẩu 1'}. Không rõ tác giả.''

Tôi bắt đầu đọc:

\txtbx{\arial "Này, biết gì chưa?\\
  Nghe nói là có hồn ma của một cô gái đã chết, cứ mỗi nửa đêm là cô ta lại xuất hiện trong phòng mỹ thuật, rồi nha, dùng máu của mình để vẽ tranh nữa cơ!"}

Tôi dừng lại, cảm giác rùng mình chạy dọc sống lưng. Lời đồn này trùng khớp đến khó tin với nội dung trong bài báo. Một câu chuyện được báo chí ghi lại, một câu chuyện truyền miệng trong học sinh, cả hai cùng chỉ về một điểm: căn phòng mỹ thuật đó.

Tôi quay sang thì thấy Saikyou vừa đọc xong mẩu thứ hai, nhưng khuôn mặt cô bé cau có thấy rõ, môi mím chặt như đang cố nuốt cơn tức giận.

``Sao vậy? Cái mẩu giấy đó làm em khó chịu à?'' tôi hỏi.

``Em thề, nếu em mà có đứa bạn như thế này thì em tát cho vài cái rồi đấy!'' em gằn giọng, rồi đưa mẩu giấy lên trước mặt tôi.

\txtbx{
  \vnmsans Lời đồn trong trường học - Mẩu 2:\\
  \arial "Cô ta đúng là lắm mồm thật, nếu muốn làm ồn thì làm ơn ra ngoài đi giùm. Mà, cô ta biến con mẹ nó đi thì càng tốt."}

Câu chữ lạnh lùng, ác ý đến mức tôi cũng phải chau mày. Không khó hiểu vì sao Saikyou nổi đóa, vì với em ấy, những kẻ thẳng thừng hắt hủi và hạ nhục người khác là loại người không đáng được nể nang.

Tôi đặt tay lên vai em gái, vỗ nhẹ. ``Thôi, anh hiểu cảm giác của em. Nhưng giờ mình cần bình tĩnh.''

Tôi biết rõ tính khí của cô em gái: sẵn sàng đứng ra bảo vệ người yếu thế, kể cả phải đối đầu với đám con trai ồn ào hay những cô nàng thích bắt nạt người khác. Có lẽ vì vậy mà em ấy ít bạn, dù thật lòng lại khao khát có những người bạn đúng nghĩa.

Cô gái chuyển trường năm đó\dots\ nếu thật sự là nạn nhân trong mẩu giấy này, chắc hẳn sẽ biết ơn Saikyou lắm.

Sau khi để em ấy hít sâu vài lần, tôi mới khẽ gật: ``Đi thôi. Chúng ta vẫn còn đường phải khám phá.'' Em gái tôi cất mấy mẩu giấy vào túi, vẫn còn hậm hực nhưng đã bước về phía cửa.

Tôi vừa định đi theo thì ánh mắt vô thức dừng lại ở một lọ hoa trắng đặt trên bàn giáo viên.
Nó đơn giản, không có gì nổi bật, nhưng\dots\ có cái gì đó trong nó khiến tôi không thể rời mắt.

Một cảm giác mơ hồ trào lên, như thể lọ hoa này không chỉ để trang trí.

Một biểu tượng? Một dấu hiệu? Hay là một phần của câu chuyện mà chúng tôi đang bị cuốn vào?

``\textbf{Nii-chan! Ta phải đi thôi!}'' Giọng Saikyou vang lên từ cửa, kéo tôi ra khỏi mạch suy nghĩ.

Tôi hít một hơi, quay lưng lại với cái lọ hoa và bước ra hành lang. Nhưng cảm giác bất an ấy\dots\ vẫn bám chặt như cái bóng không chịu rời.

\ct

Chúng tôi tiếp tục lê bước dọc theo hành lang, chỉ nghe tiếng dép chạm nền gạch lạnh lẽo vọng dài trong khoảng tối. Những ô cửa sổ cũ kỹ, dính đầy bụi, thỉnh thoảng hắt vào một chút ánh sáng mờ đục, nhưng không đủ xua bớt cái cảm giác tù túng. Dọc đường, lại thêm những lớp học im ỉm, cửa trượt hoặc là bị khóa chặt đến mức cả hai anh em dùng sức cũng chẳng nhúc nhích, hoặc mở ra thì chỉ thấy bốn bức tường trống rỗng, lạnh ngắt, chẳng một mẩu giấy hay dấu vết nào để lần theo.

Nghĩ kỹ lại, tôi bắt đầu tự hỏi\dots\ rốt cuộc mục đích ban đầu của chúng tôi ở đây là gì nữa? Chẳng phải ngay từ đầu là để tìm đường thoát hay sao? Thế mà giờ lại giống như đang chơi trò săn kho báu trong một mê cung bỏ hoang, mà ``kho báu'' ở đây thì toàn là tin đồn, giấy vụn và những câu chuyện chẳng rõ thật giả.

Đang miên man suy nghĩ, tôi giật mình khi nghe Saikyou gọi khẽ: ``Nii-chan! Cái lớp này\dots\ hình như đang mở này!''

Tôi nhìn theo tay em chỉ, hừ nhẹ: ``Dám chắc là lại chẳng tìm được gì hữu ích đâu.''

Em ấy hơi xị mặt, nhưng rồi lại tự vực tinh thần lên: ``Cũng phải\dots\ Nhưng mà, còn nước còn tát!''

Tôi cười gượng, buông một câu đùa nhạt nhẽo: ``Cảm giác như đang chơi game tìm đồ vật chứ không phải tìm lối thoát nữa nhỉ.''

``Hẳn là thế!'' Em tôi cười đáp, nhưng ánh mắt vẫn ánh lên tia hy vọng nhỏ nhoi.

Chúng tôi quyết định bỏ qua căn phòng đó, tiếp tục đi tới. Trước mặt là một cánh cửa trượt khác. Tôi thử kéo mạnh\dots

*LẠCH CẠCH\dots\ LẠCH CẠCH*

Cánh cửa chỉ rung nhẹ, không hề dịch chuyển. ``Kẹt rồi.''

Saikyou tiến lên, thử thay tôi, nhưng kết quả chẳng khác. Cánh cửa vẫn không chịu nhúc nhích.

Tôi thở dài. ``Đừng nói là lại phải quay lại phòng ban nãy nha?''

Em tôi chống hông, liếc tôi kiểu như đã đoán định trước: ``Thấy chưa! Nào, vào thử xem trong đó có gì!''

Nói là nói vậy, nhưng tôi cũng biết nếu bỏ qua thì kiểu gì sau này cũng hối tiếc. Thế là cả hai quay lại, bước vào căn phòng đang sáng đèn ấy.

Ngay khi vừa đặt chân vào, đập vào mắt chúng tôi là một quyển vở màu trắng, chỗ gáy đã sờn, mặt bìa lấm tấm những vết ố vàng như bị nước thấm lâu ngày. Ô tên ở góc bìa mờ đến mức chẳng thể đọc nổi cả chữ lẫn furigana.

Lật trang đầu tiên, chữ viết gọn gàng, nắn nót. Chắc chắn là của một nữ sinh. Dòng đầu ghi:

\txtbx{\arial "Mới ngày đầu chuyển tới nên mình khá là hồi hộp, nhưng may thay mình đã kết thân được bạn. Mình muốn họ sẽ mãi là bạn bè. Mãi mãi là thế."}

Đọc xong, tôi lặng im. Đáng ra phải mừng cho cô gái ấy mới đúng\dots\ nhưng một cảm giác tiếc nuối len vào, lẫn chút lo lắng mơ hồ. Tại sao chứ? Cảm giác ấy cứ như báo trước điều gì đó chẳng lành.

``Không hiểu sao\dots\ em cũng thấy khó tả'' Saikyou khẽ nói, ánh mắt vẫn dán vào trang giấy, ``vừa muốn vui cho cô ấy, vừa\dots\ không yên tâm.''

Tôi gật nhẹ. Cả hai đặt quyển vở xuống bàn, ánh mắt tôi lại dừng ở bàn trên cùng, nơi có một mẩu báo khác. Saikyou nhanh chân hơn, cầm lên đọc tiêu đề: ``{\vnmsans Tin nóng: Hiện tượng kỳ lạ ở thị trấn Aogami - tin thứ ba.}''

Vẫn cùng cái kiểu ``nguồn'' lá cải như hai mẩu trước, và vẫn có phần ngày tháng bị gạch xóa. Tôi bắt đầu thấy mấy tờ báo này chẳng đáng tin chút nào\dots\ nhưng thôi, đọc cho biết.

\txtbx{"Tại khu rừng ở gần thị trấn Aogami, dù trời vẫn là ban ngày nhưng xung quanh vẫn rất tối, bầu không khí đặc quánh, rợn ngợp. Người ta đồn rằng khi đi sâu vào trong, họ nhìn thấy nhiều đứa trẻ lảng vảng giữa những thân cây. Không ai có thể tiếp cận được chúng, bởi lẽ, họ tin rằng những đứa trẻ ấy không thuộc về thế giới này."}

Tôi khẽ rùng mình. Thêm một mảnh ghép mới, nhưng càng nhiều thông tin, bức tranh toàn cảnh càng rối rắm. Không biết thứ gì trong đây là thật, và thứ gì chỉ là hư cấu để câu chuyện thêm ám ảnh. Nhưng\dots\ dẫu chỉ là tin đồn, nó vẫn để lại cảm giác lạnh chạy dọc sống lưng.

Chúng tôi định tiến đến cánh cửa để thoát ra khỏi phòng, thì bất chợt\dots

*XOẸT\dots\ XOẸT XOẸT\dots*

Âm thanh như tiếng phấn cọ mạnh vào bảng vang lên đột ngột, sắc và khô khốc, khiến cả hai chúng tôi giật bắn người. Tôi theo phản xạ lia ngay đèn pin điện thoại về phía đó.

Quả nhiên, trên bảng đen đã bạc màu kia, từng nét chữ trắng đang hiện ra, rõ ràng là đang được viết, nhưng chẳng hề có lấy một bàn tay hay cây phấn nào cả.

Chữ được viết nắn nót, chậm rãi, như thể có ai cố tình để chúng tôi nhìn thấy:

\begin{center}
  ``Những kẻ không làm hài lòng những học sinh chuyển đến\dots''
\end{center}

\img{e0k}

Tôi cảm thấy một luồng lạnh chạy dọc sống lưng. Không hiểu hoàn toàn ý nghĩa, nhưng trong đầu tôi lập tức hiện ra hình ảnh những kẻ đã hắt hủi ai đó ngay từ ngày đầu họ bước vào đây. Và khi nghĩ đến mấy mẩu ghi chú rải rác mà chúng tôi tìm được trước đó\dots

``\dots chẳng phải nghe rất khớp với những gì chúng ta biết rồi sao?'' Saikyou bỗng nói, gần như đọc được ý nghĩ của tôi.

Gió lạnh (hoặc ít nhất là cái cảm giác như gió lạnh) lướt qua gáy tôi. Rõ ràng, có một thứ gì đó đang thao túng nơi này, và nó không phải thứ mà con người bình thường có thể giải thích.

Cảm giác như thời điểm này, ngôi trường đang mắc kẹt giữa quá khứ và hiện tại\dots\ ngay tại nút thắt của một chuỗi sự kiện bi thảm.

\dots

``*cười mỉm* Nghe nói\dots\ học sinh mới chuyển trường\dots\ nhỉ\dots? *cười mỉm*''

Một giọng nói vang lên từ phía sau bên phải, nhẹ nhưng rợn buốt, khiến da tôi nổi gai. Tôi lập tức quay phắt sang trái theo bản năng.

``Là em sao!?''

``Không! Em ở đây nãy giờ mà!'' em tôi đáp, giọng pha chút hốt hoảng.

Nếu không phải em ấy thì\dots

\dots

\img{e0l}

Cô gái đó.

Cô gái với mái tóc nâu dài chạm lưng, ánh mắt sâu thẳm khó đoán, từ lớp học mà chúng tôi từng thấy ở không gian trước. Vẫn bộ đồng phục đen, khăn quàng cổ xanh, và chiếc nơ cài lệch một bên như cũ.

Cô không nhìn chúng tôi, chỉ lặng lẽ bước tới cánh cửa trượt.

*XOẠCH!*

Cửa mở ra trơn tru như chưa từng bị kẹt. Cô đi qua, bóng lưng khuất dần trong ánh sáng mờ ngoài hành lang.

``Ơ, này, khoan đã, cô gì ơi!'' tôi vội chạy theo, Saikyou cũng gọi với: ``Nii-chan!?''

Nhưng vừa bước ra khỏi lớp\dots\ trống không. Cô ấy đã biến mất như thể chưa từng tồn tại. Không tiếng bước chân, không tiếng cửa, thậm chí không có cả hơi gió.

``\dots Cô ấy có thể đi đâu được nhỉ?'' tôi lẩm bẩm.

\dots

``N-Nii-chan?''

\dots

Một ý nghĩ nặng nề dần kết tinh trong đầu tôi: cô ấy không chỉ đơn thuần là một hồn ma vất vưởng. Cô gái kia là trung tâm của tất cả những gì đang diễn ra ở đây. Có thể là khởi nguồn cho cả ngôi trường\dots\ và thậm chí là cả thị trấn này.

``\textbf{Nii-chan!}''

Giọng Saikyou bỗng vang to, đầy bực bội, kéo tôi ra khỏi cơn mê man suy tưởng.

Tôi quay lại, quét ánh đèn pin vào mặt em ấy. ``\textbf{Giồi ôi! Con ma!}''

``\textbf{Ma cái gì mà ma! Em đây!}'' em ấy chống nạnh. ``Nãy giờ em gọi mà anh chẳng nghe gì! Tập trung vào đi cái!''

``Xin lỗi\dots\ từ nãy giờ anh vẫn đang nghĩ về cô ấy\dots'' tôi gãi đầu, cố tìm cách cắt mạch suy nghĩ.

``Hể\~{}? Anh bị cô ta cuốn hút rồi chứ gì?'' cô em gái lập tức nheo mắt, giọng chọc ghẹo.

Tôi chỉ thở dài, lắc đầu: ``Không có gì đâu.''

Siết chặt điện thoại trong tay, tôi nói: ``Thôi, ta tiếp tục tìm lối ra.'' Hai anh em lại bước vào hành lang tĩnh lặng, để mặc sau lưng căn phòng vẫn còn mùi phấn vương vất.

Chúng tôi tiếp tục bước đi trên hành lang vắng tanh, chỉ còn lại tiếng bước chân vọng dài, lẫn trong thứ mùi ẩm mốc đặc trưng của một ngôi trường đã bị bỏ hoang từ lâu. Mỗi lần ánh đèn pin lia qua, những chiếc cửa sổ kính bụi mờ lại phản chiếu lại ánh sáng, giống như có ai đó đang nhìn ra từ bên trong. Nhưng chẳng có ai cả, không một tiếng người, không một hơi thở, không một dấu hiệu của sự sống, ngoại trừ chúng tôi.

Lớp học nào đi qua cũng chẳng khác gì nhau: phòng thì có cửa bị khóa cứng, phòng thì bên trong trống rỗng đến lạnh lẽo. Mỗi khi thử mở, tiếng bản lề rít lên nghe như ai đó đang nghiến răng, khiến cả hai anh em đều rùng mình.

Cho đến khi, ở cuối dãy hành lang, chúng tôi bắt gặp một chướng ngại bất ngờ.

Trước mặt là một đống bàn ghế học sinh bị chất chồng lên nhau, chắn ngang gần hết lối đi. Chúng được xếp lộn xộn, như thể ai đó đã vội vàng dựng nên một bức tường tạm bợ để ngăn đường.

\img{e0m}

``\dots Nghiêm túc đấy à?'' tôi lầm bầm.

``Có vẻ như chúng ta bị chặn rồi,'' Saikyou vừa nói vừa thử với tay kéo một cái ghế ra. Nhưng ngay lập tức em ấy rụt tay lại, vẻ mặt ngạc nhiên: ``Sao\dots\ em không thể chạm vào nó được nhỉ?''

Tôi cũng thử. Bàn tay tôi xuyên qua khung gỗ như chạm vào một làn sương đặc quánh, lạnh ngắt đến buốt da. ``Anh cũng thấy lạ đấy\dots\ Không lẽ chúng bị ám bởi một sức mạnh nào đó?''

Dù tôi cố gắng dùng hết sức, cũng chẳng thể nhấc hay xô đẩy bất kỳ thứ gì, thậm chí tương tác với chúng cũng là điều bất khả thi. Cảm giác đó chẳng khác nào cố đẩy một bức tường bằng bóng tối.

Nhưng rồi ánh đèn pin của chúng tôi dừng lại ở một thứ khác, một cuốn tạp chí nằm gọn trên đỉnh đống bàn ghế, như cố tình được đặt ở đó để thu hút sự chú ý.

Bìa của nó màu đen tuyền, in hình một con ma nữ với mái tóc rũ rượi, đôi mắt trắng dã, đang bám chặt vào song sắt. Bên dưới là tựa đề in đậm bằng cả chữ Nhật và phiên âm: ``\ruby{超}{ちょう}\ruby{常}{じょう}\ruby{現}{げん}\ruby{象}{しょう}\ruby{最}{さい}\ruby{恐}{きょう}\ruby{心}{しん}\ruby{霊}{れい}'' - Những hiện tượng huyền bí tâm linh đáng sợ nhất.

Tôi cẩn thận cầm nó xuống. Giấy đã ố vàng, bìa sờn rách, và mỗi khi lật trang, từng làn bụi mịn lại bay lên, khiến ánh đèn pin lấp lánh như có hàng ngàn hạt bụi đang lơ lửng trong không khí.

Hành lang quá tối, nên chúng tôi cúi sát lại, dùng đèn pin điện thoại soi thẳng vào từng trang. Đa phần là những bài báo, ảnh chụp mờ nhòe, và những lời đồn được ghi lại từ đâu đó.

Cho đến khi tôi dừng lại ở một trang đặc biệt.

Trên đó là hình một lớp học tối om, những chiếc bàn và ghế bay lơ lửng giữa không trung, như bị một bàn tay vô hình nhấc lên. Phía dưới ảnh là đoạn chữ in:

\txtbx{\arial "Bảy điều bí ẩn trong trường có lẽ đã không còn quá xa lạ với bạn đọc. Nhưng trên thế gian này vẫn còn rất nhiều hiện tượng dị thường khác chưa có lời giải đáp. Một trong số đó là 'Lớp học hồi chuyển.'\\
  Chuyện kể rằng vào lúc nửa đêm, khi oán khí dâng cao nhất, bàn ghế trong lớp sẽ tự động di chuyển mà không có một lực tác động nào. Nhưng ngoài ra, còn có lời đồn rằng nếu bạn lọt vào một lớp học như vậy, bạn sẽ mãi mãi bị mắc kẹt, cho đến khi bạn làm một việc gì đó khác thường."}

Saikyou đọc xong, khẽ rùng mình: ``\dots Nghe đúng là đáng sợ thật đấy, Nii-chan.''

Tôi cũng không khá hơn. Ý nghĩ bị giam lỏng mãi mãi trong một căn phòng chỉ vì một trò đùa của thế lực nào đó khiến gáy tôi lạnh toát.

Chúng tôi đặt cuốn tạp chí xuống và nhìn về phía phòng học ngay bên cạnh. Cửa mở hé, bên trong sáng hơn hẳn so với những phòng trước, nhưng lại vắng lặng kỳ lạ. Chỉ có lác đác vài chiếc bàn, và tuyệt nhiên không có cái ghế nào.

Có lẽ đống bàn ghế chặn lối ngoài kia đã được lấy từ đây. Nhưng ai là người làm chuyện này? Con người, hay thứ gì khác?

Trước khi tôi kịp nghĩ tiếp, Saikyou đã gọi: ``À, Nii-chan, em tìm thấy cái mẩu báo nữa này.''

Em ấy gỡ một tờ giấy đã ngả vàng từ bảng xanh, đưa cho tôi. Trên đó ghi: ``{\vnmsans Những hiện tượng kỳ lạ ở thị trấn Aogami: Ký sự số 4}.''

Tôi cau mày. ``Kỳ lạ nhỉ\dots\ tờ số 3 đâu?''

``Em còn mang từ lớp trước này,'' em chìa ra một mẩu báo khác: ``{\vnmsans Những hiện tượng kỳ lạ ở thị trấn Aogami: Ký sự số 3}''. Không biết từ lúc nào em ấy đã cất nó trong túi.

Tôi không giấu được chút ám ảnh nhẹ với việc đọc đủ mọi phần của một tài liệu. Gọi là OCD cũng đúng. Nếu đã có số 3 và 4, tôi phải đọc cả hai, dù nội dung có thể chẳng giúp ích gì mấy.

Tờ ``{\vnmsans Ký sự số 4}'' ghi:

\txtbx{"Phòng Mỹ thuật bị nguyền rủa, hay chính cô gái bị mắc lời nguyền? Thật khó để có câu trả lời.\\
  Nếu là trường hợp đầu tiên, rất có thể nó liên quan đến toàn bộ thị trấn này, và ta sẽ cần biết nhiều hơn về lịch sử nơi đây. Trở ngại lớn nhất: gần như không có tài liệu ghi chép nào về thị trấn Aogami.\\
  Còn nếu rơi vào trường hợp thứ hai--"}

Đoạn cuối bị xé mất, để lại mép giấy răng cưa.

Tờ `{\vnmsans Ký sự số 3}'' thì khác:

\txtbx{"Có một điều mà chúng tôi vẫn còn khá băn khoăn.
  Liệu lời đồn về cô gái vẽ bức tranh mà cô ấy cho là một `kiệt tác' ấy có thật hay không?\\
  Thật sự cô ấy có bị mất trí không khi đã dùng máu của chính mình để vẽ một bức tranh?\\
  Có lẽ chúng tôi nên đào sâu hơn để hiểu rõ ngọn ngành của những hiện tượng dị thường này."}

Đọc xong, tôi khẽ thở ra. Cả hai mẩu tin đều cho thấy ngay cả những người điều tra trước đó cũng không chắc chắn, thậm chí còn nghi ngờ chính lời đồn mình ghi lại.

Nhưng nếu học sinh ở đây thực sự đã từng tận mắt chứng kiến\dots\ thì chắc chắn có điều gì đó họ chưa kể ra.

Và cảm giác ấy - rằng chúng tôi chỉ đang nhìn thấy bề nổi của một tảng băng chìm - khiến tôi không thể dứt suy nghĩ về cô gái tóc nâu bí ẩn kia.

Tôi sờ vào túi quần. Ngay lập tức, các ngón tay chạm phải một thứ gì đó ráp ráp, như thể là giấy đã nhàu nát từ lâu. Một cảm giác bất an nổi lên trong lòng, khiến tôi lập tức lôi nó ra, sợ rằng đó là một thứ gì nguy hiểm.

Một mẩu giấy bị vò lại thành cục tròn, mép đã xơ xác như đã nằm ở đây từ nhiều năm. Thật kỳ lạ\dots\ nó đã ở trong túi tôi từ bao giờ? Tôi đâu nhớ mình từng nhặt nó.

Khi mở ra, tôi nhận ra đây là một mẩu ghi chú khác của một học sinh tại ngôi trường này. Ở trên đầu có tiêu đề run rẩy bằng bút mực: ``{\vnmsans Lời đồn trong trường học - Mẩu 3}.''

Chữ viết xiêu vẹo, nhỏ và ngập ngừng như của một người quá nhút nhát hoặc đang sợ hãi khi viết. Đại ý như sau:

\txtbx{\arial Quả thực\dots\ lời nguyền ở phòng Mỹ thuật đúng là đáng sợ thật\dots\\

  N-Nếu mình là người bị mắc kẹt ở đó vào giữa ban đêm\dots\ r-rất có thể mình đã bị\dots\ sử dụng như một công cụ để vẽ rồi\dots}

Chỉ đọc thôi mà tôi đã rùng mình. Nếu ghép mẩu này với hai bài tin trước, không khó hiểu vì sao người viết lại sợ đến vậy. Nếu là tôi, chắc tôi cũng chẳng khác gì, có khi còn run như cầy sấy.

Chúng tôi cẩn thận ghim hai mẩu tin và cả tờ giấy nhàu đó lên bảng xanh, xếp ngay ngắn như một phần hồ sơ điều tra. Sau đó, tiếp tục rà soát khắp căn phòng. Không gian ở đây sáng hơn, rộng hơn, và ít đồ đạc hơn so với những lớp khác, nhưng điều đó chẳng khiến việc tìm manh mối dễ hơn là bao.

Chúng tôi chỉ tìm thấy một cuốn sách giáo khoa mới tinh, giấy còn thơm mùi mực in. Có vẻ đây là bản dùng nội bộ trong trường, thậm chí có thể chưa từng được phát cho học sinh.

Tôi lật bìa trong. Dù ánh sáng đủ để đọc, chữ vẫn nhòe một chút, có thể vì mực phai hoặc do người viết dùng bút mực rẻ tiền.

Dòng tên ghi: Misora Shino (\ruby{美}{み}\ruby{空}{そら}\ruby{時}{し}\ruby{乃}{の}, \ruby{Mi}{Mỹ}-\ruby{xô-ra}{Không}\ \ruby{Si}{Thời}-\ruby{nô}{Nãi}).

Tôi khựng lại một chút. Cái tên này\dots\ tôi chưa từng nghe bao giờ. Có thể là một trong những học sinh đã từng học ở đây. Hoặc, tệ hơn, là một trong số những người\dots\ không bao giờ rời khỏi ngôi trường này nữa. Tôi chẳng cần nói rõ hơn, bạn cũng hiểu ý tôi.

``Anh tìm được gì à, Nii-chan?'' Saikyou hỏi khi thấy tôi chăm chú.

``Ừ. Có vẻ là đồ bỏ lại\dots\ của ai đó.''

Em gái tôi liếc qua, rồi chưng hửng: ``Tên là\dots\ Misora Shino-san? Em chưa từng nghe.''

Tôi lắc đầu: ``Anh cũng thế.''

Khi tôi vừa đặt cuốn sách xuống, Saikyou lại gọi: ``Mà, em còn thấy một thứ\dots\ còn khó hiểu hơn nữa trên bàn giáo viên.''

Tôi bước lại gần. Trên bàn là một cuốn tạp chí y hệt cái chúng tôi đã đọc trước đó, với bìa đen, hình con ma nữ sau song sắt. Nhưng bên cạnh nó lại có một bức ảnh in mờ, ghi lại một lớp học với ba hàng bàn, mỗi hàng năm chiếc, trông rất chỉnh tề. Có thể là ảnh chụp lúc mới bắt đầu năm học, trước khi nơi này trở thành\dots\ như bây giờ.

``Tại sao nó lại ở đây nhỉ\dots'' tôi lẩm bẩm, đưa tay chạm vào bìa tạp chí.

Ngay lập tức, ánh sáng trong phòng chớp tắt. Trong một khoảnh khắc ngắn ngủi, mọi thứ tối đen như mực, rồi đèn bật sáng trở lại.

``\dots Chuyện quái gì vừa xảy ra vậy?'' tôi quay sang hỏi, nhưng cô em gái của tôi lại đang há hốc miệng.

``Lớp học\dots\ vừa thay đổi rồi\dots'' em ấy chỉ về phía dãy bàn bên trái. Nơi trước đây chỉ có một chiếc bàn, giờ đã xuất hiện thêm ba chiếc nữa, xếp ngay ngắn như vừa được ai đó mang tới.

``Nhưng\dots\ nó vẫn chưa đúng với ảnh,'' tôi vô thức lẩm bẩm. Căn bệnh ``ưa hoàn hảo'' của tôi lại trỗi dậy, dù đây rõ ràng không phải lúc.

``Vậy để em làm nốt phần còn lại nhé?'' em nói, rồi cũng đặt tay lên tạp chí. Một lần nữa, cả phòng tối sầm, rồi sáng bừng. Ngoài hành lang vang lên một tiếng *ROẠCH* khô khốc.

Chúng tôi nhìn quanh, và đúng như dự đoán, giờ đây cả lớp đã đủ mười lăm chiếc bàn, y hệt bức ảnh. Tiếng động ngoài kia khiến hai anh em vội chạy ra, và thấy dãy bàn ghế chặn đường đã biến mất, mở ra một hành lang dài thẳng tắp.

``\dots Anh biết không, Nii-chan, cảm giác này cứ như chúng ta vừa giải xong một màn trong trò chơi kinh dị ấy,'' Saikyou lau mồ hôi.

``Anh mà biết nó thật sự hoạt động như vậy thì đã không đùa ngay từ đầu rồi\dots'' tôi đáp, cố nén tiếng thở mạnh.

``Thôi thì\dots\ tiếp tục chứ?''

``\dots Ừ. Đi thôi.''

Chúng tôi rời căn phòng kỳ quái ấy, bước tiếp vào hành lang tối mờ, nơi bóng đèn trên trần kêu rè rè, như báo hiệu rằng những chuyện kỳ lạ vẫn chưa kết thúc.

\pov{Haragen Saikyou}

Trên đường đi, chúng tôi lướt qua thêm nhiều lớp học khác nữa. Cũng giống như những cái trước, cửa đều đóng then cài, bên trong tối om, im lìm như chưa từng có ai đặt chân vào từ rất lâu. Giờ thì tôi chẳng còn buồn ngạc nhiên nữa, chuyện này lặp đi lặp lại đến mức thành thói quen rồi.

Ánh sáng đèn huỳnh quang trên trần bắt đầu xuất hiện dày hơn, như thể ai đó đã bật thêm một nửa số đèn còn lại. Thế là chúng tôi quyết định dựa vào thứ ánh sáng trắng lạnh lẽo ấy để tiết kiệm pin điện thoại. Dù sao thì, nói là sáng hơn chứ thực ra chúng cũng chẳng đủ sức xua bớt cái cảm giác tối tăm, ngột ngạt của hành lang dài bất tận này.

Nhưng thôi, có còn hơn không.

Ít nhất thì\dots\ tôi vẫn có Nii-chan bên cạnh. Hy vọng là vậy.

Tôi chẳng muốn ở lại đây thêm một giây nào nữa. Cái cảm giác ẩm lạnh của gió lùa qua các khe cửa, tiếng rè rè của đèn huỳnh quang, và cả những khoảng tối sâu hút như nuốt trọn mọi thứ, tất cả khiến tôi chỉ muốn chạy một mạch ra lối thoát.

\dots

\dots

Đột nhiên, một mảnh giấy vở được dán chéo trên cây cột cạnh một phòng học bị khóa khiến tôi chú ý. Giấy đã ố vàng một chút, mép quăn lại như bị hơi ẩm gặm nhấm.

Tôi kéo nhẹ tay áo Nii-chan, hất cằm về phía đó: ``Nii-chan, hình như lại có một mẩu ghi chú nữa kìa.''

Chúng tôi tiến lại gần. Ánh sáng ở đây đủ để đọc rõ dòng tiêu đề viết bằng mực xanh đã nhạt màu: ``{\vnmsans Lời đồn trong trường học - Mẩu 4.}''

Nội dung nguệch ngoạc, chữ hơi run như của người viết trong lúc vội:

\txtbx{\arial Người đó nha, lại đi cãi nhau với học sinh với mấy đứa ở trường khác rồi. Thật đúng không biết xấu hổ.\\
  Mình nhớ người đó trước đây đâu có như vậy\dots\ Có đúng thật không nhỉ?}

Tôi im lặng một lát. Có lẽ người mà tờ giấy nói đến đã thay đổi nhiều lắm. Kiểu thay đổi này thường không phải ngẫu nhiên: hoặc gia đình có vấn đề, hoặc họ đã trải qua những cú va chạm rất tệ ở trường.

Nhưng nếu để tôi đoán, chắc chắn là do trường học rồi. Không chừng cậu ta bị bắt nạt đến mức phải trở nên dữ tợn để chống lại.

Nii-chan đột nhiên lên tiếng: ``Anh nghĩ\dots\ em cũng giống vậy. Ngày xưa em đâu có hay cãi nhau thế này.''

Tôi nhăn mặt, phồng má: ``Cũng tại cái đám kia toàn chọc ghẹo anh em mình chứ bộ!'' Nếu không phải anh là anh trai em thì em đã ``xử'' anh rồi đấy, Nii-chan ngốc ạ!''

Thấy anh vẫn đứng trầm ngâm nhìn mẩu giấy, tôi thở dài, bước lên trước. Việc gì phải bận tâm nhiều đến thế chứ\dots

Nhưng mới đi được vài bước, tôi chợt khựng lại.

Bên trong phòng học ngay phía trước, qua lớp cửa kính mờ mịt bụi, có hai nữ sinh đang đối diện nhau.

\img{e0n}

Một người có mái tóc xanh nước biển, buộc đuôi ngựa lệch sang một bên, ngồi thẳng trên bàn ở dãy trong cùng. Cô ta không cử động, ánh mắt nghiêm nghị như đang chờ một điều gì.

Người còn lại có tóc trắng pha xám, tết thành hai bím, đứng đối diện. Khuôn mặt cũng căng thẳng chẳng kém.

Họ đứng đó im lìm, tĩnh như tượng. Giống hệt những ``người'' mà chúng tôi từng chạm trán trước đây\dots\  không hẳn sống, cũng chẳng hẳn là ảo ảnh.

*XOẸT*

Chớp mắt, cả hai biến mất, để lại căn phòng trống rỗng. Tim tôi hẫng một nhịp, miệng buột bật ra tiếng kêu khẽ. ``Eek!''

``Có chuyện gì sao?''

``\textbf{Kyaaah!}''

Tôi giật mình hét toáng, quay ngoắt lại, chỉ để thấy Nii-chan đang trợn mắt nhìn tôi.
``Trời đất\dots\ là anh à!'' tôi ôm ngực thở hổn hển. ``Tưởng\dots\ tưởng cái gì khác cơ chứ!''

``Thế rốt cuộc em vừa thấy gì mà ngẩn ra thế?''

Tôi vội xua tay, lắc đầu: ``Không\dots\ Không có gì, chắc là em nhìn nhầm thôi\dots''

Tôi không muốn kể về hai nữ sinh đó. Chỉ cần nghĩ đến chuyện giải thích cho anh là đã thấy rắc rối rồi, và tôi cũng đoán được họ chỉ là một phần khác của cái bí ẩn nơi này.

``T-Ta đi tiếp thôi nhỉ?'' tôi nói lảng.
Nii-chan nghiêng đầu, nhưng rồi cũng gật.

Trong lòng tôi chợt dấy lên một linh cảm chẳng lành: không biết chúng tôi còn phải chạm trán bao nhiêu kẻ như thế này nữa\dots

\ct

Chúng tôi bước chậm từng bước một, như đang đi trên một cây cầu tre ọp ẹp sắp gãy. Mỗi lần bàn chân tôi đặt xuống, tiếng kẽo kẹt lại dưới những thanh sàn gỗ vang lên rõ mồn một, nghe chẳng khác gì một lời cảnh báo ``đừng bước nữa nếu không muốn xuống thẳng âm phủ''.

Ánh sáng từ đèn pin điện thoại chỉ đủ chiếu ra một khoảng mờ mịt phía trước. Những đốm bụi lơ lửng bay qua ánh sáng, xoáy thành từng vòng nhỏ như những con côn trùng li ti. Không khí thì đặc quánh mùi gỗ mục và thứ gì đó ngai ngái, như mùi ẩm mốc trộn với kim loại gỉ.

``Này\dots\ Nii-chan\dots'' tôi lên tiếng khe khẽ.
``Gì?''

``Anh nghĩ\dots\ trần nhà chỗ này sẽ trụ được bao lâu?''
``Im lặng đi, đừng hỏi mấy câu kiểu ấy lúc này,'' anh ấy trả lời, nhưng tôi thấy rõ bàn tay anh siết chặt hơn quanh tay mình. Có vẻ như câu hỏi của tôi cũng chạm đúng nỗi lo của anh ấy.

Tiếng gió rít đâu đó vang lên. Không rõ nó luồn qua khe cửa hay qua những lỗ hổng vừa mới xuất hiện trên trần, chỉ biết nó tạo ra một âm thanh rền rĩ khó chịu, cứ như có ai đó đang thì thầm tên chúng tôi trong bóng tối.

Tôi lén liếc nhìn vào những khung cửa sổ bị bụi phủ mờ. Ở ngoài kia tối đen như mực, chẳng thấy trăng sao, chẳng thấy gì ngoài bóng đêm đặc quánh. Từng vết nứt trên kính giống như những cái mạng nhện khổng lồ, chỉ chờ một lực đủ mạnh để vỡ tung.

Tôi thậm chí đã mất hoàn toàn khái niệm về thời gian. Có lẽ đã hơn nửa tiếng trôi qua, nhưng với không khí ở đây, nó dài lê thê như cả một ngày. Những ánh đèn huỳnh quang lác đác trên trần chỉ đủ soi ra một mảng mờ nhạt, còn phần còn lại của hành lang thì vẫn đặc quánh bóng tối\dots

*PHỤT!*

\dots và tôi vừa mới nghĩ xong. Tất cả đèn vụt tắt cùng lúc, như thể muốn dằn mặt chúng tôi. Đúng thời điểm ghê.

``Đèn ở đây có vẻ chập chờn nhỉ?'' anh trai tôi vừa nói vừa lôi điện thoại ra bật đèn pin.

Tôi thở dài, làm theo: ``Em thấy nơi này như đang trêu ngươi chúng ta thì đúng hơn\dots''

Bóng tối nơi này đặc quánh, dày đặc đến mức cảm giác nó như đang đè lên ngực. Cái lạnh cứ len lỏi vào tận xương sống, khiến tay tôi run lên từng hồi. Đây chắc chắn là một trong những trải nghiệm rợn người nhất mà tôi từng có, vừa là một mê cung, vừa là một cái bẫy chết người.

\dots

*PHỤT!*

Một bóng đèn huỳnh quang phía trước bỗng lóe sáng, rồi lập lòe như đang hấp hối. Các đèn còn lại vẫn tối om. Cái ánh sáng yếu ớt ấy chẳng làm cho hành lang bớt đáng sợ hơn, ngược lại còn khiến nó trông quái dị hơn.

Tôi liếc sang Nii-chan: `Em bắt đầu cảm thấy không thể tin tưởng vào đèn ở đây nữa.''

Anh ấy khẽ gật đầu, tiếp tục bước với đèn pin điện thoại trong tay.

Chúng tôi dừng lại trước một lớp học khác. Ngay dưới cái đèn vừa chập chờn sáng đó, tôi cũng nhận ra có bóng người bên trong. Không phải chỉ một, hay hai người, mà là số người đủ trong một lớp học.

Hai cánh cửa kéo bị đóng chặt, chẳng thể nào mở ra. Nhưng\dots\ tất cả những người ở trong lớp đều ngồi ngay ngắn, quay mặt về phía bảng đen. Không ai nói, không ai nhúc nhích. Không có giáo viên. Không lý do gì để họ ở đây vào giờ này.

Trừ khi\dots\ họ không còn là ``người'' nữa.

Nii-chan đứng sững lại. Cũng dễ hiểu thôi, vì tôi vừa nhận ra trong số đó có cả cô gái tóc xanh và cô gái tóc xám ban nãy. Nhưng giờ thì còn nhiều gương mặt khác nữa: tóc nâu, tóc vàng với hai cục búi tròn, hay một cô gái tai cáo ở thế giới trước - mà thực chất là cái bờm cài tóc.

Chúng tôi đồng thanh: ``Bọn họ\dots\ cũng là học sinh ở nơi này sao?''

Anh trai tôi chớp mắt, bỗng chau mày: ``Kỳ lạ\dots\ sao anh thấy họ quen lắm\dots\ Anh đã từng thấy ở đâu rồi\dots''

Anh liệt kê một loạt mái tóc và khuôn mặt, rồi\dots\ lắc đầu bỏ cuộc. Tôi không mắng anh ấy như thường lệ, vì thật ra tôi cũng thấy lạ lùng y như thế.

Chợt, tôi để ý một tờ giấy dán bên ngoài lớp. ``Có cái gì đó ở đây này.''

``Một tài liệu à?''

``Không biết nữa. Nhưng hay là đọc thử đi.''

``Nếu nó giúp ta hiểu rõ hơn về nơi này, thì được.''

Chúng tôi tiến lại gần.

Chỉ có đúng một chữ duy nhất được viết bằng mực đỏ: ``\vie{\emph{\textbf{CHẠY.}}}''

Nhưng nhìn kỹ, những vệt đỏ ấy không đều, có chỗ loang lổ và chảy thành dòng. Mùi tanh thoảng lên\dots\ Máu?

Tôi định ghé sát hơn.

Tất cả những người trong lớp đồng loạt quay đầu về phía chúng tôi. Đồng loạt.

Những nụ cười nở rộng đến méo mó, để lộ hàm răng trắng nhợt, còn đôi mắt thì vô hồn như những con búp bê vỡ.

``Cái quái--''

\gif{e0o}{1}{125}

*XOẢNG!*

Tiếng kính vỡ xé toang không gian. Một ô cửa sổ rạn nứt, bóng tối tràn vào lớp học. Tiếng cười bắt đầu vang lên: cao, chói, và kéo dài như muốn chui thẳng vào màng nhĩ. Một tiếng cười từ một thứ gì đó đã hoàn toàn mất trí.

Tôi đã nghe kiểu này ở chiều không gian trước. Nó luôn đồng nghĩa với\dots\ \emph{thứ tồi tệ sắp xảy ra}.

``\textbf{Chạy!}'' Nii-chan không nói hai lời, kéo tay tôi lao về phía cánh cửa trượt ở hành lang trước mặt, thứ vừa mới mở từ lúc nào.

Chúng tôi vừa qua được thì ngay lúc đó\dots

*RẮC! RẦM! XOẢNG!*

Một đoạn trần nhà trước mặt đổ sập. Gỗ, bê tông, bụi mù mịt. Những thanh gỗ vỡ vụn cắm xiên xuống sàn, như những cái bẫy chông chết người.

``Chỗ này\dots\ sắp sập đến nơi rồi\dots'' Nii-chan vừa ho vừa nói. Tôi phải che mũi lại, mắt cay xè.

``Bám sát anh.''

Tôi nắm chặt lấy tay anh ấy. Từng bước một, chúng tôi dò dẫm qua hành lang đầy chướng ngại vật. Sàn nhà kẽo kẹt dưới chân, những thanh gỗ mục nát rên rỉ như sắp gãy bất cứ lúc nào.

Chúng tôi không dám chạy, vì chỉ cần một bước sai, rất có thể sẽ rơi xuống khoảng tối sâu hun hút bên dưới.

Không gian trước mặt như nuốt chửng tất cả âm thanh. Tiếng bụi rơi lạo xạo từ trần nhà cũng dần thưa đi, rồi biến mất hẳn. Cả tiếng bước chân của chúng tôi, vốn vẫn vang lên khe khẽ trên nền gỗ kẽo kẹt, giờ nghe rõ mồn một như bị phóng đại giữa khoảng không im lặng.

Tôi nuốt nước bọt. Không khí đặc quánh mùi gỗ mục, mùi đất ẩm và mùi gì đó hơi ngai ngái, khó chịu. Cứ mỗi hơi thở, bụi lại len vào mũi và cổ họng, buộc tôi phải đưa tay che miệng.

Chúng tôi vẫn bước chậm rãi, gần như dò dẫm trong lớp ánh sáng yếu ớt của đèn pin.

``Ồ.''

Nii-chan bất chợt khẽ thốt lên. Giọng anh không to, nhưng giữa cái sự tĩnh mịch này, nó vang lên đến mức tôi giật mình.

``Sao thế, Nii-chan?''

Tôi nhìn theo hướng anh chỉ. Một tờ giấy\dots\ không, chính xác hơn là một mẩu báo, đang được dán trên khung cửa gỗ đã cong vênh của một lớp học xuống cấp. Mấy ô kính cửa sổ phía sau nó vỡ vụn, viền quanh vẫn còn dính những mảnh thủy tinh xước xát.

Ánh đèn pin hắt lên tờ báo, để lộ ra dòng tiêu đề đã ố vàng theo năm tháng: ``{\vnmsans Thảm họa sạt lở đất ở thị trấn Aogami.}''

\dots Nghĩ lại thì, đúng là những gì chúng tôi vừa trải qua hoàn toàn có thể khớp với một trận sạt lở đất. Cảnh trần nhà đổ sập, bụi mù mịt, tiếng gỗ gãy\dots\ tất cả đều giống như một mô phỏng thu nhỏ của thảm họa ấy.

Chúng tôi ghé sát lại, xem kỹ từng chữ. Phần ngày tháng bị gạch xóa một phần, nhưng vẫn đủ để nhận ra: ``{\arial Tháng Bảy, năm 1946.}''

Tôi sững lại. Nghĩa là nó xảy ra trước cả trận động đất kia, tận những năm tám mươi\dots\ dù khoảng cách thời gian giữa các biến cố này vốn đã mơ hồ trong trí nhớ tôi.

Nii-chan bắt đầu đọc phần nội dung, giọng anh nhỏ và đều, xen giữa là vài chỗ chữ đã nhòe hoặc bị mờ đi:

\txtbx{\arial Vào lúc 22:10, tại thị trấn Aogami đã xảy ra một trận sạt lở đất nghiêm trọng, gây thiệt hại thảm khốc với XX người chết và mất tích.
  Nguyên nhân được các nhà điều tra xác định là do mưa lớn kéo dài, khiến kết cấu đất yếu đi. Tuy nhiên, một số người dân tin rằng đây là tai ương do vong linh đã ám nơi này từ rất lâu\dots}

22:10\dots\ tức 10 giờ 10 phút tối. Trùng hợp với thời điểm chiếc đồng hồ của lớp học đầu tiên chúng tôi ghé qua khi bị dịch chuyển tới đây.

Nếu đây là sự thật, tôi nghiêng về phía giả thuyết của những người dân kia nhiều hơn. Những gì chúng tôi đang trải qua rõ ràng không thể chỉ là sự trùng hợp ngẫu nhiên. Và nếu đúng như vậy\dots\ thì trường học này có thể chính là một trong những điểm khởi nguồn.

Tôi rùng mình. Một lời nguyền xuất phát từ những học sinh từng chuyển đến đây, nghe qua thì giống mấy câu chuyện hù dọa trẻ con, nhưng khi chính mình ở giữa nó, mọi thứ chẳng còn buồn cười chút nào.

``Anh không nghĩ thị trấn này lại hứng nhiều tai họa như vậy,'' Nii-chan khẽ thở dài.

Tôi cũng không đáp, chỉ lặng lẽ gật đầu.

Chúng tôi tiếp tục bước, tiếng giày va lên nền gỗ xen giữa tiếng kẽo kẹt của sàn.

Hành lang tối tăm, bụi mịt mù, từng thanh gỗ gãy nằm ngổn ngang dưới chân. Nó thực sự mang dáng dấp của một tòa nhà bị bỏ hoang hàng chục năm.

Bất chợt, giữa làn bụi mờ đục, một bóng người hiện ra.

Tôi khựng lại. Đó là dáng của một cô gái, đứng bất động, như thể đã ở đó từ lâu.

Tóc nâu buông dài gần chạm lưng, và\dots\ cái cài tóc đỏ.

\dots Chẳng phải chính là cô gái mà chúng tôi đã thấy ban nãy sao?

``Nii-chan!''

``Ừ. Anh thấy rồi. Sao cô ấy lại\dots\ đứng đó?'' Giọng anh hơi chùng xuống, đầy nghi hoặc.

Chúng tôi định tiến lại gần hơn, nhưng đột nhiên\dots

*VỤT!*

Bóng dáng ấy biến mất ngay trước mắt, để lại khoảng không trống rỗng. Kèm theo đó, đâu đó vang lên tiếng cửa mở, khẽ nhưng rất rõ.

``Thật đấy à\dots'' Nii-chan lầm bầm, đưa tay gãi đầu, khó hiểu với màn ``bốc hơi'' vừa rồi.

Tôi bắt đầu thấy bực, nhưng đồng thời cũng thấy bất an. Nếu cô ta có thể xuất hiện bất cứ lúc nào, ở bất kỳ đâu\dots\ thì có nghĩa là chúng tôi không bao giờ biết trước chuyện gì sắp xảy ra.

``Em nghĩ\dots\ chắc chúng ta sẽ gặp lại cô ta thôi,'' tôi nói, cố giữ giọng bình thản, nhưng tim đập nhanh hơn bình thường.

Không đợi anh đáp, tôi nắm tay kéo anh bước tiếp. Tiếp tục đi trên hành lang phía trước đang tối đen như đang chực chờ nuốt chúng tôi sâu hơn vào cái im lặng đặc quánh của nó.

\ct

Chúng tôi ``ngụp lặn'' trong làn sương bụi đặc quánh đến mức mỗi bước chân đi đều nghe tiếng xào xạc. Những tiếng ho sặc sụa của hai chúng tôi tiếp vục vang lên.

Mãi sau mới lờ mờ thấy một cánh cửa lớp đang hé mở. ``Hay là\dots\ vào đây tránh bụi đi?''

Tôi gật đầu. Ít ra bên trong kia có thể dễ thở hơn. Cả hai vội chạy vào.

Bên trong, không gian là một mớ hỗn độn. Bàn ghế xếp chồng chéo quanh một khoảng trống tròn trịa ở giữa, như thể từng có ai đó dựng nên một vòng tròn đặc biệt để làm gì đó rồi bỏ lại. Bóng đèn treo lủng lẳng nhưng tắt ngóm, khung cửa sổ vỡ nát như đã qua một trận chiến.

\img{e0p}

Chúng tôi lia đèn pin khắp nơi. Trên một chiếc bàn gần cuối lớp, tôi thấy một mảnh giấy học trò, chữ viết nguệch ngoạc, tiêu đề: ``{\vnmsans Lời đồn trong trường học - Mẩu 5.}''

Nội dung bên dưới không nhắc tới lời nguyền như các mẩu trước, mà lại giống một đoạn nhật ký ngắn:

\txtbx{\arial Không biết cứ mỗi ngày trốn tiết như thế thì mình nên làm gì ở thư viện nhỉ?\\
  Bình thường thì mình có thể tìm đọc các loại sách, nhưng\dots\ có lẽ tốt nhất mình nên tránh \emph{giờ điềm} tới.}

Tôi không hiểu ``giờ điềm'' nghĩa là gì, nhưng dựa vào những gì đã trải qua, tôi có linh cảm đó chẳng phải điều tốt đẹp gì.

``Anh tìm thấy cái này,'' Nii-chan chìa ra một mảnh giấy khác, từ bàn giáo viên. Nó giống một trích đoạn bài báo:

\txtbx{\vnmsans Những hiện tượng kỳ lạ ở thị trấn Aogami - Ký sự số 5\\
  \arial "Những câu chuyện dị thường ở đây đang lan ra với tốc độ chóng mặt.\\
  Tuy vậy, điều khiến tôi băn khoăn là tại sao chúng lại lan nhanh đến vậy\dots\\
  Không lẽ thị trấn này vốn đã bị những thế lực siêu nhiên bám rễ từ lâu rồi sao?\\
  Nếu quả thật vậy, thì nó lại là--"}

Phần sau bị xé mất.

Tôi có cảm giác người viết đã bắt đầu run tay khi gõ ra những dòng cuối. Có lẽ chính hắn cũng không dám viết tiếp. Nếu cứ thế này, việc gọi Aogami là ``thị trấn bị quỷ nhập'' cũng chẳng sai. Nhưng rốt cuộc, gốc rễ của mọi chuyện là gì?

\dots Và rồi mắt tôi dừng lại ở thứ đang hiện hữu ngay trước mặt từ nãy tới giờ: một bình hoa nhỏ đặt ngay ngắn trên chiếc bàn học sinh, bên trong là những bông hoa bỉ ngạn đỏ.

\img{e0q}

Thoạt nhìn, chúng thật sự rất đẹp. Những cánh hoa mỏng manh như tơ lụa, xòe ra thành những tia lửa đỏ rực giữa căn phòng u tối. Tôi từng nghe loài hoa này còn được gọi là hoa lycoris hay hoa thiên điểu đỏ, vốn gắn liền với hình ảnh đoàn tụ, lời hứa sẽ gặp lại nơi bờ bên kia. Ở một số nơi, nó còn mang ý nghĩa về sự bảo hộ linh hồn, như một ngọn đèn soi đường cho người lạc bước.

Nhưng đó mới chỉ là một nửa câu chuyện. Ở Nhật Bản, bỉ ngạn đỏ thường gắn với tang lễ, với chia ly vĩnh viễn và cái chết. Người ta bảo nó nở rực rỡ nhất trên những con đường dẫn xuống hoàng tuyền. Một vẻ đẹp đến mức đau lòng, như mỉm cười chào tạm biệt trước khi tan biến.

Đặc biệt, trong môi trường học đường, một bình hoa đặt trên bàn học sinh mang hai ý nghĩa hoàn toàn trái ngược: hoặc là một cách để tưởng nhớ một học sinh đã khuất, hoặc là một hình thức bắt nạt tàn nhẫn, như thể ngầm chúc người đó sớm chết đi. Và với những bông bỉ ngạn đỏ này\dots\ tôi không dám nghĩ đến ý nghĩa thực sự của nó.

``\dots Anh không nghĩ đây là loài hoa bình thường đâu,'' Nii-chan trầm giọng.

Tôi gật đầu. ``Nó vừa mang vẻ đẹp của sự đoàn tụ\dots\ nhưng cũng mang lại cảm giác lạnh sống lưng.''

``Ừ. Vừa là lời tiễn biệt, vừa như một dấu hiệu của oan hồn,'' anh ấy đáp.

Tôi bỗng rùng mình. Nếu ai đó cố ý đặt nó ở đây, thì chắc chắn không phải để trang trí cho vui. Nó giống một thông điệp, hoặc\dots\ một lời nguyền được gửi tới ai đó.

\dots

*XOẢNG!*

``Gì vậy!?''

``C-Chết! Nó vỡ rồi, Nii-chan!''

Bình hoa bất ngờ đổ nhào xuống nền gạch, vỡ tan, dù chúng tôi chưa hề chạm vào.

``Làm sao bây giờ!?''

``Anh biết chết liền à!? Hay là em lại--''

``Vớ vẩn! Không phải em! Chắc anh--''

``Anh còn chưa động tới nó!''

Giữa lúc chúng tôi còn tranh cãi\dots

``Hi hi hi\dots\ Hai người cũng là học sinh mới chuyển tới, phải không?''

Một giọng con gái khe khẽ vang lên từ phía sau.

Cả hai đồng loạt quay đầu, và rồi chết lặng.

``L-Là cô ta!''

``Cô ấy\dots\ vào đây từ khi nào vậy?''

Là cô gái tóc nâu, cài chiếc kẹp đỏ mà chúng tôi đã gặp trước đó. Nhưng lần này, ánh mắt cô ta cụp xuống, môi mím như nén một nụ cười, và không nói gì thêm. Cô chỉ đung đưa mũi chân như đang chơi một trò riêng của mình.

``Cô\dots\ là ai vậy?'' tôi hỏi, giọng khô khốc.

Không có câu trả lời.

Tôi vừa bước lên một bước\dots

*XOẸT!*

Cô biến mất. Lại nhanh như khi xuất hiện. Nhưng ngay sau đó\dots

*Ù Ù Ù\dots*

Một cơn chấn động bất ngờ ập tới, mạnh hơn hẳn cái trước. Trần nhà rung bần bật, từng thanh gỗ nứt toác rồi rơi ầm ầm xuống sàn như mưa. Bụi mù đặc quánh, mùi gỗ mục và sắt rỉ xộc thẳng vào mũi khiến tôi nghẹn lại.

``Chạy thôi!'' cả hai chúng tôi hét lên gần như cùng lúc.

Tiếng bước chân dồn dập vang trên sàn gỗ ọp ẹp khi chúng tôi lao ra hành lang. Nhưng chưa kịp hoàn hồn, cả dãy hành lang cũng bắt đầu sụp xuống từng đoạn. Tấm kính cửa sổ bên lớp học nổ tung thành hàng trăm mảnh, bay vèo vèo như dao lam sát bên tai, rơi loảng xoảng xuống nền. Sàn dưới chân rung bần bật rồi nứt ra, tạo thành những hố sâu đen ngòm.

``\textbf{Gyaah!}''  tôi thét lên khi bàn chân hụt xuống một khoảng trống lạnh buốt.

``\textbf{Cẩn thận!}'' Nii-chan kịp chộp lấy cổ tay tôi, kéo mạnh lên. Cánh tay anh run lên vì lực kéo, nhưng vẫn giữ chặt cho tới khi tôi đặt được cả hai chân lên sàn.

``Em\dots\ ổn chứ?''

Tôi gật đầu, ho sặc sụa vì bụi. Thứ bụi mịn này bám dính khắp mặt, vừa rít trong cổ họng, vừa khiến mắt cay xè. Không còn đường nào khác ngoài việc tiếp tục chạy, chúng tôi vừa bật đèn pin soi đường vừa bám vào nhau. Mỗi bước chân đều phải dè chừng vì những tấm ván mục có thể gãy bất cứ lúc nào.

``Đằng trước! Cửa trượt kìa!'' Nii-chan chỉ về một cánh cửa mở hé, ánh sáng mờ mờ lọt qua khe.

Chúng tôi lao tới, nhưng cả đoạn hành lang phía sau sập xuống ngay khi vừa nhảy qua ngưỡng cửa. Tiếng gỗ gãy răng rắc hòa cùng tiếng kính vỡ chát chúa, tất cả biến thành một mớ âm thanh hỗn loạn nuốt chửng khoảng không sau lưng.

``Anh không sao chứ?''

``Không sao\dots\ nhưng\dots\ trời ạ\dots'' anh thở hắt ra, cố trấn tĩnh.

Chúng tôi phủi tạm lớp bụi và mảnh vụn bám trên quần áo, rồi nhìn quanh hành lang mới bước sang. Nơi này cũng chẳng khá hơn: trần nhà thủng lỗ chỗ, vài thanh xà ngang gãy đôi chênh vênh như chực rơi. Nền gỗ mục nát để lộ vài khoảng trống đen ngòm bên dưới, chỉ vừa đủ cho một người bước lọt qua.

\img{e0r}

Điều may mắn duy nhất là ở đây không còn bụi mù mịt, nhưng không vì thế mà bớt nguy hiểm. Mỗi bước đi đều kèm theo tiếng *RẮC* khe khẽ, như lời cảnh báo sẵn sàng nuốt chửng chúng tôi.

Hai bên hành lang là những lớp học đóng kín, nhiều phòng vẫn sáng đèn, nhưng bên trong chỉ có một chiếc bàn trơ trọi và một lọ hoa đã héo quắt. Cửa bị phong ấn bằng những dải giấy cũ, chữ trên đó đã nhòe mờ nhưng vẫn tỏa ra thứ cảm giác lạnh gáy. Dù tò mò, tôi chẳng dại gì mà gỡ thử.

\img{e0s}

Chúng tôi bước chậm, tai vẫn căng ra để nghe ngóng. Tim tôi đập dồn dập, và rồi\dots\ một ý nghĩ len lỏi vào.

Không biết có phải là do tôi tưởng tượng, nhưng tôi thấy\dots\ mọi thứ đang diễn ra giống như một đoạn phim tua lại của trận sạt lở đất kia. Cảm giác rung chấn, tiếng gỗ nứt, bụi bay vừa rồi\dots\ chẳng khác nào đang sống lại khoảnh khắc đó.

Chắc chắn không thể chỉ là trùng hợp.

``Saikyou? Ở dưới chân em có gì kìa,'' giọng Nii-chan bất ngờ cắt ngang luồng suy nghĩ của tôi.

``S-Sao ạ?'' tôi nhìn xuống theo hướng anh chỉ. Dưới chân tôi là một tờ giấy nhỏ, xộc xệch, cạnh mép bị sờn như đã nằm đây hàng chục năm. Loại giấy quen thuộc ấy\dots\ giống hệt mấy mẩu ghi chú của gã nhà báo mà chúng tôi từng thấy ở ``thế giới kia''.

Tôi nhặt lên.

\txtbx{\vnmsans Ghi chú của nhà báo - Tờ số 4\\
  \arial "Từ hôm qua, tôi có cảm giác như có ai đó đang dõi theo từng cử động của mình\dots\ Thực sự lạnh sống lưng. Điều lạ là, cảm giác này chỉ bắt đầu từ khi tôi đến đây, hay từ khi tôi bắt đầu điều tra vụ này?\\
  \dots Không, chắc là do tôi mệt thôi. Phải vậy\dots"}

Tôi rùng mình. Không, chắc chắn gã ta không tưởng tượng. Ở thị trấn này, thứ ``theo dõi'' ấy hoàn toàn có thể là có thật.

``\dots Anh cũng thấy lạnh sống lưng,'' Nii-chan khẽ nói. ``Em có thấy kiểu\dots\ giống như đang bị ai đó nhìn chằm chằm không?''

Tôi nuốt khan, từng nước bọt như đang ứ nghẹn lại đằng sau cổ: ``Có. Em cảm giác nó ở ngay sau lưng, nhưng quay lại thì lại chẳng thấy gì\dots\''

Cả hai im lặng vài giây. Chỉ còn tiếng đèn pin lách tách và tiếng bước chân. Tôi đặt lại mẩu giấy xuống, cố tập trung vào việc tìm đường ra.

Nhưng vừa quay lại, tôi thấy Nii-chan đã đi tít sang phía cuối hành lang, tay cầm một thứ gì đó.

``N-Này! Đợi em với!'' tôi luống cuống bước qua những tấm ván rung rinh để tới chỗ anh.

Anh giơ ra trước mặt tôi một tập giấy bị rách góc. ``Xin lỗi, tại anh mải đọc cái này\dots\ Một bản báo cáo về vụ sạt lở đất.''

Tôi cầm lấy, lướt mắt qua những dòng chữ mờ:

\txtbx{\vnmsans Báo cáo kết quả về vụ việc sạt lở đất ở thị trấn Aogami - Bản ghi số 1\\
  \arial "Vào ngày X tháng Bảy năm 19, trận sạt lở đất xảy ra đêm hôm trước đã gây ra thiệt hại nghiêm trọng cho thị trấn vốn đã ít dân này. Trong số đó, ngôi trường trung học của thị trấn đã sụp đổ hoàn toàn\dots"}

Những tấm phong ấn ngoài kia bỗng trở nên dễ hiểu. Ngôi trường này đã từng chứng kiến cảnh tan hoang như trong báo cáo, và giờ nó chỉ còn là cái xác bị bỏ lại, phủ đầy bụi và ký ức chết chóc.

``Có vẻ cũng chẳng giúp mình hiểu rõ hơn chuyện gì đang diễn ra\dots'' Nii-chan thở dài.

Tôi khẽ gật, nhưng trong lòng lại thấy thứ gì đó đang dần khớp nối lại. Một bức tranh ghép vừa quen vừa đáng sợ.

Chúng tôi tiếp tục dò dẫm dọc theo hành lang thủng lỗ chỗ, mỗi bước chân đều tạo ra những âm thanh lạo xạo, hòa cùng tiếng kẽo kẹt bất an từ mấy thanh gỗ trên trần và dưới sàn. Mùi gỗ mục và bụi mốc len lỏi vào từng hơi thở, khiến cổ họng tôi khô khốc. Nơi này vừa nguy hiểm, vừa chẳng khác gì một cái bẫy có thể sập xuống bất cứ lúc nào mà không ai trong chúng tôi ngờ tới.

Chúng tôi đi mãi cho đến cuối hành lang, và rồi\dots

\dots cô ta lại ở đó. Cô gái tóc nâu, đứng đủng đỉnh như thể đang chờ sẵn, nét mặt chẳng hề có chút lo lắng nào. Trái lại, khóe môi cong lên thành một nụ cười như thể đây là trò chơi của riêng cô ta.

Ngay khi chúng tôi vừa chạm đến khoảng cách gần với cánh cửa trượt, cô ta ngẩng lên, mắt sáng lên một tia hứng thú kỳ lạ, và chậm rãi bước về phía chúng tôi.

``Các cậu thấy trường học có vui không?'' giọng nói ấy vang lên, nhẹ nhàng nhưng rợn gai ốc.

Tim tôi như bị bóp chặt. Lời nói thì nghe như một câu hỏi xã giao, nhưng ở đây, trong bầu không khí đặc quánh này, nó giống như một câu trêu ngươi đầy ẩn ý.

``\textbf{Vui cái con khỉ mốc!}'' tôi bật quát, giọng cao hơn bình thường, như để át đi nỗi bất an đang dâng lên. ``\textbf{Định dọa cho bọn này chết khiếp à!?}''

Nii-chan vội đặt tay lên vai tôi, giọng ôn hòa: ``Thôi nào, đừng nóng\dots\ Anh không nghĩ cô ấy muốn làm hại chúng ta đâu.''

Tôi quay phắt sang: ``Anh nói cái gì chứ!? Chúng ta vừa suýt bỏ mạng ở đây đấy!''

Tôi cảm thấy máu dồn lên mặt, vừa muốn lao tới tóm cổ cô ta, vừa muốn lùi lại vì cảm giác\dots\ có gì đó rất sai.

Nii-chan phải dùng cả hai tay chặn trước ngực tôi: ``Bình tĩnh, Saikyou!''

Cô gái vẫn đứng đó, như thể chuyện tôi hét vào mặt chẳng hề liên quan. Cô mỉm cười, hơi nghiêng đầu: ``Tớ dám chắc là các cậu thấy vui mà.''

Câu nói ấy như cố tình đổ thêm dầu vào lửa. Tôi siết chặt nắm đấm, vừa tức giận vừa thấy da gà nổi khắp lưng. Cô ta\dots\ trng không bình thường chút nào.

Nii-chan, vẫn giữ tôi lại, lên tiếng: ``Nhưng\dots\ cô thực sự là ai vậy?''

*BÍNH BONG\dots\ BÍNH BONG\dots*

Tiếng chuông báo vang vọng khắp nơi, đột ngột và vô định, khiến cả hai chúng tôi khựng lại. Tiếng ngân kéo dài như từ một khoảng trống sâu thẳm nào đó, rồi lại lặp lại.

Cô gái vẫn nở nụ cười đều đều, ánh mắt như xuyên qua cả tiếng ồn: ``\dots Hết tiết rồi đó. Thôi nào, ta hãy cùng về thôi.''

``Vấn đề là chúng ta về kiểu gì chứ\dots'' Nii-chan thở dài bất lực.

``\textbf{Về cái quái gì!? Chúng tôi còn chẳng biết đây là đâu nữa mà về!}'' tôi hét, cảm giác như giọng mình vang vọng trong khoảng không mờ đặc ấy.

Cô ta chẳng đáp, chỉ ngoảnh lại, bước về phía cánh cửa trượt, rồi mở ra. Không quay đầu, không nói thêm một câu.

``\textbf{Ê! Đứng lại đó!}'' tôi hét với theo, vùng vẫy khỏi tay Nii-chan.

Nhưng chỉ vài bước, cô đã biến mất vào làn sương mù dày đặc phía bên kia hành lang.

``Bỏ em ra, Nii-chan! Em phải-''

``Không, đừng! Anh còn chẳng chắc cô ấy có phải là nguồn cơn của mấy thứ này không nữa!''

Phải mất một lúc tôi mới dứt khỏi vòng tay anh, nhưng khi lao tới\dots\ chẳng còn gì ngoài sương và bóng tối, cùng vài mảnh trần nhà vừa rơi lộp bộp xuống sàn.

Tôi cắn môi, tim vẫn đập thình thịch. Tôi vừa giận điên lên, vừa thấy lạnh sống lưng\dots\ như thể ánh mắt cô ta vẫn còn dõi theo từ đâu đó.

Nii-chan bước tới, giọng trầm: ``Đừng để cảm xúc lấn át, kẻo em tự hại mình.''

``\dots Vâng,'' tôi đáp lí nhí, nhưng trong đầu lại lẩm bẩm: Nếu gặp lại, mình sẽ bắt cô ta nói rõ mọi chuyện, dù có phải ba mặt một lời.

*RÈ*

``\textbf{Eeeek!}'' tôi bật kêu, khi ánh đèn chớp sáng lên bất ngờ. Trong ánh sáng ấy\dots\ bóng cô gái hiện ra ở cuối hành lang.

Tim tôi siết chặt, cảm xúc uất ức trào dâng. Được rồi, lần này\dots

*PHỤT!*

Bóng tối nuốt chửng mọi thứ. Không còn ai ở đó nữa.

Tôi đứng chôn chân vài giây, hơi thở gấp gáp, rồi tự ép mình nuốt xuống cục tức và nỗi sợ vừa dâng lên.

``Em có ổn không vậy? Trông em như sắp lao vào vồ mồi ấy\dots'' Nii-chan gãi đầu, vừa lo vừa ngán.

``\dots Không sao,'' tôi hít sâu, tự trấn tĩnh. ``Giờ thì\dots\ quay lại nhiệm vụ chính. Tìm đường ra.''

``Ờm\dots\ được rồi. Tiện thể thì\dots''

``Sao nữa?'' tôi quay sang, vẫn chưa hết bực từ vụ vừa rồi.

``\dots Anh lại nhặt được một cuốn ghi chú nữa này.'' Giọng Nii-chan bình thản, nhưng cái cách anh chìa mảnh giấy ra trước mặt tôi khiến tôi lập tức chú ý.

``Ở đâu vậy?''

``Ở trên cái bàn bên ngoài lớp học đây này.''

Tôi giật lấy mảnh ghi chú từ tay anh, lia mắt xuống tiêu đề: ``{\vnmsans Ghi chú của nhà báo - Tờ số 5.}''

Cảm giác gai người chạy dọc sống lưng. Mấy tờ ghi chú trước vốn đã đủ rùng rợn, và tôi biết chắc gã nhà báo này hẳn đã moi ra thứ gì còn khủng khiếp hơn về cái thị trấn chết tiệt này.

Tôi đọc tiếp:

\txtbx{\arial "Tôi đã điều tra được từ các tài liệu cổ xưa về nơi này. Từ rất lâu trước đây, thị trấn này đã tồn tại một hủ tục quái dị: để xoa dịu thần linh, người dân phải thế mạng một người từ bên ngoài.\\
  \dots Liệu điều này có liên quan tới vụ các nữ sinh mất tích gần đây không?
  Ngoài ra, hôm nay tôi lại cảm giác như có một ánh mắt dõi theo mình, mặc dù xung quanh không hề có ai."}

Tôi ngẩng lên: ``Nghe giống như ở đây từng có một nghi lễ hiến tế người sống ấy nhỉ, Nii-chan?''

Anh gật đầu, nhăn mặt: ``Dường như là vậy\dots\ giống kiểu mấy bộ lạc xưa ở rừng sâu hay trên núi cao. Nghĩ thôi cũng rợn gáy.''

Tôi gấp tờ giấy lại, nhét vào túi áo. ``Rợn thì rợn thật, nhưng càng rợn càng phải biết nguyên nhân.''

\ct

Chúng tôi bước sang căn lớp học kế bên. Ánh đèn pin lia qua một khung cảnh hỗn độn: bàn ghế xô lệch, giấy tờ rơi rải rác, đồ treo tường ngã nhào xuống sàn. Một chiếc đồng hồ treo tường đứng im, kim ngừng ở đâu đó từ rất lâu.

Ngay dưới đồng hồ là một mẩu giấy khác, chữ đã nhòe và giấy ố vàng. Không phải bản in, mà là chữ viết tay, khó đọc và ngoằn ngoèo. Tôi ghé sát đọc:

\txtbx{\vnmsans Báo cáo điều tra vụ việc tại ga Aogami - bản ghi số 1\\
  \arial Tháng Bảy, năm 2018.\\
  Số lượng người mất tích tại thị trấn Aogami đã tăng nhanh đột biến trong những tháng vừa qua. Điểm chung duy nhất giữa các vụ việc là--"}

Dòng chữ bị cắt ngang bởi một vết xé.

Tôi nhếch mép. Không cần đọc hết, tôi cũng đoán ra. Chắc chắn là dính líu tới ngôi trường quái quỷ này\dots\ hoặc thứ gì còn lớn hơn đang nuốt chửng cả thị trấn.

Bên cạnh đó, trên bàn có một cuốn vở cũ kỹ, bìa không ghi tên, mép giấy bị bong tróc. Tôi định lật bừa thì ánh đèn chiếu vào một dòng chữ nhỏ viết ở góc: ``Giờ thì mình cũng có thể giống như cậu rồi. Mình sẽ gặp cậu sớm thôi. Mình thực sự xin lỗi cậu nhiều nha, Shino.''

Tôi đọc thầm cái tên ấy.

Shino\dots?

Trí óc tôi lập tức lục lại ký ức. Cô gái với cuốn sách giáo khoa mới tinh, ánh mắt như nhìn xuyên vào tôi\dots\ và cái giọng tươi cười hỏi ``Trường học có vui không?''

\dots Khoan đã.

Từ khi lạc vào đây, không chỉ một lần mà rất nhiều lần, hình bóng ấy xuất hiện trước mắt chúng tôi. Luôn là mái tóc dài nâu, luôn là nụ cười không đổi, và\dots\ luôn biến mất khi tôi định tiến lại gần.

Tôi nuốt khan, như không muốn tin vào điều đó. Nhưng có lẽ nào\dots?

Không thể nhầm được. Đó chính là Misora Shino.

Nếu đúng là cô ta\dots\ thì có thể toàn bộ ngôi trường này bị ám là vì cô ta. Hoặc ít nhất, cô ta có liên quan trực tiếp.

Câu hỏi ``Trường học có vui không?'' lại vang lên trong đầu tôi. Những học sinh mới chuyển trường thường bị hỏi câu đó\dots\ và nếu giả thuyết đúng, nghĩa là chúng tôi đang bị coi như ``học sinh mới'' ở nơi này?

Tôi rùng mình. Liệu đây có phải là một trò tái hiện méo mó của một lễ nhập học\dots\ hay là nghi thức mở màn cho nghi lễ hiến tế?

Không, tôi chưa thể kết luận ngay được. Nhưng tôi chắc chắn một điều: lần tới gặp lại cô ta, tôi sẽ lôi cho ra hết câu trả lời. Và tôi sẽ nói giả thuyết này với Nii-chan, dù anh có tin hay không.

``\textbf{Nii-chan!}'' tôi gọi gấp, giọng hơi cao vì kích động.

Anh tôi quay lại, trên tay là bốn mảnh giấy. ``Anh lại tìm thêm mấy cái tài liệu này.''

``Cái gì đây?'' tôi hơi sốt ruột.

``Hai bản báo cáo và hai tờ ghi chú của gã nhà báo kia. Nhưng mà\dots'' anh ngập ngừng, liếc sang tôi rồi lại cúi xuống.

``Sao vậy?`` tôi cau mày.

Anh chỉ thở dài, né tránh: ``\dots Không có gì. Mà em vừa định nói gì à?``

Tôi ngập ngừng một nhịp, định bật ra cái tên kia, nhưng ánh mắt anh trông như đang muốn vội chuyển chủ đề.

``Em đọc cái báo cáo này trước đi.'' Anh chìa cho tôi một mẩu báo cáo, như thể cố tình đẩy tôi vào một hướng khác.

\dots

Lần này là giấy in, chắc chắn thuộc cùng tập với bản trước, thậm chí chỗ rách ở phần đầu còn khớp hoàn toàn. Có thể hai tờ này vốn là một.

\txtbx{\vnmsans Báo cáo kết quả về vụ việc sạt lở đất ở thị trấn Aogami - Bản ghi số 2\\
  \arial "Thời điểm sạt lở đất diễn ra, đáng nhẽ sẽ không có thiệt hại nào về người, nhưng bất ngờ thay, cơ quan chức năng đã phát hiện ra những thi thể vẫn còn mặc đồng phục học sinh ở dưới đống đổ nát."}

Một luồng lạnh chạy dọc sống lưng tôi.

Học sinh? Ở đó? Lúc nửa đêm? Hoặc là họ bị kéo đến đây\dots\ hoặc là đã có hai vụ sạt lở, với một ngoài đời, một ở ``thế giới'' này.

Ý nghĩ ấy càng khiến cái tên Shino lởn vởn trong đầu.

\dots

``Còn tờ báo cáo kia thì sao?'' tôi hỏi. Nii-chan không trả lời.

Anh vẫn cắm cúi vào hai mẩu ghi chú, như thể bị chúng hút hồn.

``Kyuen-nii-chan!'' tôi gọi mạnh hơn.

``G-Gì nữa?''

``Cái tờ báo cáo kia! Đưa đây cho em xem!''

``À, ờ\dots\ phải rồi.''

Nii-chan đưa tờ còn lại. Một tờ viết tay, cùng loại với bản tìm thấy ở ga tàu. Có lẽ đây là phần tiếp.

Báo cáo điều tra vụ việc tại ga Aogami - Bản ghi số 2
``--khi họ tiếp cận tới nhà ga Aogami, họ đã không thể liên lạc được.''

Vỏn vẹn một câu, nhưng lại khiến lòng tôi trĩu xuống. Cắt liên lạc\dots\ chẳng khác gì biến mất khỏi thế giới thực.

``Anh nghĩ sao?'' tôi hỏi.

Kyuen vò đầu bứt tai: ``Anh đọc hai cái tờ này mấy lần rồi mà chẳng hiểu gã ấy muốn nói gì nữa. Nhìn chữ thôi cũng thấy ghê.''

Tôi chìa tay: ``Đưa em xem.''

``{\vnmsans Ghi chú của nhà báo - Tờ số 6}'' chỉ toàn chữ ``Không'' (\ruby{嫌}{イヤ}だ hay イヤ) lặp đi lặp lại, như một câu thần chú tuyệt vọng.

Còn ``{\vnmsans Ghi chú của nhà báo - Tờ số 7}'' thì còn\dots\ khủng khiếp hơn:

\txtbx{\arial
  "青ノ神抗ウナ……せイシン二アラガウナ……せイシン二……アハハハハハハハハハハ―"
  (Đừng chống lại Aogami\dots\ Đừng chống lại những linh hồn\dots\ ở đây\dots\ Đừng\dots\ Ahahahahaha--)
}

Tôi lạnh toát người. Chắc đây là thứ cuối cùng gã viết trước khi biến mất\dots\ hoặc chết.

Tôi ném chúng xuống bàn, tay vẫn run: ``Anh\dots\ kiếm ở đâu vậy?''

``Một cái trên ghế, một cái ở ngoài kia. À, hình như trên bàn kia còn có cái gì nữa.''

Chúng tôi rọi đèn tới nó. Một bông bỉ ngạn đỏ nằm chơ vơ.

Bỉ ngạn\dots\ loài hoa tượng trưng cho cái chết. Và trong đầu tôi, hình ảnh Shino lại hiện rõ hơn bao giờ hết.

Hay là một ai đó khác mà chúng tôi còn không biết tới?

Thêm nữa, tại sao ả ta lại bị mọi người ghét đến thế?

Đến giờ tôi vẫn không không hiểu được, và cũng không thể nào biết được.

Có vẻ như chúng tôi cũng sắp phát điên vì cái nơi quái thai này rồi. Quả đúng như những gì gã ta đã viết, nơi này đúng là có ma ám!

Phải thực sự tìm đường rời khỏi đây thật nhanh!

``Đi thôi nào, Nii-chan!'' tôi hối thúc anh trai tôi chạy khỏi đây. Không cần để chần chừ gì thêm, chúng tôi nhanh chóng rời khỏi căn phòng này và bắt gặp một cánh cửa trượt.

*XOẠCH*

Cánh cừa được mở một cách trơn tru mà không gặp phải trở ngại gì. Tuy nhiên\dots

\dots thứ hiện ra sau đó lại khiến tim tôi chùng xuống: một hành lang nữa, cũng bị đổ sập từng mảng, và vẫn ngập ngụa trong bóng tối.

Chúng tôi đã quen với cảnh này, nhưng rồi\dots

``Hihihi\dots'' ``Haha\dots''

\dots những âm thanh khúc khích vang lên từ xa, âm sắc trẻ con nhưng vỡ vụn như vọng qua lớp nước đục.

Tiếng cười không phải của ai đang đứng gần, mà như thể chúng thấm vào tường, len lỏi qua trần, bò sát dọc ống tai.

Tôi bám chặt tay Nii-chan, cảm nhận từng nhịp run khẽ của anh. Người mạnh mẽ nhất trong nhóm mà còn như vậy, huống gì tới tôi.

``Biến đi\dots'' ``Lại học sinh mới à? Haha\dots\ đúng là đần.'' ``Mày không xứng đáng ở đây, nhóc ạ.''

Giọng nói đan xen với tiếng cười, lúc gần lúc xa, như cả hành lang đang sống, đang ghim lấy chúng tôi bằng những đôi mắt vô hình.

Tôi chợt hiểu: ở đây, ``khách'' không bao giờ được chào đón. Những linh hồn này dường như mang đầy oán hận, và thứ ``nghiệp'' của họ đã biến ngôi trường thành một hầm mộ biết nói.

Khi ánh đèn pin lia xuống sàn, có gì đó lấp ló dưới lớp bụi.

Hai mẩu giấy, trong đó có một tờ in máy và một tờ viết tay. Chúng nằm vắt ngang nhau như dấu ``X'' của định mệnh.

Tôi cúi xuống, nhặt lên. Giấy lạnh buốt, như vừa được lấy ra từ một ngăn đá. Chúng tôi vừa đi vừa đọc.

\txtbx{\vnmsans Báo cáo kết quả về vụ việc sạt lở đất ở thị trấn Aogami - Bản ghi số 3\\
  \arial Kể từ sau vụ tai nạn đó, nhiều hiện tượng kỳ lạ bắt đầu xảy ra tại thị trấn Aogami. Người ta chứng kiến những hồn ma từ ----------- vất vưởng khắp nơi, và những người sợ hãi đã bắt đầu rời bỏ thị trấn.}

Lưng tôi toát mồ hôi. Ở đây, ``vất vưởng khắp nơi'' nghĩa là\dots\ có thể chúng tôi đang bị nhìn từ mọi phía ngay lúc này.

Mẩu giấy viết tay thì có phần đầu bị xé, nhưng mép rách trùng khớp với các bản trước, ghép lại thành một chuỗi u ám.

\txtbx{\vnmsans Báo cáo điều tra vụ việc tại ga Aogami - Bản ghi số 3\\
  \arial "Những người mất tích ở thị trấn Aogami biến mất không còn dấu vết. Điều khó hiểu hơn, các nhà chức trách không tìm được bất cứ manh mối nào khác tại ga Aogami này."}

Tôi ngẩng lên, hướng mặt về anh trai đang lọ mọ soi đường. ``Nii-chan\dots\ nếu nối các mẩu này lại, thì chuyện ở trường và chuyện ở thị trấn đều bắt nguồn từ cùng một thứ. Và\dots''

Anh gật đầu. ``Cô gái đó. Shino.''

Cái tên vừa thoát ra đã khiến hành lang như lạnh thêm một bậc. Có lẽ Nii-chan cũng đã luận ra được cô gái ấy là người mà tôi cần tìm rồi.

Chúng tôi tiếp tục bước, nhưng càng đi càng nhận ra: những dãy hành lang này\dots\ lặp lại.

``Anh có thấy lạ không khi cứ đi qua mấy hành lang giống hệt nhau?'' tôi hỏi.

Anh tôi khựng lại, đôi mắt ánh lên vẻ dò xét. ``Đúng thật. Một trường bình thường đâu có thế\dots''

``Có lẽ nào\dots''

``\dots Chúng ta bị kẹt trong một vòng lặp,'' cả hai đồng thanh. Và nếu thế, thủ phạm chỉ có thể là Shino. Chung quy lại, vẫn chỉ xoay quanh ả ta mà thôi.

Tôi nhớ lại cuốn tạp chí đọc ban nãy: Một khi bị mắc kẹt tại đây thì sẽ không thể thoát ra, trừ khi làm một việc đặc biệt.

Nhưng việc đó là gì? Chúng tôi đâu phải thầy trừ tà.

``Được rồi, tóm lại nhé,'' tôi bắt đầu.

``Một cơn động đất xảy ra\dots''

``\dots giết chết mọi người ở đây, để hồn ma họ lang thang, và gọi những người đến sau là `học sinh mới'.''

``Những kẻ như thế đã làm gì đó khiến mọi chuyện rối tung, kèm theo những hiện tượng quái dị khắp nơi\dots''

``\dots từ ga tàu, trường học, tới cả thị trấn, khiến nơi này bị xa lánh.''

Nhưng vẫn còn hai câu hỏi treo lơ lửng như dây thòng lọng vẫn còn đau đáu: con tàu ở ga Aogami và bức tranh ``kiệt tác'' kia, rốt cuộc để làm gì?

Tôi không biết. Nhưng tôi chắc một điều: có ai đó đang cười khúc khích sau lưng chúng tôi ngay lúc này, đang chứng kiến những khoảnh khắc mà chúng mà tôi đang dần mất trí vì cái nơi khỉ ho cò gáy chết dẫm này.

Chúng tôi vẫn lầm lũi bước trong hành lang tối tăm, tiếng bước chân khẽ vang vọng giữa không gian đặc quánh như mùi bụi bặm trộn lẫn hơi ẩm mốc. Bỗng ánh đèn pin lia vào một tấm biển lớp, chữ khắc trên đó rõ ràng hơn hẳn so với những biển lớp mờ mịt trước đây.

``Phòng\dots\ Mỹ thuật?'' tôi lẩm bẩm.

Một luồng khí lạnh lướt qua gáy tôi khi nhớ lại những mẩu ghi chép trước đó. Đây chẳng phải là nơi đã sản sinh ra ``kiệt tác'' kia, và cũng là nơi mà cô gái họa sĩ ấy bỏ mạng sao? Nếu đúng vậy\dots\ có khi nào bức tranh đó vẫn ở đây, im lặng chờ ai đó tìm đến?

Ánh sáng bên trong phòng le lói, mờ đến mức gần như tan vào bóng tối, nhưng lại vừa đủ để kích thích trí tò mò của tôi. Không biết vì sao, nhưng tim tôi đập nhanh hơn\dots\ như có thứ gì đó đang gọi tên tôi từ bên trong.

\dots Tôi phải nhìn nó.

Không suy nghĩ, cũng chẳng báo trước, tôi bước tới, đẩy cánh cửa khẽ phát ra tiếng két khô khốc. Lúc này, mọi lời của anh trai tôi ở phía sau bị tôi bỏ ngoài tai.

``Ơ-Ơ!? Saikyou!?'' giọng anh vang lên đầy hoảng hốt, nhưng tôi đã bước qua ngưỡng cửa rồi.

\pov{Haragen Kyuen}

``Saikyou! Này, em định đi đâu thế!?'' tôi vội vàng chạy theo, tay với ra nhưng chạm vào khoảng không.

Không một lời đáp lại. Chỉ có bước chân của em ấy, chậm rãi nhưng kiên quyết, tiến sâu vào căn phòng mờ tối. Cảm giác bất an bỗng trào dâng, chẳng phải chúng ta vừa mới kết luận được mọi chuyện ở đây nguy hiểm thế nào sao?

''\dots Nii-chan, anh nhìn này.''

Giọng em ấy vang lên bình thản đến lạnh người, khác hẳn vẻ sợ sệt ban nãy. Tôi bước vào và thấy Saikyou đang đứng bất động trước một giá vẽ. Trên đó là một tấm toan trắng, cũ kỹ đến mức mép đã sờn, bề mặt hơi ngả vàng vì thời gian.

Khác với các lớp học bình thường, phòng Mỹ thuật của trường có vẻ rộng rãi hơn hẳn. Thay vì những dãy bàn ghế ngay ngắn, ở đây là những giá vẽ được xếp thành hình bán nguyệt, mỗi giá có một chiếc ghế đẩu nhỏ kèm theo. Trên tường treo đầy những bức tranh của học sinh, từ phong cảnh cho đến chân dung, tĩnh vật. Góc phòng là một tủ đựng dụng cụ vẽ với đủ loại màu vẽ, cọ, bút chì. Không khí trong phòng thoang thoảng mùi sơn dầu và mực tàu.

Điều đặc biệt nhất là ở chính giữa phòng: một giá vẽ lớn hơn hẳn những cái còn lại, trên đó là một tấm vải trắng tinh khôi. Nó nổi bật giữa căn phòng như một ngọn hải đăng giữa biển đêm.

Tôi tiến tới gần hơn bức vẽ mà cô em gái đang đứng, nhưng khi đang định gọi tên để hỏi, thì bỗng nhiên, một cơn nhói nhẹ lan khắp đầu tôi. Không rõ từ đâu, nhưng nó khiến tôi khẽ nhăn mặt. Tôi định mở miệng thì Saikyou đã khụy nhẹ, hai tay ôm đầu.

``Nii-chan\dots\ sao tự dưng đầu em đau quá vậy\dots!'' giọng em run run, nhưng vẫn dán mắt vào khung vẽ.

\begin{wrapfigure}[21]{r}{10cm}
  \includegraphics[width=10cm]{../../../../Image/Book?/canvas.png}
\end{wrapfigure}

Toàn bộ căn phòng bỗng nhuộm đỏ. Tấm toan trắng biến mất, thay vào đó là một bức tranh hiện lên, một thiếu nữ, được phác bằng nét chì đen, điểm những vệt đỏ tươi trông như máu loang. Đôi mắt cô ta nhìn thẳng về phía chúng tôi. Nụ cười\dots hay tiếng gào thét? Tôi không phân biệt nổi.

Một luồng cảm xúc lạ lùng đâm thẳng vào ngực tôi, và cả Saikyou.

``\dots Sao\dots\ em thấy đau trong lòng vậy?'' em nói, giọng nghẹn lại, tay ôm chặt lấy ngực. Tôi cũng cảm nhận được, nhưng ở em, nó mạnh mẽ hơn nhiều: một nỗi đau tinh thần nặng trĩu, vừa uất ức, vừa tuyệt vọng, như đang bị nhấn chìm.

``Không\dots\ tại sao chứ\dots\ \textbf{Tại sao!?}'' tiếng hét của em vang lên xé toạc sự im lặng, khiến tôi giật mình. Em ngẩng đầu, ôm đầu mình, đôi mắt mở to trông hệt như bị thứ gì đó chiếm lấy. ``Đ-Đau quá\dots\ Nii-chan\dots\ cứu em\dots''

Tôi lập tức lao đến, ôm chặt em vào lòng, che tầm mắt em khỏi bức tranh đó. Nhịp thở của em gấp gáp, run bần bật. Còn tôi, chỉ có thể nhìn sang khung tranh kia với một câu hỏi: Đây là ``kiệt tác'' mà cô gái họa sĩ ấy đã để lại ư? Và tại sao\dots\ lại tạo ra một thứ ám ảnh đến mức này?

*PHỤT!* *ROẠT!*

Chưa kịp tìm câu trả lời, ánh sáng trong phòng vụt tắt. Một âm thanh vang lên từ chính khung tranh. Tôi cảm thấy cơ thể mình cũng nhẹ bẫng đi một chút, như vừa thoát khỏi sức nặng vô hình nào đó.

Saikyou ngẩng đầu, hơi thở đã bình ổn hơn: ``Ể\dots\ Đầu em\dots\ hết đau rồi\dots\ Chuyện gì vừa xảy ra vậy?''

Tôi thở dài, xoa nhẹ vai em: ``Anh\dots\ cũng không biết nữa.''

Ánh đèn trong phòng vụt sáng trở lại, không theo một cách nhẹ nhàng, mà là một cú ``bật'' khô khốc, như thể có ai đó vừa quất công tắc bằng cả bàn tay.

Trước mắt chúng tôi\dots\ bức tranh. Hoặc đúng hơn, đã từng.

Nó đã bị xé đôi. Vết rách ngoằn ngoèo chạy dọc từ trên xuống dưới, xé toạc khuôn mặt thiếu nữ kia thành hai nửa chẳng còn ghép lại được. Những sợi vải toạc ra như vết thương mới, vẫn còn rung nhẹ, như thể vừa bị một bàn tay vô hình bấu mạnh.

Tôi nhìn quanh. Không ai ở đây ngoài chúng tôi. Chắc chắn không ai đụng vào nó.

Có cái gì\dots\ hoặc ai đó\dots\ vừa làm việc này.

Thật quái gở là, thay vì sợ hãi hơn, tôi lại thấy một cơn nhẹ nhõm tràn qua, như thể thứ đó vừa cứu chúng tôi khỏi việc biến thành\dots\ thứ gì đó không còn là con người.

Không biết là tôi có nên biết ơn hay run rẩy vì vừa được ``một bàn tay khác'' cứu mạng.

Cái cảm giác từ bức tranh vẫn chưa hoàn toàn rời khỏi đầu óc tôi, giống như một dư âm, một cái bóng của cảm xúc đau đớn vẫn đè lên tim, nặng trĩu và tanh tưởi.

``Ta phải rời khỏi đây thôi!'' tôi gằn giọng, vừa như ra lệnh, vừa như khẩn cầu.

Saikyou chỉ khẽ ậm ừ, gật đầu, nhưng tôi nhận ra ánh mắt vẫn liếc lại phía bức tranh, như thể vẫn bị hút bởi một sức mạnh nào đó. Tôi nắm tay, kéo em ra khỏi căn phòng Mỹ thuật, cố không để em ấy ngoái đầu lại.

Cánh cửa phòng khép lại sau lưng, nhưng tim tôi vẫn đập nhanh hơn bình thường. Tôi không biết vì sao, nhưng ngay khi cánh cửa đóng, một luồng khí lạnh chạy dọc sống lưng, giống như ai đó đang đứng sát bên, thở khẽ vào tai tôi.

Trước mặt là một cánh cửa trượt khác. Nếu sau cánh cửa này vẫn là một hành lang\dots

*XOẠCH.*

\dots Và đúng là thế thật.

Tôi có thể đã hét lên, nhưng liệu có ích gì? Dù gì, Saikyou làm thay tôi: ``\textbf{À, trời ạ! Lại hành lang nữa sao!?}''

Tiếng vọng của em ấy vang dọc lối đi, rồi tắt ngấm trong lớp bụi mờ.

Cảm giác chán ghét nơi này len vào từng bước chân.

Chúng tôi biết đây là trường trung học Aogami, biết nó nằm giữa thị trấn Aogami, biết đã từng có hàng loạt vụ việc rùng rợn xảy ra quanh nó\dots\

Nhưng câu hỏi lớn nhất vẫn chưa thay đổi từ giây phút đầu tiên tôi tỉnh dậy ở đây: Tại sao chúng tôi lại ở đây?

Đầu tôi rối như tơ vò. Mỗi khúc quanh, mỗi cánh cửa chỉ mở ra thêm một mớ bí ẩn. Tôi thật sự cần một lời giải thích, không phải ngay mai, không phải khi tỉnh dậy như sau một cơn ác mộng, mà là ngay bây giờ.

Nhưng trước cả điều đó, còn một câu hỏi khác bức bách hơn: Đường ra ở đâu?

\dots

\dots

Chúng tôi tiếp tục lê bước qua cái hành lang dài bất tận ấy, tâm trí vẫn nặng trĩu vì những câu hỏi chưa lời đáp.

Bất chợt, giọng Saikyou khẽ vang lên, như sợ phá vỡ sự im lặng đáng sợ bao trùm nơi này: ``\dots Nii-chan.''

Tôi rướn mày, mắt vẫn lia về phía trước. ``Lại gì nữa đây?''

``\dots Cô ta kìa.''

``Hả?''

Cô em gái khẽ nghiêng cằm, chỉ về phía cuối hành lang.

\begin{wrapfigure}[13]{l}{4cm}
  \includegraphics[height=8cm]{../../../../Image/Book?/shino_1.png}
\end{wrapfigure}

Ở đó, trong vùng sáng mờ hắt ra từ bóng đèn huỳnh quang chập chờn, là một cô gái. Mái tóc nâu dài đến lưng, đồng phục đen gọn gàng, đứng thẳng như một pho tượng, thỉnh thoảng lắc lư như một đứa trẻ. Nhưng cái khí chất ấy\dots\ tôi nhận ra ngay: Misora Shino.

Không chớp mắt, cô ta nhìn thẳng về phía chúng tôi. Lạnh lùng. Im lặng.

``\dots Chúng ta có nên tiếp cận ả ta không?'' em gái tôi thì thầm, ánh mắt vừa cảnh giác vừa bồn chồn.

Tôi thực sự không biết cô ta đang toan tính gì. Nhưng có một điều chắc chắn: cô ta coi chúng tôi là ``học sinh chuyển trường'' và tôi không hiểu vì lý do quái gì. Có phải cô ta ghét bọn tôi? Hay chỉ đơn giản là ghét những người giống bọn tôi?

Dù thế nào, Shino vẫn là người ``ám'' cả cái nơi chết tiệt này. Và nếu ai có câu trả lời, thì là cô ta.

``\dots Không còn cách nào nữa. Chẳng còn đường lui nào,'' tôi đáp, giữ mắt nhìn về phía trước. ``Chính em cũng bảo là muốn ba mặt một lời còn gì.''

``\dots Vâng, em nói thế thật. Đã đến nước này, ta phải hỏi cô ta cho ra trò!'' giọng em tôi run run, nhưng trong đó lại có chút quyết tâm.

Quyết định đã được đưa ra. Nếu Shino không chịu nói, chúng tôi sẽ không rời đi.

Chúng tôi tắt đèn pin điện thoại, bước từng bước chậm mà chắc. Hành lang dường như dài thêm ra theo từng nhịp chân, và cái bóng của Shino cứ đứng đó, không hề thay đổi.

\dots

Khi khoảng cách chỉ còn vài mét, Shino khẽ nhướng mày, vẻ mặt thoáng ngạc nhiên, rồi nhanh chóng lấy lại vẻ điềm tĩnh.

``Các cậu đang làm gì ở đây vậy?'' giọng cô ta trầm, đều.

Em gái tôi lập tức phản công: ``Chúng tôi phải hỏi cô câu ấy mới đúng đó! Chuyện quái gì đang xảy ra vậy!?''

Tôi tiếp lời, chất vấn thẳng: ``Với cả, tại sao cô lại gọi bọn tôi là `học sinh chuyển trường'? Tại sao chúng tôi lại ở đây?''

Cô ta im lặng một hồi, nhưng đôi mắt có vẻ như đang không nhìn chúng tôi.

Shino đáp mà như chẳng đáp: ``Nhanh nào, ta cùng về thôi chứ?'' Giọng điệu nhẹ tênh, nhưng tôi nhận ra: hoặc cô ta đang né tránh, hoặc cố tình lờ đi.

``Vấn đề là về thế nào được chứ--'' tôi chưa kịp nói hết câu, thì\dots

``\textbf{\dots Tôi không về!}''

Một giọng con gái vang lên, vọng khắp hành lang. Không lẫn vào đâu được: một giọng đầy cảm xúc, dứt khoát.

Chúng tôi lập tức quay ngoắt lại, giơ ánh đèn điện thoại quét ngang quét dọc.

``Ai đó!?'' tôi gằn giọng. Rồi quay sang Saikyou: ``Em nói đấy à?''

``Đâu có! Giọng này không phải của em!''

\dots Khoan. Sao tôi lại thấy giọng này giống hệt Shino?

Tôi chưa kịp thắc mắc thêm, thì thấy Shino sững lại, ánh mắt thoáng run. Cô ta\dots\ không nói chuyện với chúng tôi nữa, mà đang hướng về phía giọng nói kia.

\begin{wrapfigure}[17]{r}{5cm}
  \centering
  \includegraphics[height=8cm]{../../../../Image/Book?/shino_2.png}
  \caption{Lời từ người biên dịch: Đừng hỏi tôi vì sao mặt của cô ấy trông hề hước vậy. Tôi lấy trực tiếp từ trang \textit{hololive} ERROR đấy.}
\end{wrapfigure}

Và rồi, từ cuối hành lang sau lưng chúng tôi, một bóng người tiến tới. Cũng mái tóc nâu dài, cũng chiếc cài tóc đỏ lệch trên mái, nhưng đôi mắt là màu nâu chứ không phải xanh dương. Trên gương mặt ấy là một nỗi buồn âm ỉ, nhưng xen lẫn một sự kiên định không thể nhầm lẫn.

Bộ đồng phục cô ấy mặc giống hệt Saikyou: áo sơ-mi trắng ngắn tay, nơ xanh dương, váy ca-rô, chỉ thêm chiếc blazer đen khoác ngoài.

Shino mở miệng trước: ``C-Cậu đang nói gì chứ, Sora?''

À\dots\ vậy ra tên cô nàng mắt nâu ấy là Sora à. Hoặc, có lẽ tôi đoán vậy.

``Chuyện dài lắm, tớ không thể nói bây giờ,'' Sora đáp gấp, giọng khàn khàn như phải kìm nén gì đó.

Saikyou chen vào: ``Vậy cô với ả ta là--''

``Bây giờ không tiện giải thích đâu! Hai cậu lùi lại đi! Để tớ lo chuyện này!'' Sora cắt ngang, ánh mắt vẫn không rời Shino.

Chúng tôi dù chẳng hiểu gì, nhưng vẫn nghe theo, lùi lại vài bước. Không khí như đặc quánh, và giữa hành lang, hai bóng hình giống nhau đến kỳ lạ đang đối mặt.

``Tại sao\dots? Cậu vẫn chưa chịu bỏ cuộc sao\dots?'' Shino hỏi, giọng như một vết rạn đang lan.

Sora siết chặt nắm tay: ``Cậu biết tớ vẫn còn mong muốn chưa thể thực hiện được mà! Và một khi tớ đã định, tớ sẽ làm đến cùng!'' nói đến đây, nước mắt đã tràn ra, nhưng giọng cô vẫn rắn rỏi.

``Những gì cậu đang làm\dots\ chỉ vô nghĩa thôi! \textbf{Đừng ảo tưởng nữa! Hãy chấp nhận thực tại và quay về đi, Sora!}'' Shino gần như quát, đôi mắt xanh lóe lên.

Sora lắc đầu, từng chữ bật ra như dao cắt: ``Tớ không muốn cứ ra đi như vậy! Tớ không muốn kết thúc đời học sinh tươi đẹp vẫn còn dang dở ấy! Tớ sẽ không quay về đâu!''

``Cậu\dots\ vẫn không chịu từ bỏ sao, Sora?'' giọng Shino vang lên, nghe vừa run rẩy vừa chất chứa giận dữ bị nén lại. ``Tại sao cậu lại cố chấp ở đây làm gì? Chuyện đó\dots\ chẳng thể thay đổi được gì cả.''

Sora đứng thẳng người, siết chặt tay, ánh mắt như muốn xuyên thấu qua lớp mặt nạ lạnh lùng kia. ``Tớ không thể bỏ mặc mọi thứ như vậy được. Tớ không muốn rời khỏi nơi này\dots\ khi mọi thứ vẫn còn dang dở. Những ngày tháng ở đây\dots\ vẫn còn điều mà tớ trân trọng.''

``Điều trân trọng ư?'' Shino bật cười nhạt, nụ cười đầy mỉa mai. ``Cậu gọi việc tự nhốt mình vào quá khứ là `trân trọng' sao? Cậu đang bám víu vào một cái bóng mà thôi.''

``Ít ra cái bóng ấy vẫn còn ấm áp hơn sự trống rỗng mà cậu muốn tớ chấp nhận!'' Sora đáp trả ngay, giọng cao hơn hẳn, nhưng đôi mắt vẫn hơi rưng rưng. ``Nếu phải rời đi, thì tớ muốn ra đi khi bản thân đã thực sự sẵn sàng. Còn bây giờ\dots\ tớ chưa thể.''

``Nhưng sự thật vẫn là sự thật!'' Shino gằn từng chữ, bước lên một bước, khoảng cách giữa hai người chỉ còn vài mét. ``Cậu càng ở lại, cậu càng tự dày vò mình thôi! Quay về đi, trước khi quá muộn.''

``Không.'' Sora lắc đầu dứt khoát. ``Tớ sẽ ở lại, dù cho cậu có nghĩ thế nào đi nữa.''

Khoảng lặng kéo dài, chỉ còn tiếng quạt thông gió rì rì vang vọng trong hành lang trống trải. Tôi và Saikyou nhìn nhau, chẳng biết nên xen vào hay lùi lại. Cả hai đều hiểu rõ, chúng tôi không phải người trong cuộc, và chuyện này\dots\ không hề đơn giản.

*CẠCH!*

Tiếng gì đó vang lên từ cô gái mặc bộ đồng phục đen, như thể một sợi dây cuối cùng trong cô ấy đã đứt phựt. Dường như cô ta đã vượt quá giới hạn.

Bầu không khí lập tức trở nên đặc quánh, nặng nề đến mức ngay cả hơi thở cũng khó khăn. Tôi thực sự không hiểu hết ý nghĩa trong lời qua tiếng lại kia, nhưng trực giác mách bảo, đây không còn là cuộc tranh cãi thông thường nữa.

``Ờm\dots\ Shino-san\dots\ cô ổn chứ?'' tôi hỏi, giọng nhỏ hơn bình thường.

Saikyou quay phắt sang tôi: ``Anh lo cho cô ta làm gì chứ, Nii-chan!?''

Không một lời đáp lại.

Chỉ có sự im lặng, và một thứ gì đó như đang âm ỉ bên trong Shino.

Bỗng Sora khựng lại, mắt mở to, hơi thở gấp gáp. Tôi liếc theo hướng nhìn của cô ấy, và rồi thấy Shino đang run bần bật, từ từ ngẩng đầu, đôi mắt giờ đã trống rỗng, không còn chút ánh sáng.

``\dots Cậu đã chạy trốn.''

Một bước.

``Chạy trốn khỏi tội lỗi của mình.''

Hai bước.

``Lặp lại\dots\ hết lần này đến lần khác.''

Ba bước.

``Cậu thực sự cứng đầu đấy.''

Thêm một bước.

``Cậu đang cố trở thành cái gì? Và\dots\ điều đó có nghĩa lý gì không?''

Lại một bước nữa.

``\dots Bởi sau cùng, cậu không thể trốn chạy mãi được đâu\dots''

Và rồi, câu cuối cùng, lạnh lẽo đến mức sống lưng tôi tê cứng: ``Tôi sẽ kết thúc cái ảo tưởng hão huyền đó\dots\ ngay tại đây.''

Cùng lúc ấy, cả hành lang bỗng nhuộm đỏ như máu, như thể ánh sáng xung quanh đã bị biến dạng. Shino tiến thẳng về phía chúng tôi, bước đi chậm nhưng nặng nề, giống hệt một kẻ bị nguyền rủa.

\img{e0t}

``\textbf{Không ổn rồi!}'' Sora hét lên. ``\textbf{Cậu ta phát điên rồi! Chạy đi, mau!}''

Không cần hiểu rõ, chúng tôi cũng biết đây là lúc phải rút lui. Cả ba quay người bỏ chạy về phía cánh cửa trượt đang mở sẵn. Tôi liếc lại, thấy Shino vẫn tiến về phía chúng tôi với tốc độ không nhanh\dots\ nhưng cảm giác như khoảng cách chẳng hề xa ra.

``Cậu cứ chạy trốn, chạy trốn mãi\dots'' giọng cô ta vọng theo, khàn đặc và vang vọng như phát ra từ bức tường chính nơi này.

``Đi thôi, Nii-chan!'' Saikyou kéo tay tôi.

Dù chân đang chạy hết sức, tôi vẫn thấy như hơi thở nóng của Shino đang áp sát sau lưng, có lẽ chỉ là ảo giác, hoặc\dots\ cũng có thể không.

Phía trước chúng tôi lại là một hành lang nhuốm một màu đỏ tang tóc nữa.

Nhưng mới chạy được vài bước, chúng tôi đã bị chặn lại bởi một đống bàn ghế chất chồng cao đến mức không thể nào leo qua.

\img{e0u1}

``\textbf{Mẹ kiếp chứ!}'' tôi bật ra một câu chửi. ``Sao tự dưng lại có bàn ghế chân thế này!?''

Em gái tôi nhíu mày: ``Trước đó có thấy nó ở đây đâu nhỉ?''

Chúng tôi thử kéo bàn ghế ra, nhưng từng chiếc như bị đóng đinh xuống sàn, nặng đến mức khó tin. Trong khi đó, tiếng bước chân chậm rãi, đều đặn, nhưng nặng nề của Shino đã vang lên sát ngay sau lưng.

``Mọi người, hướng này!'' Sora đột ngột chỉ vào cánh cửa lớp học bên cạnh và lao tới. Không kịp nghĩ, chúng tôi ùa vào và đóng sầm cửa lại.

Bên ngoài, giọng trầm của cô ta như xuyên qua tấm cửa gỗ: ``Đã kết thúc rồi\dots\ Đừng có chạy trốn nữa\dots''

Từng chữ như đè nặng vào không khí, khiến tôi thấy khó thở.

Bên trong lớp, những chồng bàn ghế khác xếp thành những bức tường méo mó, chỉ để lại một lối đi ngoằn ngoèo như mê cung.

\img{e0u2}

``\textbf{Lại cái quái gì nữa đây!?}'' – Saikyou hét lên.

``Men theo lối này!'' Sora dẫn đầu, luồn lách qua những khe hẹp. Ba chúng tôi đi sát nhau, tiếng thở gấp gáp hòa lẫn tiếng va chạm lạch cạch của bàn ghế khi bị đẩy nhẹ sang hai bên.

Từ phía sau, giọng Shino vẫn vọng tới, đều đều, lạnh lẽo như đang thở mạnh sau gáy: ``\dots Dù có đau khổ thế nào, cậu cũng sẽ phải đối mặt với thực tại thôi\dots''

Ra được cửa bên kia, chúng tôi lao ra một hành lang ngổn ngang như vừa trải qua động đất: sàn nứt lỗ chỗ, mảnh kính vỡ rải khắp nơi, vài thanh xà gỗ trần lơ lửng như chực rơi xuống.

Sora không dừng lại, kéo chúng tôi tới một hành lang khác, nơi hàng loạt ống sắt chắn ngang đường.

\img{e0u3}

``Anh đâu nhớ là có cái này ở đây bao giờ\dots'' tôi lẩm bẩm.

Saikyou chỉ vào mấy lá bùa niêm phong ở các khớp nối: ``Chắc phải tháo mấy cái này thôi.''

``Nhưng vấn đề là''

*XOẸT!* *KENG KENG!*

Chưa kịp nói xong, cô em gái đã giật phăng một lá bùa. Một đoạn ống sắt lập tức rơi xuống, để lại một khoảng trống vừa đủ để đi qua.

``\dots Thật luôn?'' tôi trố mắt.

``Không tin thì anh thử xem,'' em tôi vừa đáp vừa nhặt đoạn ống sắt lên, như thể đó là vũ khí sẵn sàng để dùng.

Tôi chưa kịp phản ứng thì Sora hét lên: ``Cô ấy đuổi kịp rồi!''

Quay lại, tôi thấy Shino đã lật tung đống bàn ghế như chẳng nặng gì. Vừa bước tới, cô vừa nói, giọng vẫn bình thản một cách đáng sợ: ``Cứ trốn mãi như thế\dots\ bộ cậu thấy hạnh phúc lắm sao?''

Chúng tôi chạy tiếp, gặp thêm một dãy ống chắn khác. Lần này, tôi và Saikyou cùng tháo lá bùa, nhưng một đoạn ống rơi thẳng xuống cái hố sâu gần đó, vang lên tiếng *KENG* lạnh gáy.

Hệ thống ống thứ ba hiện ra trước mặt với ba lá bùa. Tôi giật mình khi ngoảnh lại, Shino đã ngay sát, ánh mắt xanh nhạt như soi thấu từng mạch máu của tôi.

Sora và Saikyou gỡ được hai lá bùa, nhưng lá cuối cùng quá chặt. ``Chặt\dots\ quá\dots!'' Saikyou rít lên.

Cùng lúc đó, trần hành lang rung mạnh, những thanh xà gỗ rơi xuống, bụi bay mù mịt. Một thanh gỗ đập mạnh vào vai tôi, khiến tôi ngã sõng soài.

``\textbf{Ặc\dots!}''

``Kyuen-kun!?'' Sora hoảng hốt, vừa lúc bóc xong miếng bùa cuối cùng.

Ngay khi tôi cố chống tay ngồi dậy, một bàn tay lạnh ngắt túm lấy cổ áo tôi và nhấc bổng lên.

``Bắt được rồi\~{}\dots''

Giọng nói ấy khiến tôi sởn gai ốc. Không phải Sora. Không ấm áp, không run rẩy\dots\ mà như một tiếng thì thầm của vực sâu.

Shino.

Ánh mắt cô trống rỗng, miệng vẫn nở một nụ cười ma mị, nhưng lực tay thì không chút do dự, đẩy tôi về phía miệng hố gần đó.

\img{e0v}

``\textbf{Này! Tôi đã làm gì chứ!?}'' tôi gào lên, cố gắng chống lại, nhưng hoàn toàn vô ích.

``Chúng ta\dots\ phải quay về,'' Shino thì thầm.

``\textbf{Không! Chúng ta phải ở lại!}'' Sora hét lên, như thể cố nói cho cả hai cùng nghe.

Nhưng cô ta vẫn tiến từng bước, buộc tôi sát mép hố. Sức tôi đã kiệt, hơi thở nóng rát cổ họng. Tôi biết\dots\ chỉ cần thêm nửa giây nữa\dots

*BỐP!*

Shino khựng lại, lảo đảo về trước. Một bóng người lao tới kéo tôi lùi lại. Là Sora.

``Đừng có hòng động đến Nii-chan của tôi!'' Saikyou gầm lên, tay cầm đoạn ống sắt vừa nãy.

``Cậu ổn chứ!?'' cô ấy nắm chặt lấy tay, dồn hết sức kéo tôi khỏi cái lỗ sâu hoắm ấy, như cách mà tôi làm với cô em gái tôi.

Tôi thở dốc, cố mỉm cười: ``\dots Nghĩ lại thì\dots\ đúng là nên thủ sẵn cái gì đó thật.''

``Đã bảo rồi mà!'' em hất cằm đầy đắc thắng. ``Anh cũng nên rèn luyện sức khỏe đi!''

``\dots Ý hay.''

Sora chỉ đứng im nhìn Shino đang ngã xuống sàn, ánh mắt pha lẫn hối hận và đau lòng: ``Bọn tớ\dots\ thật sự không muốn làm thế này đâu\dots''

``Không cần phải xin lỗi ả ta đâu!'' Saikyou gắt, đẩy hai chúng tôi lên trước, rồi nhặt thanh ống ban nãy lên. ``Đi trước khi cô ta tỉnh lại!''

Chúng tôi lao tới cánh cửa trượt phía trước, bỏ lại sau lưng màu đỏ u tối cùng tiếng bước chân có thể sẽ lại vang lên bất cứ lúc nào.

*XOẠCH*

Bên kia là một hành lang khác, sàn gỗ thủng lỗ chỗ như bị mối gặm, từng mảng gỗ bị bong lên, để lộ ra khoảng không đen ngòm bên dưới. Không khí vẫn đỏ quạch như máu đông, trùm lên mọi thứ một màu đáng sợ. Nhưng\dots\ có gì đó khác. Ngoài ra, trông nó còn sạch sẽ và tươm tất hơn - nêu tôi muốn dùng từ như vậy với cái hành lang kiểu này - so với hành lang trước đó.

\img{e0u4}

Tôi khẽ cau mày. ``Có vẻ\dots\ đây là hành lang của thời hiện tại.''

Hoặc ít nhất, tôi muốn tin là vậy.

Nhưng ngoại trừ một chi tiết khiến toàn bộ suy đoán của tôi lung lay.

Ngay giữa những lỗ lớn trên sàn, từng cánh cổng torii đỏ tươi nhấp nhô như xương sống của một con quái vật khổng lồ. Kế đó, chen giữa các khe hở là những miếu đá nhỏ, giống hệt loại hay thấy ở đền thờ, nhưng ở đây\dots\ chúng mọc lên như thể ai đó cố tình đặt để chặn đường.

``\textbf{Chuyện quái gì đang xảy ra thế này!?}'' tôi và Saikyou đồng thanh kêu lên.

Nhảy qua ư? Cái lỗ rộng đến mức tôi không dám chắc chân mình sẽ đáp xuống bên kia. Đẩy miếu? Chúng làm bằng đá đặc, nặng gấp mấy lần cơ thể tôi, đẩy thế quái nào được?

``Để tớ lo!'' Sora nói nhanh, giọng không chút do dự.
Cô nhắm mắt, chắp tay, môi mấp máy niệm những âm tiết lạ lẫm như gió cuốn.

Trong nháy mắt, chiếc miếu đá chắn ngay trước mặt chúng tôi rung lên rồi trượt ngang, rơi *RẦM* xuống vực tối bên dưới, như bị một bàn tay vô hình xô ngã.

``Đường thông rồi đó!'' Sora mở mắt, ánh nhìn như thúc giục.

Chúng tôi men theo con đường hẹp, len qua những cánh cổng torii vươn ra sát mặt, mỗi lần lách qua là một lần nghe tiếng gỗ khẽ kêu cót két như có ai chạm vào.

Phía trước lại là một miếu đá chắn ngang. Từ đâu đó, một giọng thì thầm khẽ vọng tới: ``Mãi mãi\dots\ mãi mãi\dots\ mãi mãi\dots''

Tôi cảm giác máu mình như đông lại.

Khốn thật. Shino. Cô ấy đã đuổi tới đây sao?

Tôi ngoảnh lại. Không có ai.

Chỉ là hành lang đỏ đặc quánh và sự im lặng khiến tai tôi ù đi. Nhưng sống lưng tôi lạnh ngắt, như có một ánh mắt đang dán chặt vào gáy.

Sora lại chắp tay niệm. Miếu đá trước mặt khẽ nghiêng, rồi rơi khỏi mép sàn. Chúng tôi lập tức vượt qua.

Có lẽ vì quá sợ, tôi bất giác vượt lên trước hai người con gái, lao một mạch về phía cánh cửa trượt tiếp theo.

``Này, đợi em với, Nii-chan!'' Saikyou gọi giật lại.

``Chạy nhanh lên nào, Saikyou-chan!'' Sora tiếp lời.

Và đó là lúc tôi chợt nhận ra\dots

Khoan đã. Cô ấy biết tên chúng tôi từ khi nào?

Nhưng đây không phải lúc để hỏi.

*XOẠCH!*

Cánh cửa trượt mở ra, để lộ một hành lang tối hơn hẳn. Không có dấu hiệu sập đổ, nhưng cũng không hề dễ chịu hơn. Ánh sáng từ những khung cửa sổ rọi vào chẳng đủ xua bớt bóng tối đang tràn ngập.

*HỘC HỘC\dots*

*HỘC HỘC\dots*

``Đằng trước\dots\ có một cánh cửa trượt kìa!'' Sora thở gấp, chỉ về phía ánh sáng trắng rực hắt ra từ khe cửa.

\img{e0w}

``Lối ra\dots?'' cả hai chúng tôi vừa chạy vừa thở, cảm tưởng như sắp hết hơi tới nơi.

``Tớ nghĩ là vậy!'' cô đáp nhanh, đôi mắt ánh lên hy vọng.

``Cuối cùng\dots\ thì\dots'' Saikyou thở hắt ra, nhưng vẫn không giảm tốc độ.

Tôi quay lại nhìn một lần nữa.

Cánh cửa mà tôi mở ban nãy\dots\ đã đóng lại từ lúc nào.

Có lẽ Shino sẽ không đuổi theo nữa\dots\ đúng chứ?

\dots

\dots

Tiếng thở ngắt quãng của ba chúng tôi vẫn tiếp tục.

Cánh cửa sáng trắng đã ngay trước mắt.

Chúng tôi cứ tiếp tục chạy, hơi thở gấp gáp hòa cùng tiếng bước chân vang vọng trong hành lang tối. Khoảng cách giữa chúng tôi và cánh cửa trượt sáng rực kia càng lúc càng ngắn lại, như thể từng bước đi đang kéo gần lại ước mơ rời khỏi cái nơi quái quỷ này.

``Chỉ cần chúng ta mở cánh cửa này\dots'' tôi thở hổn hển, vẫn chạy.

``\dots là chúng ta sẽ thoát khỏi đây, đúng không?'' em gái tôi tiếp lời, mắt vẫn dán vào ánh sáng trắng ở cuối đường, rồi nhìn sang Sora như chờ một sự khẳng định.

Nhưng\dots\ đáp lại chỉ là một khoảng im lặng kéo dài đến khó chịu.

Tôi tưởng cô ấy không nghe thấy, cho đến khi\dots

``\dots Hả? Cái gì cơ?'' cô ấy bất chợt lên tiếng, khiến cả hai anh em giật bắn mình. Tim tôi suýt nữa thì rụng ra ngoài.

``Tôi hỏi, đây có phải là lối thoát hay không?'' Saikyou nhắc lại, giọng bắt đầu lộ vẻ sốt ruột.

``À\dots\ ừ\dots\ Tớ nghĩ\dots\ chắc vậy\dots'' Sora trả lời, nhưng giọng cô ấy không hề chắc chắn. Không khó để nhận ra sự lưỡng lự ấy, và tôi không thể trách cô ấy. Chúng tôi đâu biết Sora là ai, chỉ biết cô ấy giống Shino một cách kỳ lạ, như hai bản sao hoàn hảo.

Khoảng lặng lại phủ xuống. Chỉ còn tiếng thở nặng nề của chúng tôi vang vọng giữa hành lang.

Rồi, Sora khẽ cất tiếng, chậm rãi, như thể phải lấy hết can đảm mới nói ra: ``\dots Tớ\dots\ thực sự muốn được tận hưởng tuổi học trò.''

Cả tôi và Saikyou khựng lại, không phải vì đôi chân, mà vì câu nói ấy. Một câu nói được bộc phát ra một cách bất ngờ, ở một thời điểm mà chúng tôi không ngờ tới.

``Tớ muốn được kết bạn\dots\ được chơi cùng mọi người\dots\ được cười nói sau giờ học\dots\ Tớ muốn biết cảm giác đó là thế nào\dots\ khi mình thực sự có bạn\dots'' giọng cô ấy run run, đôi bàn tay nắm chặt tà váy như đang cố giữ lại chút can đảm mong manh.

Tôi thoáng hình dung ra một quãng đời trống rỗng của cô ấy: không sân trường, không tiếng cười, không có cả những trò nghịch dại. Có lẽ đã có một biến cố nào đó… thứ đã lấy đi của Sora tất cả những điều ấy.

Trong khoảnh khắc ấy, tôi muốn tin rằng đây là ước nguyện của Shino. Nhưng\dots\ không, cảm giác mách bảo tôi rằng Sora là một người khác, với nỗi khao khát rất riêng của mình.

``Thế nên\dots'' cô ấy ngẩng đầu lên, đôi mắt nâu ấy hướng thẳng vào chúng tôi, ánh nhìn vừa khẩn thiết vừa sợ hãi bị từ chối.

``\dots Các cậu\dots\ \emph{có thể làm bạn với tớ} không?''

Lời cầu xin. Không màu mè, không giả tạo. Chỉ là một trái tim trần trụi đưa ra trước mặt chúng tôi. Tôi bỗng thấy cổ họng mình nghẹn lại.

``Anh nghĩ sao, Nii-chan?'' Saikyou quay sang tôi hỏi.

``\dots'' tôi im lặng. Thật sự, để nói lời đồng ý ngay thì rất khó. Chúng tôi chỉ mới gặp cô ấy chưa đầy mười phút. Thậm chí ở một nghịch cảnh đầy oái oăm và rối rắm này nữa.

Nhưng\dots\ ước vọng ấy, tôi có thể cảm nhận được. Và nó quá chân thành để mà làm ngơ.

\dots Dù thế, tôi không thể không thấy cảm giác quen thuộc. Như thể trước đây tôi đã từng đứng trước một lựa chọn tương tự\dots\ nhưng ở đâu, và khi nào, thì tôi không tài nào nhớ nổi.

Sora vẫn cúi đầu, chờ đợi câu trả lời từ hai chúng tôi. Cái dáng vẻ ấy khiến tôi cảm thấy mình mà từ chối thì sẽ đạp nát một điều gì đó mong manh lắm.

Tôi thở dài một hơi. Có vẻ như không còn lựa chọn nào khác nhỉ?

``\dots Được chứ,'' tôi nói.

Tôi nhìn sang Saikyou: ``Còn em?''

``Anh--'' Em ấy mở miệng định phản đối, nhưng tôi đã chặn lại: ``Dù gì cô ấy cũng có thành ý. Và\dots\ trông tội lắm. Anh không nghĩ cô ấy có ý định làm hại chúng ta đâu.''

Saikyou chỉ biết thở dài: ``Anh đúng là đồ ngốc tốt bụng\dots\ Thôi được rồi.''

Em gái tôi bước lên, chìa tay ra: ``Miễn là cậu đừng có giành Nii-chan của tôi, thì tôi sẽ coi cậu là bạn.'' Má em ấy phồng lên rõ rệt, chẳng khác nào một lời tuyên bố chủ quyền, như thể sợ ai đó cướp mất vật quý của mình.

Sora thoáng sững lại, rồi bật khóc. Không phải khóc vì buồn, mà là vì nhẹ nhõm và vui mừng. Cô ấy nắm lấy tay Saikyou, siết chặt như sợ nó biến mất.

``Được chứ! Tớ\dots\ tớ mong chúng ta sẽ mãi là bạn của nhau!''

``Tất nhiên rồi,'' Saikyou mỉm cười, ``bạn bè thì phải giúp đỡ nhau, cả với anh trai tớ nữa.''

``Ừ! Tớ hứa.''

Và như vậy, giữa cái nơi quỷ dị này, chúng tôi đã tìm được một người bạn mới,  người đầu tiên kể từ khi lạc vào đây.

Chúng tôi cùng nhau quay lại nhìn cánh cửa trượt rực sáng, nở nụ cười mà đã lâu rồi tôi mới cảm thấy thật sự.

``Vậy\dots\ chúng ta đi thôi,'' tôi nói.

Cả hai cô gái gật đầu. Tôi tiến lên, bàn tay chạm vào cánh cửa, đẩy nó ra, để ánh sáng ùa vào.

Một khởi đầu mới. Và có lẽ, một tình bạn vừa chớm nở\dots

\pov{-}

Shino đứng trong bóng tối nơi hành lang vắng, lặng lẽ dõi theo ba người họ bước qua cánh cửa trượt ngập ánh sáng.

Không tiếng gọi với theo, không bước chân đuổi theo.

Chỉ là đôi mắt xanh ấy khẽ khép lại, và trên môi nở một nụ cười.

Không phải là nụ cười chiến thắng.

Cũng chẳng phải sự cay cú vì cú đấm của cô em gái ban nãy.

Mà là một nụ cười rất nhỏ\dots\ chất chứa nỗi buồn sâu hun hút, hòa lẫn với thứ gì đó giống như sự tiếc nuối, tiếc cho số phận nghiệt ngã của mình, cho những lỗi lầm mà cô tin rằng tất cả đều bắt nguồn từ chính bản thân.

Ánh sáng từ khe cửa hắt lên khuôn mặt Shino, để lộ làn da tái nhợt và khóe mắt ươn ướt.

\img{e0x}

Cô nhìn khoảng trống mà ba người họ vừa đi qua, khẽ thì thầm với giọng trầm lẫn run:

\begin{center}
  ``Những ảo tưởng đó\dots\ ai sẽ cứu được cậu?\\
  Chính cậu\dots\ chính cậu đã khởi đầu đại họa này\dots''
\end{center}

Cô ngẩng đầu lên, như đang nói với một người vô hình đứng ở đâu đó trong khoảng không:

\begin{center}
  ``Cậu có thể chạy, chạy mãi\dots\\
  Nhưng rồi cậu cũng không thể nào trốn thoát khỏi sự thật.''
\end{center}

Tiếng gió lạnh rít qua những ô cửa vỡ, kéo dài khoảng lặng như muốn nuốt chửng lấy mọi lời nói.

Shino khẽ cười một lần nữa, nhưng lần này yếu ớt hơn, như đang tự trách mắng một phiên bản khác của mình, cái ``bản sao'' mang đôi mắt nâu kia.

\begin{center}
  ``Đến bao giờ\dots\ cậu mới chịu thừa nhận tôi đây?\\
  Đến bao giờ\dots\ cậu mới chịu chấp nhận cái số phận đau khổ này?\\
  Và\dots\ đến bao giờ\dots\ chúng ta mới thoát được khỏi cái vòng luân hồi chết tiệt này?''
\end{center}

Giọng cô vỡ ra ở những từ cuối, hòa vào một tiếng thở dài buốt lạnh.

Trong giây lát, ánh sáng từ cánh cửa phản chiếu lên đôi mắt xanh, biến chúng thành thứ ánh sáng u uất, như chứa đựng cả một đại dương đầy khao khát và tuyệt vọng.

Rồi, như một bóng hình chưa từng tồn tại, cơ thể Shino tan vào khoảng không, để lại hành lang đỏ đặc quánh sự im lặng.

Chỉ còn lại cánh cửa kia, vẫn mở rộng, nhưng ánh sáng bên trong đã không còn dành cho cô nữa.

\end{document}